% Options for packages loaded elsewhere
% Options for packages loaded elsewhere
\PassOptionsToPackage{unicode}{hyperref}
\PassOptionsToPackage{hyphens}{url}
\PassOptionsToPackage{dvipsnames,svgnames,x11names}{xcolor}
%
\documentclass[
  spanish,
  letterpaper,
  DIV=11,
  numbers=noendperiod]{scrreprt}
\usepackage{xcolor}
\usepackage{amsmath,amssymb}
\setcounter{secnumdepth}{5}
\usepackage{iftex}
\ifPDFTeX
  \usepackage[T1]{fontenc}
  \usepackage[utf8]{inputenc}
  \usepackage{textcomp} % provide euro and other symbols
\else % if luatex or xetex
  \usepackage{unicode-math} % this also loads fontspec
  \defaultfontfeatures{Scale=MatchLowercase}
  \defaultfontfeatures[\rmfamily]{Ligatures=TeX,Scale=1}
\fi
\usepackage{lmodern}
\ifPDFTeX\else
  % xetex/luatex font selection
\fi
% Use upquote if available, for straight quotes in verbatim environments
\IfFileExists{upquote.sty}{\usepackage{upquote}}{}
\IfFileExists{microtype.sty}{% use microtype if available
  \usepackage[]{microtype}
  \UseMicrotypeSet[protrusion]{basicmath} % disable protrusion for tt fonts
}{}
\makeatletter
\@ifundefined{KOMAClassName}{% if non-KOMA class
  \IfFileExists{parskip.sty}{%
    \usepackage{parskip}
  }{% else
    \setlength{\parindent}{0pt}
    \setlength{\parskip}{6pt plus 2pt minus 1pt}}
}{% if KOMA class
  \KOMAoptions{parskip=half}}
\makeatother
% Make \paragraph and \subparagraph free-standing
\makeatletter
\ifx\paragraph\undefined\else
  \let\oldparagraph\paragraph
  \renewcommand{\paragraph}{
    \@ifstar
      \xxxParagraphStar
      \xxxParagraphNoStar
  }
  \newcommand{\xxxParagraphStar}[1]{\oldparagraph*{#1}\mbox{}}
  \newcommand{\xxxParagraphNoStar}[1]{\oldparagraph{#1}\mbox{}}
\fi
\ifx\subparagraph\undefined\else
  \let\oldsubparagraph\subparagraph
  \renewcommand{\subparagraph}{
    \@ifstar
      \xxxSubParagraphStar
      \xxxSubParagraphNoStar
  }
  \newcommand{\xxxSubParagraphStar}[1]{\oldsubparagraph*{#1}\mbox{}}
  \newcommand{\xxxSubParagraphNoStar}[1]{\oldsubparagraph{#1}\mbox{}}
\fi
\makeatother


\usepackage{longtable,booktabs,array}
\usepackage{calc} % for calculating minipage widths
% Correct order of tables after \paragraph or \subparagraph
\usepackage{etoolbox}
\makeatletter
\patchcmd\longtable{\par}{\if@noskipsec\mbox{}\fi\par}{}{}
\makeatother
% Allow footnotes in longtable head/foot
\IfFileExists{footnotehyper.sty}{\usepackage{footnotehyper}}{\usepackage{footnote}}
\makesavenoteenv{longtable}
\usepackage{graphicx}
\makeatletter
\newsavebox\pandoc@box
\newcommand*\pandocbounded[1]{% scales image to fit in text height/width
  \sbox\pandoc@box{#1}%
  \Gscale@div\@tempa{\textheight}{\dimexpr\ht\pandoc@box+\dp\pandoc@box\relax}%
  \Gscale@div\@tempb{\linewidth}{\wd\pandoc@box}%
  \ifdim\@tempb\p@<\@tempa\p@\let\@tempa\@tempb\fi% select the smaller of both
  \ifdim\@tempa\p@<\p@\scalebox{\@tempa}{\usebox\pandoc@box}%
  \else\usebox{\pandoc@box}%
  \fi%
}
% Set default figure placement to htbp
\def\fps@figure{htbp}
\makeatother


% definitions for citeproc citations
\NewDocumentCommand\citeproctext{}{}
\NewDocumentCommand\citeproc{mm}{%
  \begingroup\def\citeproctext{#2}\cite{#1}\endgroup}
\makeatletter
 % allow citations to break across lines
 \let\@cite@ofmt\@firstofone
 % avoid brackets around text for \cite:
 \def\@biblabel#1{}
 \def\@cite#1#2{{#1\if@tempswa , #2\fi}}
\makeatother
\newlength{\cslhangindent}
\setlength{\cslhangindent}{1.5em}
\newlength{\csllabelwidth}
\setlength{\csllabelwidth}{3em}
\newenvironment{CSLReferences}[2] % #1 hanging-indent, #2 entry-spacing
 {\begin{list}{}{%
  \setlength{\itemindent}{0pt}
  \setlength{\leftmargin}{0pt}
  \setlength{\parsep}{0pt}
  % turn on hanging indent if param 1 is 1
  \ifodd #1
   \setlength{\leftmargin}{\cslhangindent}
   \setlength{\itemindent}{-1\cslhangindent}
  \fi
  % set entry spacing
  \setlength{\itemsep}{#2\baselineskip}}}
 {\end{list}}
\usepackage{calc}
\newcommand{\CSLBlock}[1]{\hfill\break\parbox[t]{\linewidth}{\strut\ignorespaces#1\strut}}
\newcommand{\CSLLeftMargin}[1]{\parbox[t]{\csllabelwidth}{\strut#1\strut}}
\newcommand{\CSLRightInline}[1]{\parbox[t]{\linewidth - \csllabelwidth}{\strut#1\strut}}
\newcommand{\CSLIndent}[1]{\hspace{\cslhangindent}#1}

\ifLuaTeX
\usepackage[bidi=basic]{babel}
\else
\usepackage[bidi=default]{babel}
\fi
% get rid of language-specific shorthands (see #6817):
\let\LanguageShortHands\languageshorthands
\def\languageshorthands#1{}


\setlength{\emergencystretch}{3em} % prevent overfull lines

\providecommand{\tightlist}{%
  \setlength{\itemsep}{0pt}\setlength{\parskip}{0pt}}



 


\usepackage{booktabs}
\usepackage{longtable}
\usepackage{array}
\usepackage{multirow}
\usepackage{wrapfig}
\usepackage{float}
\usepackage{colortbl}
\usepackage{pdflscape}
\usepackage{tabu}
\usepackage{threeparttable}
\usepackage{threeparttablex}
\usepackage[normalem]{ulem}
\usepackage{makecell}
\usepackage{xcolor}
\KOMAoption{captions}{tableheading}
\makeatletter
\@ifpackageloaded{bookmark}{}{\usepackage{bookmark}}
\makeatother
\makeatletter
\@ifpackageloaded{caption}{}{\usepackage{caption}}
\AtBeginDocument{%
\ifdefined\contentsname
  \renewcommand*\contentsname{Tabla de contenidos}
\else
  \newcommand\contentsname{Tabla de contenidos}
\fi
\ifdefined\listfigurename
  \renewcommand*\listfigurename{Listado de Figuras}
\else
  \newcommand\listfigurename{Listado de Figuras}
\fi
\ifdefined\listtablename
  \renewcommand*\listtablename{Listado de Tablas}
\else
  \newcommand\listtablename{Listado de Tablas}
\fi
\ifdefined\figurename
  \renewcommand*\figurename{Figura}
\else
  \newcommand\figurename{Figura}
\fi
\ifdefined\tablename
  \renewcommand*\tablename{Tabla}
\else
  \newcommand\tablename{Tabla}
\fi
}
\@ifpackageloaded{float}{}{\usepackage{float}}
\floatstyle{ruled}
\@ifundefined{c@chapter}{\newfloat{codelisting}{h}{lop}}{\newfloat{codelisting}{h}{lop}[chapter]}
\floatname{codelisting}{Listado}
\newcommand*\listoflistings{\listof{codelisting}{Listado de Listados}}
\makeatother
\makeatletter
\makeatother
\makeatletter
\@ifpackageloaded{caption}{}{\usepackage{caption}}
\@ifpackageloaded{subcaption}{}{\usepackage{subcaption}}
\makeatother
\usepackage{bookmark}
\IfFileExists{xurl.sty}{\usepackage{xurl}}{} % add URL line breaks if available
\urlstyle{same}
\hypersetup{
  pdftitle={lisa-book},
  pdfauthor={Laboratorio de Investigación Social Abierta},
  pdflang={es},
  colorlinks=true,
  linkcolor={blue},
  filecolor={Maroon},
  citecolor={Blue},
  urlcolor={Blue},
  pdfcreator={LaTeX via pandoc}}


\title{lisa-book}
\author{Laboratorio de Investigación Social Abierta}
\date{2025-11-01}
\begin{document}
\maketitle

\renewcommand*\contentsname{Tabla de contenidos}
{
\hypersetup{linkcolor=}
\setcounter{tocdepth}{2}
\tableofcontents
}

\bookmarksetup{startatroot}

\chapter*{Presentación}\label{presentaciuxf3n}
\addcontentsline{toc}{chapter}{Presentación}

\markboth{Presentación}{Presentación}

Libro colaborativo del Laboratorio de Investigación Social Abierta

\bookmarksetup{startatroot}

\chapter{Introducción}\label{introducciuxf3n}

El interés de realizar investigaciones científicas con base en la
Ciencia Social Abierta ha ido aumentando en la última década. Distintos
centros de investigaciones y organismos nacionales e internacionales han
posicionado como un estándar mínimo de calidad, confianza y
transparencia que los(as) científicos sociales suscriban sus
investigaciones bajo los principios de la Ciencia Abierta. En este libro
se abordan los conceptos principales que explican de qué se trata la
Ciencia Abierta, sin embargo, el objetivo final es convertirse en una
guía para los(as) investigadores de las Ciencias Sociales. A partir de
la revisión de los principios a los que suscriben los principales
organismos internacionales y de las principales herramientas utilizadas
para cumplir con este propósito es que este libro permite consolidar las
experiencias actuales para transformarlas en una serie de
recomendaciones que les permitan a los(as) investigadores de las
Ciencias Sociales empíricas suscribirse bajo los estándares de la
Ciencia Social Abierta. Este libro se guía bajo cuatro principios clave
de la Ciencia Abierta y que constituyen cada uno de los capítulos
tratados a continuación: 1) Diseño transparente; 2) Apertura de datos;
3) Análisis reproducibles y 4) Publicaciones libres.

En el primer capítulo se aborda uno de los principales estándares que
guían este tipo de investigaciones y que tiene que ver con la
transparencia. En este libro, transparencia se asocia con la apertura
pública del diseño de investigación, lo que refiere a la práctica
científica de abrir al público todo tipo de información que sea
relevante en la ejecución del estudio, desde las hipótesis hasta el plan
de análisis. Así, quienes lean los hallazgos de la investigación (tanto
científicos como el público en general) podrán evaluar la credibilidad
de estos y descartar la influencia de prácticas cuestionables de
investigación.

En el segundo capítulo se aborda un elemento clave para fomentar la
colaboración en las investigaciones de las Ciencias Sociales. La
publicación de los datos de investigación (open data) se suscribe como
un estándar de calidad internacional que permite contribuir a que
los(as) investigadores se centren en la investigación y creación de
hallazgos relevantes, antes que en la producción de datos. Más aún, se
señala la importancia de publicar los datos de investigación cumpliendo
con estándares internacionales sobre los datos y la documentación, de
modo tal que cualquier investigador sea capaz de reutilizar los datos
sin la necesidad de contactarse con el equipo que los produjo y con la
información suficiente para reconocer o citar el aporte del equipo
productor de los datos.

En el tercer capítulo se aborda la discusión existente sobre la
reproducibilidad y replicabilidad de una investigación, poniendo el
principal foco en la reproducibilidad de los resultados de
investigación. Reproducibilidad refiere a una publicación detallada de
la documentación, código y datos que permitan obtener los mismos
resultados publicados por los autores, enfatizando que los análisis sean
transparentes y claros con el objetivo de ser verificados por sus pares.
En el escenario actual, la falta de reproducibilidad de los resultados
afecta la robustez y credibilidad de la evidencia reportada por las
investigaciones, principalmente por la falta de precisión en los
procedimientos y las barreras de acceso a materiales clave del proceso
de análisis (Freese y Peterson 2017).

Finalmente, el cuarto capítulo está enfocado en la parte final de una
investigación y se centra en el proceso de divulgación científica. Para
este propósito se abordan dos puntos clave. Por un lado, se aborda la
idea de propiedad intelectual junto a los elementos centrales de la
producción científica y el resguardo legal que poseen los(as) autores
sobre sus investigaciones. Por otro lado, se aborda el movimiento del
Open Access (OA), que tiene por objetivo promover el acceso libre a la
producción de conocimiento científico digital y hacer uso de él sin
restricciones de \emph{copyright}, instalando la posibilidad de que
cualquier persona pueda reutilizar la información de la investigación en
cuanto a datos, procesos y resultados.

\bookmarksetup{startatroot}

\chapter{Diseño transparente}\label{diseuxf1o-transparente}

Uno de los principios que guía el quehacer científico es el escepticismo
organizado (Merton 1973). Ningún hallazgo se acepta porque sí, sino que
los científicos deben demostrarle a otros miembros de la comunidad
científica que \textbf{sus hallazgos son válidos}. En orden de lograr
esto es que se han desarrollado varios tipos de herramientas como los
indicadores de impacto, la revisión por pares o las convenciones para el
reporte de análisis estadístico. Sin embargo, casos como los de Diderik
Stapel han dejado al descubierto que este tipo de herramientas pueden no
ser suficientes para asegurar la credibilidad de los hallazgos
científicos.

Diderik Stapel fue un investigador en el campo de la psicología social.
Su carrera se caracterizó por una trayectoria ejemplar: entre 1995 y
2015 publicó aproximadamente 150 artículos en revistas científicas,
algunas de las más prestigiosas (e.g.~\emph{Science}). También, fue
galardonado con el premio a la trayectoria académica por la Sociedad de
Psicología Social Experimental (SSEP, por sus siglas en inglés). Stapel
llegó a ser una figura reconocida en su área, sin embargo, todo cambió
el año 2011 cuando se confirmó que gran parte de sus artículos contenían
\textbf{datos inventados}. La investigación final de su caso determinó
que más de 50 de sus artículos eran fraudulentos (Carey 2011). De hecho,
Stapel figura dentro de los diez investigadores con más artículos
retractados (58) en \href{https://retractionwatch.com/}{\emph{Retraction
Watch}}, una plataforma que se dedica a sistematizar y mantener una
lista actualizada sobre retracciones de artículos.

A efectos de este escrito, lo más destacable es que Stapel cometió mala
conducta académica durante más de 15 años. \emph{¿Cómo fue esto posible?
¿cómo ninguno de sus colegas, alumnos o coautores se dio cuenta de sus
malas prácticas?} La respuesta tiene que ver con el concepto central de
este capítulo: por la falta de \textbf{transparencia} durante el proceso
de investigación. Stapel se caracterizaba por hacer el trabajo solo y
``a puertas cerradas'', es decir, nadie más que él tenía acceso a los
datos brutos, ni tampoco a la ejecución de las pruebas estadísticas.
Generalmente, sólo compartía con sus colegas los datos procesados y con
los análisis ya efectuados. A raíz de esta forma de trabajo es que se
abría la oportunidad para que Stapel inventara datos a su conveniencia.
Como nadie más colaboraba con el procesamiento de datos, ni tampoco
parecía extraño que así fuera, Stapel pudo sostener una trayectoria
académica destacable en base a la falta de transparencia.

El ejemplo de Stapel es solo una forma de poner en discusión la
importancia de la transparencia en el proceso de investigación.
Siguiendo esta línea, este capítulo busca presentar una primera
aproximación, tanto teórica como práctica, sobre la transparencia en la
investigación social. El objetivo es que investigadoras e investigadores
de las ciencias sociales empíricas se introduzcan en la temática de la
transparencia, así como también que conozcan las principales
herramientas que se están utilizando para adoptarla.

\section{Sobre la transparencia}\label{sobre-la-transparencia}

La transparencia es un concepto multidimensional. (Elliott 2020) propone
entender la transparencia en torno a cuatro preguntas: \emph{¿por qué?,
¿quién? ¿qué? y ¿cómo?} (ver Figura~\ref{fig-taxonomy}). Cada una de
estas preguntas se relaciona a una dimensión. La primera pregunta
(\emph{¿por qué?}) se refiere a las razones y propósitos por los cuales
es necesario adoptar la transparencia; la segunda pregunta
(\emph{¿quién?}) apunta a la audiencia que está recibiendo la
información; la tercera pregunta (\emph{¿qué?}) hace alusión al
contenido que es transparentado y la cuarta pregunta (\emph{¿cómo?})
consiste en cuatro dimensiones distintas sobre cómo adoptar la
transparencia: actores, mecanismos, tiempo y espacios. También, esta
taxonomía estipula una dimensión sobre las amenazas que podrían afectar
a las iniciativas que busquen promover la transparencia. A raíz de estas
dimensiones podemos comprender qué tipo de prácticas y situaciones caben
dentro del concepto de transparencia.

\begin{figure}

\centering{

\includegraphics[width=0.75\linewidth,height=\textheight,keepaspectratio]{images/taxonomy.png}

}

\caption{\label{fig-taxonomy}Taxonomía de Transparencia}

\end{figure}%

Basándonos en esta taxonomía, presentaremos una propuesta sobre qué se
entiende por transparencia en las ciencias sociales y cómo podemos
adquirir prácticas orientadas a este principio. \textbf{Nuestra
propuesta es entender la transparencia como la apertura pública del
diseño de investigación}, lo que incluye todo tipo de información
importante para la ejecución del estudio, desde las hipótesis hasta el
plan de análisis. Esto permitirá que las personas que lean los hallazgos
de investigación (ya sean científicos o la ciudadanía) puedan evaluar la
credibilidad de estos y descartar la influencia de prácticas
cuestionables de investigación. Para llevar a la práctica esta
propuesta, presentaremos la elaboración de preregistros como una de las
principales herramientas que contribuyen a hacer público el diseño,
además de otras que complementan su uso.

\section{¿Por qué?}\label{por-quuxe9}

La pregunta del \textbf{por qué} es necesario avanzar hacia la
transparencia encuentra su respuesta en la existencia de prácticas de
investigación que merman la credibilidad de los hallazgos científicos. A
modo de entender qué es lo problemático de ciertas prácticas, es que
(Steneck 2006) propone el concepto de \emph{Conducta Responsable de
Investigación (RCR, por sus siglas en inglés)} como un ideal que engloba
prácticas éticas e íntegras dentro de la investigación científica (ver
Figura~\ref{fig-rcr}). Según este autor, la distinción entre la ética y
la integridad recae en que la primera tiene que ver con seguir
principios morales dentro de la investigación (e.g.~usar consentimientos
informados), en cambio, la segunda está más relacionada con el
seguimiento de códigos de conductas y estándares profesionales (Abril
Ruiz 2019).

\begin{figure}

\centering{

\includegraphics[width=1\linewidth,height=\textheight,keepaspectratio]{images/rcr.png}

}

\caption{\label{fig-rcr}Conducta Responsable de Investigación. Imagen de
Abril Ruiz (2019) basada en Steneck (2006)}

\end{figure}%

Las prácticas de investigación se pueden evaluar en un continuo que
representa cuánto adhieren los investigadores a los principios de
integridad científica (Steneck 2006). La Figura~\ref{fig-grad}
esquematiza esta idea mostrando dos extremos, donde a la izquierda está
el mejor comportamiento (RCR) y a la derecha el peor comportamiento
(FFP). Las FFP son una abreviación en lengua inglesa para referirse a
\emph{Fabrication, Falsification, Plagiarism} (Invención, Falsificación
y Plagio), también conocidas como \emph{mala conducta académica}. En el
medio del continuo están las \emph{prácticas cuestionables de
investigación} (QRP, por sus siglas en inglés) las cuales refieren a las
``acciones que violan los valores tradicionales de la empresa de
investigación y que pueden ser perjudiciales para el proceso de
investigación'' (\emph{National Academies of Science} 1992 en Steneck
2006, 58). Las QRP se caracterizan porque tienen el potencial de dañar
la ciencia, en tanto las FFP la dañan directamente.

\begin{figure}

\centering{

\includegraphics[width=1\linewidth,height=\textheight,keepaspectratio]{images/grad.png}

}

\caption{\label{fig-grad}Gradación del comportamiento integro en
investigación. Imagen de Abril Ruiz (2019) basada en Steneck (2006)}

\end{figure}%

La mala conducta académica implica la violación de los principios de
integridad científica. Esto se relaciona con el falseamiento de datos de
cualquier forma: invención de datos, alteración de gráficos o tablas,
entre otros. El caso de Diderik Stapel presentado en la introducción de
este capítulo es un ejemplo que cabe dentro del concepto de mala
conducta académica. Según la literatura, la prevalencia de este tipo de
prácticas no es alta (e.g. Fanelli 2009; John, Loewenstein, y Prelec
2012), sino que más bien son casos aislados (ver Abril Ruiz 2019, 23-128
para una revisión). En contraste, las QRP han demostrado ser más
prevalentes.

\subsection{Prácticas cuestionables de investigación
(QRP)}\label{pruxe1cticas-cuestionables-de-investigaciuxf3n-qrp}

Existen una serie de estudios que han intentado medir directamente la
prevalencia de estas prácticas a través de encuestas. (Fanelli 2009)
hizo un metaanálisis que tenía por objetivo sistematizar los resultados
de estudios que hasta esa fecha habían abordado las prácticas de
investigación desde encuestas cuantitativas. Los resultados mostraron
que un 1.97\% de investigadores había inventado datos al menos una vez
(FFP) y que un 33.7\% había realizado alguna vez una QRP como ``borrar
puntos de los datos basados en un sentimiento visceral''. Unos años más
tarde, (John, Loewenstein, y Prelec 2012) efectuaron otro estudio
similar, demostrando que un 36.6\% de quienes participaron alguna vez
habían practicado alguna QRP. En detalle, analizando los porcentajes
práctica a práctica los resultados muestran que el 50\% de los
psicólogos encuestados alguna vez reportaron selectivamente estudios que
apoyaran sus hipótesis; un 35\% alguna vez reportaron resultados
inesperados como esperados; y un 2\% alguna vez reportó datos falsos.
Este estudio ha sido replicado en Italia (Agnoli et~al. 2017), en
Alemania (Fiedler y Schwarz 2016) y en Brasil (Rabelo et~al. 2020).

Más recientemente, han emergido estudios similares en otras áreas
disciplinarias. Por ejemplo, a través de encuestas online (Makel et~al.
2021) logran constatar que existe una prevalencia considerable de las
QRP en el campo de investigación educacional. De los participantes de la
muestra, un 10\% admitió haber rellenado datos faltantes (NA's) y un
67\% señaló alguna vez haber omitido ciertos análisis de manera
intencional. Siguiendo el mismo método, en el campo de la investigación
comunicacional se ha encontrado evidencia parecida: 9\% de los
encuestados señala haber imputado datos faltantes sin reportarlo, un
34\% declaró alguna vez haber excluido casos extremos de forma
arbitraria y un 60\% señala no haber reportado análisis con variables
clave que no funcionaron. Del mismo modo, en los estudios cuantitativos
sobre criminología existe un uso extendido de las QRP: un 87\% de la
muestra de (Chin et~al. 2021) ha utilizado múltiples QRP, siendo el
reporte selectivo de resultados el más común (53\%). Por último, fuera
del área de las ciencias sociales, pero siguiendo la misma línea,
(Fraser et~al. 2018) también encuentran evidencia a favor de la
existencia de distintas QRP en el campo de la ecología y evolución.

Los estudios mencionados arriba corresponden a la evidencia existente
sobre la medición de QRP a través de encuestas. A modo de resumen, la
Tabla~\ref{tbl-tabqrp2} adaptada del trabajo de (Chin et~al. 2021)
agrupa la información que hemos revisado (y más), ordenándose por
prácticas, campos de estudio y los artículos correspondientes a cada
campo de estudio. Los números dentro de las casillas representan el
porcentaje de prevalencia de cada práctica reportado por los
participantes del estudio, y entre paréntesis el número de casos de la
muestra. Los guiones significan que este estudio no incluyó esa práctica
en el cuestionario.

\begin{verbatim}
Warning in styling_latex_scale(out, table_info, "down"): Longtable cannot be
resized.
\end{verbatim}

\begin{table}

\caption{\label{tbl-tabqrp2}Porcentajes de prevalencia por QRP según
distintos trabajos}

\centering{

\begingroup\fontsize{10}{12}\selectfont

\begin{longtabu} to \linewidth {>{\raggedright\arraybackslash}p{12em}>{\centering}X>{\centering}X>{\centering}X>{\centering}X>{\centering}X>{\centering}X>{\centering}X}
\toprule
Práctica & Psicología¹ & Psicología² & Psicología³ & Ecología⁴ & Evolución⁴ & Educación⁵ & Comunicación⁶\\
\midrule
Omitir estudios o variables no significativas & 46 (485) & 40 (217) & 55 (232) & - & - & 62 (783) & 60\\
Subreportar resultados & 63 (486) & 48 (219) & 22 (232) & 64 & 64 & 67 (871) & 64\\
Subreportar condiciones & 28 (484) & 16 (219) & 35 (232) & - & - & - & -\\
Muestreo selectivo & 56 (490) & 53 (221) & 22 (232) & 37 & 51 & 29 (806) & 23\\
Excluir datos selectivamente & 38 (484) & 40 (219) & 20 (232) & 24 & 24 & 25 (806) & 34\\
\addlinespace
Excluir covariables selectivamente & - & - & - & - & - & 42 (773) & 46\\
Cambiar análisis selectivamente & - & - & - & - & - & 50 (811) & 45\\
HARK & 27 (489) & 37 (219) & 9 (232) & 49 & 54 & 46 (880) & 46\\
Redondear valores p & 22 (499) & 22 (221) & 18 (232) & 27 & 18 & 29 (806) & 24\\
Mal orientar respecto a efectos sociodemográficos & 3 (499) & 3 (223) & 4 (232) & - & - & - & -\\
\addlinespace
Esconder problemas & - & - & - & - & - & 24 (889) & -\\
Esconder imputaciones & 1 (495) & 2 (220) & 1 (232) & 5 & 2 & 10 (898) & 9\\
\bottomrule
\end{longtabu}
\endgroup{}

}

\end{table}%

Como se observa en la Tabla~\ref{tbl-tabqrp2}, las encuestas sobre QRP
han incluido varias prácticas relativas al tratamiento y análisis de los
datos. No obstante, consideramos que existen tres términos que, en su
mayoría, logran sintetizar esta tabla y que están relacionados a la
transparencia en los diseños de investigación. Estas son: 1) los sesgos
de publicación, 2) el \emph{p-hacking} y el 3) \emph{HARKing}.

\textbf{Sesgo de publicación}

El \textbf{sesgo de publicación} ocurre cuando el criterio determinante
para que un artículo sea publicado es que sus resultados sean
significativos, en desmedro de la publicación de resultados no
significativos. Un ejemplo ilustrativo que usan (Christensen, Freese, y
Miguel 2019) para explicar esta práctica es el cómic \emph{xkcd}
titulado \emph{Significant}. Este cómic (Figura~\ref{fig-significant})
muestra cómo es que un hallazgo de investigación sufre del sesgo de
publicación. Al publicarse únicamente el resultado significativo e
ignorándose los otros 19 no significativos, cualquier lector tendería a
pensar que efectivamente las gominolas verdes causan acné, cuando
probablemente sea una coincidencia. El trabajo de (Rosenthal 1979) fue
de los primeros en llamar la atención respecto de esta práctica,
adjudicando el concepto de \emph{file drawer problem} (en español:
problema del cajón de archivos), el que hace alusión a los resultados
que se pierden o quedan ``archivados'' dentro de un cuerpo de
literatura. Desde ese estudio en adelante varios autores han contribuido
con evidencia empírica sobre el sesgo de publicación. Por ejemplo, el
estudio de (Franco, Malhotra, y Simonovits 2014) logra cuantificar esta
situación encontrando que los resultados nulos tienen un 40\% menos de
probabilidades de ser publicados en revistas científicas, en comparación
a estudios con resultados significativos. Es más, muchas veces los
resultados nulos ni siquiera llegan a ser escritos: más de un 60\% de
los experimentos que componen la muestra del estudio de (Franco,
Malhotra, y Simonovits 2014) nunca llegaron a ser escritos, en contraste
con el 10\% de resultados significativos.

\begin{figure}

\centering{

\includegraphics[width=0.75\linewidth,height=\textheight,keepaspectratio]{images/significant.png}

}

\caption{\label{fig-significant}Comic \emph{Significant} de xkcd}

\end{figure}%

El principal problema del sesgo de publicación es que puede impactar en
la credibilidad de cuerpos enteros de literatura. Por ejemplo, en
economía se han hecho varios metaanálisis que han buscado estimar el
sesgo de publicación en distintos cuerpos de literatura (e.g. Brodeur
et~al. 2016; Vivalt 2015; Viscusi 2014). Uno de los trabajos más
concluyentes es el de (Doucouliagos y Stanley 2013), quienes efectúan un
metaanálisis de 87 artículos de metaanálisis en economía. En este
trabajo encuentran que más de la mitad de los cuerpos de literatura
revisados sufren de un sesgo ``sustancial'' o ``severo''. Si bien en
economía se ha avanzado mucho en este tipo de estudios, también a partir
del desarrollo de distintos métodos de detección, se ha podido
diagnosticar el sesgo de publicación en importantes revistas en
sociología y ciencias políticas (Gerber y Malhotra 2008a, 2008b).

\textbf{P-hacking}

Otra práctica bastante cuestionada es el \emph{p-hacking}. El p-hacking
suele englobar muchas de las prácticas que vimos en un inicio,
especialmente las que refieren al manejo de datos: excluir datos
arbitrariamente, redondear un valor p, recolectar más datos posterior a
hacer pruebas de hipótesis, entre otras. Lo que tienen todas estas
prácticas en común y lo que define el p-hacking es que se da cuando el
procesamiento de los datos tiene por objetivo obtener resultados
significativos. Si el sesgo de publicación afecta la credibilidad de un
cuerpo de literatura, el \emph{p-hacking} afecta a la credibilidad de
los artículos mismos, ya que al forzar la significancia estadística la
probabilidad de que en realidad estemos frente a un falso positivo
aumenta. Un trabajo que da sustento a esta idea es el de (Simmons,
Nelson, y Simonsohn 2011), quienes calculan la posibilidad de obtener un
falso positivo (error Tipo I) de acuerdo con el nivel de \emph{grados de
libertad} que son utilizados por parte de los investigadores. El
resultado principal es que a medida que aumenta el uso de grados de
libertad, la posibilidad de obtener un falso positivo aumenta
progresivamente.

El \emph{p-hacking} también contribuye a sesgar cuerpos enteros de
literatura. Para diagnosticar esto se ha utilizado una herramienta
denominada \emph{p-curve}, la cual ``describe la densidad de los
\emph{p-values} reportados en una literatura, aprovechando el hecho de
que si la hipótesis nula no se rechaza (es decir, sin efecto), los
p-values deben distribuirse uniformemente entre 0 y 1'' (Christensen,
Freese, y Miguel 2019, 67.). De esta manera, en cuerpos de literatura
que no sufran de p-hacking, la distribución de valores p debería estar
cargada a la izquierda (siendo precisos, asimétrica a la derecha), en
cambio, si existe sesgo por p-hacking la distribución de p-values
estaría cargada a la derecha (asimetría a la izquierda). (Simonsohn,
Nelson, y Simmons 2014) proponen esta herramienta y la prueban en dos
muestras de artículos de la \emph{Journal of Personality and Social
Psychology (JPSP)}. Las pruebas estadísticas consistieron en confirmar
que la primera muestra de artículos (que presentaban signos de
p-hacking) estaba sesgada, en cambio la segunda muestra (sin indicios de
p-hacking), no lo estaba. Los resultados corroboraron las hipótesis,
esto es: los artículos que presentaban solamente resultados con
covariables resultaron tener una p-curve cargada a la derecha
(asimétrica a la izquierda).

\textbf{HARKing}

Por último, pero no menos importante, está la práctica del
\emph{HARKing}. El nombre es una nomenclatura en lengua inglesa:
\emph{Hypothesizing After the Results are Known}, que literalmente
significa establecer las hipótesis del estudio una vez que se conocen
los resultados (Kerr 1998). El principal problema de esta práctica es
que confunde los dos tipos de razonamiento que están a la base de la
ciencia: el exploratorio y el confirmatorio. El objetivo principal del
razonamiento exploratorio es plantear hipótesis a partir del análisis de
los datos, en cambio, el razonamiento confirmatorio busca plantear
hipótesis basado en teoría y contrastar esas hipótesis con datos
empíricos. Como señala (Brian A. Nosek et~al. 2018), cuando se confunden
ambos tipos de análisis y se hace pasar un razonamiento exploratorio
como confirmatorio se está cometiendo un sesgo inherente, ya que se está
generando un razonamiento circular: se plantean hipótesis a partir del
comportamiento de los datos y se confirman las hipótesis con esos mismos
datos

\section{¿Qué? y ¿Quién?}\label{quuxe9-y-quiuxe9n}

Preguntarse por el \textbf{qué} es básicamente cuestionarse sobre lo que
se está poniendo a disposición del escrutinio público. Dentro de la
literatura sobre ciencia abierta, son varios los contenidos que pueden
ser transparentados. Por ejemplo, la apertura de los datos es una
alternativa que, ciertamente, contribuye a hacer un estudio más
transparente (Miguel et~al. 2014). Al estar los datos abiertos al
público, la posibilidad de escrutinio por parte de las audiencias es
mayor. Otro ejemplo es la idea de \emph{disclosure}, la cual se refiere
a la declaración que hacen los investigadores respecto a si la
investigación cuenta con conflictos de interés o si está sesgada de
alguna manera. En base a esta idea, recientemente han emergido
propuestas cómo las de (O'Boyle, Banks, y Gonzalez-Mulé 2017) que,
además de la declaración de conflicto interés, llaman a adoptar
declaraciones relacionadas a las QRP. Básicamente declarar si se ha
cometido QRP o no durante el transcurso de la investigación. No
obstante, \textbf{la apertura pública de los diseños de investigación}
es, sin duda, la práctica que mejor representa la transparencia en el
último tiempo.

Como hemos revisado anteriormente, existe una emergente preocupación en
las ciencias sociales respecto a la credibilidad de los hallazgos
científicos y la influencia que pueden tener las QRP. En disciplinas
como la psicología la apertura de los diseños de investigación ha sido
la principal herramienta ante el diagnóstico que una parte considerable
de los hallazgos eran irreproducibles o irreplicables (Motyl et~al.
2017; Wingen, Berkessel, y Englich 2020; Camerer et~al. 2018). Ante un
diagnóstico similar en sociología (Breznau et~al. 2021), en ciencias de
la gestión (Byington y Felps 2017) y en economía (Gall, Ioannidis, y
Maniadis 2017), por dar algunos ejemplos, también se ha reconocido la
apertura del diseño de investigación como una herramienta importante
para avanzar hacia la transparencia. A razón de esto, es que nuestra
propuesta se focaliza en la apertura de los diseños de investigación y
el uso de preregistros.

Además de preguntarse por los contenidos que se transparentan, también
es relevante hacer un alcance respecto al \textbf{quién}. Esto es tener
en consideración quiénes van a recibir la información que sea
transparentada. En principio, los hallazgos de la investigación
científica son un bien público -o están pensados para serlo-. Si bien en
algunos países o áreas especificas del conocimiento el financiamiento de
la ciencia es por parte de privados, la investigación en ciencias
sociales en Chile se financia con los impuestos de los habitantes.
Organismos como la Agencia Nacional de Investigación y Desarrollo de
Chile (ANID), quienes están encargados de la administración de fondos
para la ciencia, tienen un rol público, por lo que los hallazgos no solo
sirven para alimentar el avance del conocimiento dentro de la academia,
sino que también son insumos relevantes para los tomadores de decisiones
y/o para cualquier ciudadano que se podría beneficiar de esos hallazgos.

En consiguiente, cualquier tipo de información que sea transparente es
un insumo más para que distintos tipos de audiencias puedan beneficiarse
de hallazgos más creíbles y robustos. Ya sean investigadores, tomadores
de decisión o la ciudadanía.

\section{¿Cómo?}\label{cuxf3mo}

Una buena forma de introducir el \textbf{cómo} adoptar la transparencia
son las \emph{Transparency and Openess} Promotion (TOP) Guidelines\_
(Guías para la Promoción de la Transparencia y la Accesibilidad). Las
\emph{TOP Guidelines} son una iniciativa del Centro para la Ciencia
Abierta (COS, por sus siglas en inglés) que busca fomentar la ciencia
abierta a partir de la adopción de distintos principios. Su objetivo es
que tanto los investigadores como las revistas científicas adhieran a
prácticas transparentes. Los principios van desde temas de citación,
hasta la replicabilidad (ver el detalle sobre esta propuesta en
\url{https://osf.io/9f6gx/}). Si bien esta es una iniciativa que incluye
principios que escapan del enfoque de este capítulo, es un buen punto de
partida ya que permite situar la transparencia en el diseño de
investigación y el uso de preregistros como parte de una red de
principios más grandes a instalar en la ciencia. Además, permite dar
cuenta que ya existe un conjunto de actores preocupados por temáticas de
ciencia abierta, que han reconocido la transparencia en los diseños de
investigación como un principio importante por adoptar. En lo que sigue,
revisaremos cómo los preregistros logran fomentar la transparencia en
los diseños de investigación.

\subsection{Preregistros}\label{preregistros}

Los preregistros son una marca temporal sobre las decisiones del diseño,
el método y el análisis de un artículo científico y se suelen hacer
antes del levantamiento de datos (Stewart et~al. 2020). Básicamente,
preregistrar un estudio implica que un grupo de investigadores dejarán
por escrito una pauta de investigación, la cual seguirán cuando
desarrollen la investigación, especialmente la recopilación y el
análisis de los datos. El objetivo del preregistro es reducir el margen
de flexibilidad que tienen los investigadores a la hora de analizar los
datos. El que exista una guía sobre qué hipótesis probar y qué análisis
realizar hace menos probable que los científicos recurran en prácticas
como el sesgo de publicación, p haking o el HARKing.

Como vimos, el sesgo de publicación se trata de publicar selectivamente
los resultados de investigación: resultados que no hayan sido
significativos, o hipótesis que ``no funcionaron'' simplemente se
omiten. Sin embargo, cuando existe un documento como un preregistro, el
cual deja estipulado claramente las hipótesis que deben ponerse a prueba
y los análisis que se emplearán para ello, se torna más difícil reportar
selectivamente los resultados. Dicho de otra forma, cuando existe una
pauta a la cual apegarse, la discrecionalidad en el reporte de los
resultados disminuye.

En el caso del p-hacking, el efecto del preregistro es parecido. El
p-hacking consiste en abusar de las pruebas estadísticas para obtener
resultados significativos. ``Abusar'' en el sentido de buscar toda vía
posible para obtener un valor p que confirme las hipótesis planteadas.
El hecho de preregistrar el plan de análisis y el procesamiento que se
le efectuará a las variables permite evitar este tipo de búsqueda
intencionada: como hay una guía que seguir, cualquier desviación debe
ser justificada. En esta misma línea, un preregistro evita el HARKing ya
que las hipótesis están previamente planteadas y no es posible
cambiarlas una vez que se han visto los resultados. En suma, el plantear
un registro \emph{a priori} de la investigación, disminuye la
flexibilidad que suele dar paso a las QRP.

Existen resquemores respecto del uso de preregistros de los que es
preciso hacerse cargo. Una de las principales preocupaciones es que el
uso de preregistros tendería a coartar la creatividad y la producción de
conocimiento exploratorio (Moore 2016). La lógica es que, como cada
parte de la investigación debe ser registrada detalladamente previo a la
recopilación, no queda espacio para la espontaneidad durante el análisis
de datos. Sin embargo, más que inhibir la investigación exploratoria, el
objetivo de preregistar un estudio es separar la investigación
confirmatoria (pruebas de hipótesis) y la exploratoria (generación de
hipótesis) (Brian A. Nosek et~al. 2018). En ese sentido, es posible la
investigación exploratoria bajo el modelo de preregistros, solo que hay
que especificarla como tal.

Una segunda creencia es que realizar un preregistro añade un nivel de
escrutinio mayor del necesario, es decir, como se conoce cada detalle,
la investigación se vuelve un blanco fácil de críticas. Sin embargo, la
situación es todo lo contrario (Moore 2016). Por ejemplo, haber
preregistrado un plan de análisis para una regresión logística binaria
con datos que originalmente eran ordinales hará más creíble los
resultados, ya que quienes evalúen la investigación tendrán pruebas de
que el nivel de medición no se cambió solo para obtener resultados
significativos. Una tercera idea en torno a los preregistros es que
conllevan una gran inversión de tiempo y energía. Si bien es cierto que
se añade un paso más al proceso de investigación, el avance en la
temática ha logrado que existan una variedad de plantillas que hacen el
proceso más rápido y eficiente. Desde una lógica racional, el tiempo que
toma este paso nuevo en la investigación es un costo bajo en contraste a
los beneficios que trae.

La característica más importante de los preregistros es que sean
elaborados previo al análisis de datos y deseablemente previo a su
recopilación. Este requisito es lo que permite asegurar la credibilidad
de los resultados, ya que si no hay datos que alterar, entonces las
probabilidades de que ocurra una QRP son básicamente nulas.
Generalmente, para las ciencias médicas o la psicología experimental
(disciplinas donde cada vez se usan más los preregistros), esto no suele
ser un problema ya que se utilizan diseños experimentales. Estos se
apegan al método científico clásico: se plantean hipótesis basadas en la
teoría, se diseña un experimento para probar esas hipótesis y luego se
recopilan y analizan los datos para ver si dan soporte a las hipótesis
planteadas. Sin embargo, en muchas disciplinas de las ciencias sociales
los diseños experimentales son una pequeña fracción del conjunto de la
literatura {[}e.g.~según (Card, DellaVigna, y Malmendier 2011) en 2010,
solo un 3\% de los artículos en las mejores revistas de economía eran
experimentales{]}, y lo que prima son los diseños observacionales con
datos secundarios. A diferencia de los estudios experimentales, en los
estudios con datos preexistentes se afecta el principal componente de
credibilidad de los preregistros. Nada puede asegurar que los datos
fueron analizados antes de la escritura del preregistro y que, por
ejemplo, las hipótesis se están planteando una vez conocidos los
patrones significativos (HARKing). De ahí que nace la pregunta sobre la
posibilidad de utilizar preregistros en estudios con datos
preexistentes.

En la literatura sobre preregistros se han discutido los desafíos que
implica preregistrar estudios que utilicen datos preexistentes (e.g.
Editors 2014). Existen posturas que proponen que, en realidad, no existe
una forma creíble para preregistrar este tipo de estudios (Christensen y
Miguel 2018). No obstante, otras posturas han profundizado en las
situaciones en las que aún es posible preregistrar estudios con datos
elaborados previamente. (Burlig 2018) propone tres escenarios donde el
preregistro de datos observacionales es valioso. El primero es,
básicamente, cuando los investigadores que diseñaron la investigación
generan sus propios datos, en este caso, los investigadores sí pueden
elaborar un preregistro previo a la recolección de datos. El segundo
escenario se da cuando se preregistra un estudio que tiene como objeto
de interés un suceso que aún no ha ocurrido, lo que se conoce como
estudios prospectivos. Por ejemplo, un grupo de investigadores puede
estar interesado en el efecto que tendrá la introducción de una ley en
las prácticas sociales, o el efecto de un tratado en las variaciones del
PIB. Para esos casos, el preregistro aún mantiene su validez original ya
que, si bien los datos ya existen, no es posible hacer los análisis
antes del preregistro porque el evento de interés no ha ocurrido. El
tercer escenario ocurre cuando los datos existen, pero no están abiertos
al público. En estos casos, es la accesibilidad lo que determina la
credibilidad del preregistro. Por ejemplo, el grupo de investigadores
que elaboraron los datos pueden establecer que serán accesibles con
previo contacto y que se solicitará un preregistro. Por ende, en orden
de analizar los datos, los investigadores interesados deberán elaborar
un preregistro para utilizar los datos.

Conforme a lo anterior, (Mertens y Krypotos 2019) proponen dos prácticas
para asegurar la credibilidad de un preregistro con datos secundarios.
Primero, que el grupo de investigadores que analiza los datos sea
distinto e independiente de quien propuso el diseño de investigación y
segundo, que el equipo realice sintaxis de análisis con datos simulados,
con tal de demostrar que las hipótesis ya existían previas a acceder a
los datos. Estas propuestas muestran que el requisito sobre la
temporalidad del preregistro puede.

La recomendación más transversal y a la vez simple para preregistrar
análisis con datos secundarios, es ser sincero y claro respecto a lo que
se ha hecho y lo que no (Lindsay 2020; Brian A. Nosek et~al. 2018). Por
ejemplo, reportar si es que se ha leído el reporte descriptivo sobre la
base de datos o si se tiene conocimiento de algún tipo de patrón de los
datos. Es preciso transparentar cualquier tipo de aproximación a los
datos previo haberlos analizado. Para lograr este nivel de detalle y ser
eficiente con los tiempos y la comunicación hacia otros investigadores,
es que existen plantillas predeterminadas para preregistrar distintos
tipos de artículos en diferentes situaciones. En la siguiente sección
presentaremos las plantillas más usadas.

\section{Herramientas para los diseños
transparentes}\label{herramientas-para-los-diseuxf1os-transparentes}

Nuestra principal apuesta para promover la transparencia en los diseños
de investigación son los preregistros, por lo que es a lo que le
dedicaremos más espacio de esta sección. De todos modos, también
revisaremos un par de herramientas que pueden complementar el uso de
preregistros. Esperamos que, posterior a esta sección, el lector pueda
ser capaz de utilizar estas herramientas para sus investigaciones.

\subsection{Plantillas de preregistro}\label{plantillas-de-preregistro}

En la práctica, preregistrar un artículo es básicamente sintetizar la
información importante sobre nuestra investigación en una plantilla
estandarizada y alojar ese documento en un lugar público. Por lo que el
primer paso para elaborar un preregistro es elegir la plantilla
correcta. Existen plantillas estandarizadas que están estructuradas de
tal forma que son útiles para preregistrar estudios de cualquier
disciplina, así como también existen plantillas dedicadas a una
disciplina o a situaciones particulares. En este apartado presentaremos
plantillas bajo cuatro categorías: a) plantillas genéricas, b)
plantillas para experimentos y previas a recolección de datos, c)
plantillas para datos existentes/secundarios, c) plantillas para
estudios de replicación y d) plantillas para \emph{registered reports}.
Además de proveer una descripción de cada una, orientaremos al lector
para elegir una plantilla de acuerdo al tipo de investigación que está
desarrollando.

El \emph{Open Science Framework} (OSF), que en español se traduce como
un marco de referencia para la ciencia abierta, actúa cómo una
herramienta y un repositorio que alberga las plantillas que trataremos
en esta sección (y más). Para ver todas las plantillas disponibles en
OSF hacer clic en el siguiente enlace:
\url{https://osf.io/zab38/wiki/home/}.

\subsubsection{Plantillas Genéricas}\label{plantillas-genuxe9ricas}

Las plantillas genéricas son aquellas que pueden ser utilizadas
independientemente de las características del estudio, ya que suelen
incluir campos que abordan distintos tipos de situaciones (e.g.~si un
estudio es experimental o observacional). Acá revisaremos dos plantillas
genéricas, la plantilla estándar de \emph{AsPredicted}, y la plantilla
estándar de OSF.

\textbf{AsPredicted}

La plantilla de AsPredicted es quizás una de las más conocidas para
hacer preregistros, dado que está estandarizada y puede ser utilizada en
cualquier disciplina. Recomendamos utilizarla cuando lo que se busque es
optimizar tiempo y energías. La plantilla cuenta solamente con nueve
preguntas clave que aplican a cualquier tipo de artículo. Esta plantilla
la podemos encontrar tanto en OSF, como en la página de AsPredicted, en
este caso, mostraremos cómo es el proceso en la página original.

Partimos por entrar a la página de AsPredicted, donde veremos algo
similar a la Figura~\ref{fig-asp}. Acá se nos da la opción de crear un
preregistro, de ver los que ya hemos hecho (si es que ese es el caso) y
también una breve descripción de AsPredicted. A grandes rasgos, la
página nos dice que AsPredicted es una plataforma que busca facilitar el
preregistro de estudios por parte de los investigadores a través de
nueve simples preguntas. La página genera un documento .pdf y una URL
asociada. También, cuenta cómo funciona el preregistro. Básicamente, un
autor elabora un preregistro de un estudio y los coautores reciben un
mail para aprobar ese preregistro. Una vez aprobado por todos los
autores, el preregistro queda alojado en la plataforma de manera
privada, y no cambia hasta que un autor decida hacerlo público. Además,
en caso de que el estudio entre en revisión por pares, se puede enviar
una versión anónima del preregistro. Por último, nos entrega una
recomendación sobre qué hacer en el caso de que el proceso de
investigación no haya podido apegarse totalmente a lo predicho.

\begin{figure}

\centering{

\includegraphics[width=1\linewidth,height=\textheight,keepaspectratio]{images/prereg1.png}

}

\caption{\label{fig-asp}Botón para comenzar un preregistro}

\end{figure}%

Para elaborar un preregistro debemos hacer click en el rectángulo azul
que dice \emph{Create}. Una vez hecho eso, nos pedirá una dirección de
email para continuar. Cuando ingresemos un email, nos enviará un enlace
al email que hayamos ingresado, con ese enlace podremos comenzar el
preregistro. Una vez hayamos entrado en el enlace, veremos la plantilla
de preregistro. Lo primero que aparece es una sección donde debemos
escribir los emails de los autores colaboradores del estudio. También,
nos da la opción de añadir otros emails además del que hemos
introducido. Una vez pasada esta parte, ya nos encontramos con las
preguntas del preregistro, que son las siguientes:

\begin{enumerate}
\def\labelenumi{\arabic{enumi})}
\tightlist
\item
  Recogida de datos. ¿Se han recogido ya datos para este estudio?
\item
  Hipótesis. ¿Cuál es la pregunta principal que se plantea o la
  hipótesis que se pone a prueba en este estudio?
\item
  Variable dependiente. Describa la(s) variable(s) dependiente(s) clave
  especificando cómo se medirán.
\item
  Condiciones. ¿Cuántos y qué condiciones se asignarán a los
  participantes?
\item
  Análisis. Especifique exactamente qué análisis realizará para examinar
  la pregunta/hipótesis principal.
\item
  Valores atípicos y exclusiones. Describa exactamente cómo se definirán
  y tratarán los valores atípicos, así como su(s) regla(s) precisa(s)
  para excluir las observaciones.
\item
  Tamaño de la muestra. ¿Cuántas observaciones se recogerán o qué
  determinará el tamaño de la muestra?
\item
  Otros. ¿Hay algo más que quiera preinscribir?
\item
  Nombre. Poner un título a este preregistro de AsPredicted Finalmente.
  A efectos de registro, indíquenos el tipo de estudio que está
  preinscribiendo.
\end{enumerate}

Las preguntas son bastante autoexplicativas, pero no está de más
entregar algunos detalles adicionales. En la pregunta de recolección de
datos, las opciones son tres: ``Sí, se han recolectado datos'', ``No, no
se han recolectado datos'' y ``Es complicado''. Es importante mencionar
que, en esta plantilla, la respuesta de que se han recolectado datos no
es válida, por lo que si se está llevando a cabo un estudio con datos
secundarios hay que responder ``Es complicado'' y en la pregunta 8 de la
plantilla especificar por qué este preregistro sigue siendo válido pese
a que los datos son preexistentes. Otro detalle importante es que cada
pregunta está pensada para ser respuesta en aproximadamente una oración.
Esta plantilla tiene el objetivo de ser lo más eficiente posible, por lo
que, en general, se recomienda que todo el documento no pase de los 3200
caracteres. Otro detalle que especificar es que la pregunta acerca del
tipo de estudio que se está preregistrando también es semicerrada,
tenemos las opciones de: ``Proyecto de clase'', ``Experimento'',
``Encuesta'', ``Estudio observacional'' y ``Otro''. Es responsabilidad
de quien hace el preregistro el seleccionar la opción que más se asemeje
a su situación. Por último, es importante señalar que el preregistro, al
menos en la página de AsPredicted, solo puede ser rellenado en inglés,
por lo que en caso de utilizar otro idioma solicitará traducirlo.

\textbf{OSF Estándar}

La plantilla estándar de OSF también es de carácter general y busca
abarcar distintas situaciones y tipos de estudios. La diferencia que
tiene con la plantilla de AsPredicted es que contiene más preguntas y
más detalladas. Por ejemplo, contiene preguntas relativas a en qué
momento específico se está efectuando el preregistro (i.e antes de
recopilar los datos, antes de cualquier observación humana a los datos,
antes de acceder a los datos, entre otros) o si se harán
transformaciones de las variables y su nivel de medición. La plantilla
puede verse en este enlace:
\url{https://osf.io/preprints/metaarxiv/epgjd/}. A continuación, veremos
el inicio del proceso para registrar en OSF, el cual sirve tanto para la
plantilla estándar como para otras que veremos más adelante.

El primer paso es acceder a la sección específica de preregistros de la
página de OSF, la cual se encuentra en el siguiente enlace:
\url{https://osf.io/prereg/} (\emph{para usar este servicio es necesario
tener una cuenta}). Si entramos al enlace, la apariencia de la página
será algo similar a la Figura~\ref{fig-osfprereg1}. Seleccionemos
\emph{Start a new preregistration}, le damos un nombre al proyecto y
hacemos click en \emph{Continue}. En la página siguiente, podemos ver
que hemos creado un proyecto nuevo en OSF, el cual nos da la opción de
preregistrarlo haciendo click en el botón \emph{New registration}.

\begin{figure}

\centering{

\includegraphics[width=1\linewidth,height=\textheight,keepaspectratio]{images/prereg2.png}

}

\caption{\label{fig-osfprereg1}Opciones para comenzar un preregistro en
OSF}

\end{figure}%

En la Figura~\ref{fig-osfprereg2} podemos ver dos cosas. Primero, la
descripción de lo que está haciendo OSF al comenzar un nuevo
preregistro, lo que en pocas palabras es una versión no modificable del
proyecto al momento que hacemos el preregistro. En otras palabras, es
una versión ``congelada'' del proyecto en OSF. Segundo, también se
aprecia una serie de opciones para preregistrar, estas son las
plantillas que mencionamos anteriormente. OSF nos ofrece distintas
plantillas de acuerdo con el carácter que tiene nuestro estudio.

\begin{figure}

\centering{

\includegraphics[width=1\linewidth,height=\textheight,keepaspectratio]{images/prereg3.png}

}

\caption{\label{fig-osfprereg2}Opciones de plantilla de pre-registro}

\end{figure}%

El primer paso es escoger la plantilla de acuerdo a la situación en la
que nos encontremos. En este caso, seleccionamos la plantilla estándar
de OSF. Una vez hayamos hecho eso, será necesario llenar los metadatos
del estudio (e.g. Figura~\ref{fig-osfprereg3}). Esta sección es
transversal a todas las plantillas y consiste en registrar el título,
descripción, contribuyentes, entre otras cosas que ayudan a identificar
el proyecto.

\begin{figure}

\centering{

\includegraphics[width=0.75\linewidth,height=\textheight,keepaspectratio]{images/prereg4.png}

}

\caption{\label{fig-osfprereg3}Ejemplo de campos de metadatos para
rellenar}

\end{figure}%

Una vez hayamos rellenado los campos correspondientes a los metadatos,
podemos rellenar la plantilla de preregistro. Con los campos rellenados
podremos proceder a preregistrar nuestro proyecto.

\subsubsection{Plantillas para diseños experimentales y previas a
recolección de
datos}\label{plantillas-para-diseuxf1os-experimentales-y-previas-a-recolecciuxf3n-de-datos}

Considerando que el sentido original de un preregistro es que sea
elaborado previo a la recolección y análisis de los datos, en principio
cualquier plantilla genérica podría servir. Por ejemplo, la plantilla
estándar de OSF ofrece preguntas detalladas que refieren al uso de
diseños experimentales (i.e.~cuáles son las condiciones de tratamiento o
quiénes están al tanto de las manipulaciones experimentales del
estudio). Sin embargo, OSF ofrece una vía alternativa para preregistrar
estudios previo a la recolección de datos a través de la plantilla
``OSF-Standard Pre-Data Collection pre-registration''. Esta plantilla es
un complemento de la plantilla estándar y se utiliza solamente si es que
el preregistro original está archivado en un documento. Esto quiere
decir que, en el caso de que el equipo de investigación no haya usado el
flujo de preregistro de OSF, sino que ya cuenta con un archivo que
alberga el preregistro del estudio, entonces puede escoger esta
plantilla de complemento para no efectuar todo el proceso de nuevo.

Esta plantilla agrega algunas preguntas cruciales sobre la recopilación
de datos, siendo las siguientes:

\begin{enumerate}
\def\labelenumi{\arabic{enumi}.}
\tightlist
\item
  ¿Ha comenzado la recogida de datos para este proyecto?
\end{enumerate}

\begin{itemize}
\tightlist
\item
  No, la recopilación de datos no ha comenzado
\item
  Sí, la recopilación de datos está en curso o se ha completado
\end{itemize}

\begin{enumerate}
\def\labelenumi{\arabic{enumi}.}
\setcounter{enumi}{1}
\tightlist
\item
  ¿Ha mirado los datos?
\end{enumerate}

\begin{itemize}
\tightlist
\item
  Sí
\item
  No
\end{itemize}

\begin{enumerate}
\def\labelenumi{\arabic{enumi}.}
\setcounter{enumi}{2}
\tightlist
\item
  Otros comentarios
\end{enumerate}

\subsubsection{Plantillas para datos
secundarios}\label{plantillas-para-datos-secundarios}

En el último tiempo se ha comenzado a considerar cómo es que
investigadores de las ciencias sociales empíricas que trabajan con datos
preexistentes pueden, de todas maneras, preregistrar sus estudios para
asegurar la credibilidad de sus hallazgos. Si bien OSF cuenta con una
plantilla tipo para este tipo de situaciones (ver en
\url{https://osf.io/x4gzt/}), nosotros recomendamos la plantilla de
Mertens y Krypotos (2019) dada la simpleza y exhaustividad de sus
preguntas. La plantilla de Mertens y Krypotos (2019) se puede ver en el
siguiente enlace: \href{https://osf.io/p26rq/}{ttps://osf.io/p26rq/}.

Esta plantilla está dirigida a investigadores que usen datos
secundarios. Hace más simple el proceso de preregistro al enfocarse en
las preguntas que se asemejan realmente a la situación del investigador
que trabaja con grandes encuestas o datos administrativos, lo cual es
muchas veces el caso en disciplinas como sociología, ciencias políticas
o economía.

Esta plantilla cuenta con las siguientes diez preguntas:

\begin{enumerate}
\def\labelenumi{\arabic{enumi}.}
\item
  ¿Cuál es la hipótesis que se investigará?
\item
  ¿Cómo se operacionalizarán las variables cruciales?
\item
  ¿Cuál es la fuente de los datos incluidos en los análisis?
\item
  ¿Cómo se obtendrán estos datos?
\item
  ¿Hay algún criterio de exclusión de los datos?
\item
  ¿Cuáles son los análisis estadísticos previstos?
\item
  ¿Cuáles son los criterios para confirmar y desconfirmar las hipótesis?
\item
  ¿Se han validado los análisis con un subconjunto de datos? En caso
  afirmativo, facilite los archivos de datos y la sintaxis de los
  mismos.
\item
  ¿Qué se sabe sobre los datos que podría ser relevante para las
  hipótesis probadas?
\item
  Facilite un breve calendario para los diferentes pasos del
  prerregistro.
\end{enumerate}

Como podemos ver, además de los campos que se pueden encontrar en
cualquier plantilla como la especificación de las hipótesis del estudio
o los criterios de exclusión de datos, esta plantilla hace preguntas
relativas al nivel de conocimiento previo de los datos. Por ejemplo, en
la pregunta 4 solicita especificar la fuente de los datos. Como son
datos secundarios, esto implica detallar cómo se accederá o serán
solicitados los datos: si es que es necesario rellenar algún formulario
o contactar a alguien en específico para obtener acceso. También, en la
pregunta 9 se solicita describir cualquier conocimiento de algún patrón
en los datos que sea relevante para la pregunta de investigación
(e.g.~la media o la mediana de una variable). Estos son ejemplos de
preguntas que hacen esta plantilla útil para los investigadores que
trabajen con datos preexistentes y que quieran preregistrar sus
estudios.

\subsubsection{Plantillas para estudios de
replicación}\label{plantillas-para-estudios-de-replicaciuxf3n}

Tratar en extenso qué es y cómo se hace un estudio de replicación escapa
de los objetivos de este capítulo, sin embargo, es una práctica de la
ciencia abierta que no podemos dejar fuera ya que también cuenta con una
plantilla para prerregistro. Replicar un estudio significa regenerar los
hallazgos de un estudio, siguiendo sus hipótesis y plan de análisis,
pero con datos distintos.

OSF ofrece dos plantillas para investigadores que tienen el objetivo de
replicar un estudio. La primera es una plantilla para estudios de
replicación previo a su ejecución, la cual incluye una serie de
preguntas relativas al estudio original, por ejemplo, qué efecto se está
buscando replicar, por qué es importante replicarlo o en qué área
geográfica fue conducido el estudio original. Esta plantilla se puede
ver en el siguiente enlace: \url{https://osf.io/4jd46/}. La segunda
plantilla es para estudios de replicación que ya han sido completados.
En este caso, las pvreguntas están relacionadas a los hallazgos del
estudio de replicación, por ejemplo, cuál fue el efecto obtenido en la
replicación y si se considera que este efecto logra replicar los
resultados originales o no. El enlace para esta plantilla se puede
encontrar en este enlace: \url{https://osf.io/9rp6j/}.

\subsubsection{\texorpdfstring{Plantillas para \emph{registered
reports}}{Plantillas para registered reports}}\label{plantillas-para-registered-reports}

Cómo veremos en la siguiente sección, los \emph{registered reports} son
un modelo alternativo de publicación donde el diseño del estudio pasa
por un proceso de revisión por pares, a diferencia del modelo
tradicional donde el documento que pasa por revisión por pares es el
artículo finalizado. La plantilla que ofrece OSF es para estudios que
han sido aceptados en la revisión por pares en una revista académica que
cuenta con el modelo de \emph{registered reports}. El objetivo de esta
plantilla es poder dejar un registro público sobre el artículo en
proceso, por lo que la plantilla de preregistro consta de, básicamente,
una sección para el título del artículo, la fecha de aceptación, el
manuscrito y archivos complementarios. Esta plantilla se puede ver en
este enlace: \url{https://osf.io/gm36s/}.

\section{Otras herramientas}\label{otras-herramientas}

Sí bien los preregistros son una de las herramientas que más ha ido
tomando protagonismo para promover la transparencia, existen otras.
Específicamente, queremos mencionar dos de ellas: el modelo de
\emph{registered reports} (en español, informes registrados) y la
\emph{transparency checklist} (en español, lista de transparencia).

\subsection{Registered Reports}\label{registered-reports}

El modelo de \emph{registered reports} es una alternativa al modelo
tradicional de publicación. Consiste en que el artículo atraviesa una
revisión por pares en etapas tempranas de la investigación,
específicamente \textbf{previo a la recolección de datos}. Esta práctica
tiene por objetivo que el estudio sea evaluado por su innovación y
calidad del diseño de investigación, más que por resultados impactantes
(Chambers et~al. 2015). Además, busca dejar sin efecto prácticas como el
sesgo de publicación, p-hacking y HARKing, ya que no solamente existe
una marca temporal que avala el diseño de investigación del estudio
(como es el caso de un preregistro), sino que también existe un grupo de
científicos anónimos que están de acuerdo con que el artículo es un
aporte al conocimiento (Chambers 2013; Brian A. Nosek y Lakens 2014;
Marsden et~al. 2018).

Los \emph{registered reports} tienen dos características principales
(Marsden et~al. 2018). Primero, un manuscrito con la justificación del
estudio, lo que incluye una introducción, revisión de antecedentes y una
pregunta de investigación, dando la posibilidad de una
\textbf{aceptación preliminar} (IPA, por sus siglas en inglés \emph{In
principle acceptance}). La segunda característica es que el IPA no puede
revocarse en base a los resultados del estudio, esto evita que trabajos
con resultados no significativos no sean publicados y así combatir el
sesgo de publicación. Relacionado a ambas características, los informes
registrados pasan por dos etapas de revisión, la primera es la del
manuscrito, siendo este el determinante de se el estudio se acepta o no,
y la segunda revisión que se da posterior a terminar la recolección y
análisis de datos. El modelo, en comparación al sistema tradicional de
publicaciones, se puede ver en la Figura~\ref{fig-rr}.

\begin{figure}

\centering{

\includegraphics[width=0.75\linewidth,height=\textheight,keepaspectratio]{images/rr.png}

}

\caption{\label{fig-rr}Método convencional y de registered reports para
publicación científica}

\end{figure}%

Para enviar un artículo bajo el modelo de \emph{registered reports},
primero se debe tener en conocimiento cuáles son las revistas que
cuentan con este tipo de revisión. El Centro para la Ciencia Abierta
cuenta con una lista actualizada de revistas
\href{https://www.cos.io/initiatives/registered-reports}{aquí}. Una vez
escogida una revista, el proceso no es tan distinto al método
convencional, en el sentido de que los investigadores envían su
manuscrito con la justificación del estudio y este puede ser aceptado o
rechazado por el editor, ya sea directamente o después de la corrección
de algunos comentarios. Una vez se cuenta con el IPA y se efectúa la
revisión en la segunda etapa, los revisores se aseguran de que el
estudio ha seguido el plan de análisis inicialmente planteado y si sus
conclusiones tienen sentido de acuerdo a los datos y resultados
obtenidos, así como también que toda desviación del plan original sea
debidamente justificada (Stewart et~al. 2020). Desviaciones muy
sustanciales y/o que no sean debidamente justificadas pueden conllevar
el rechazo del artículo, aunque puede seguir el proceso en el método
convencional de publicación (Stewart et~al. 2020).

\subsection{Transparency Checklist}\label{transparency-checklist}

La \emph{transparency checklist} es una herramienta complementaria
elaborada por Aczel et~al. (2020) que busca acompañar el proceso de
reportar un estudio, contribuyendo a que estos sean más transparentes.
Esta lista ha sido elaborada específicamente para investigadores de las
ciencias sociales y del comportamiento que trabajan con datos primarios,
aunque puede ser útil para otros enfoques y disciplinas. La lista consta
de 36 items divididos en cuatro categorías: preregistro, método,
resultados y discusiones y accesibilidad de datos, donde cada ítem
refiere a alguna característica de un reporte transparente, preguntando
si ha sido efectuada o no, es decir, las respuestas posibles de cada
ítem son ``Sí'', ``No'' y ``N/A''. Existe una versión más corta de 12
items, los cuales son los siguientes:

\textbf{\emph{Sección de preregistro:}}

\begin{enumerate}
\def\labelenumi{(\arabic{enumi})}
\item
  Antes de analizar el conjunto completo de datos, se publicó un
  prerregistro con sello de tiempo en un registro independiente de
  terceros para el plan de análisis de datos.
\item
  El estudio fue registrado\ldots{}
\end{enumerate}

\begin{itemize}
\tightlist
\item
  antes de que cualquier dato fuera recolectado
\item
  después de que algunos datos fueron recolectados, pero antes de
  explorarlos
\item
  después de que todos los datos fueron recolectados, pero antes de
  explorarlos
\item
  después de explorar los datos, pero antes de que cualquier análisis
  estadístico fuera efectuado
\item
  después de efectuar algunos análisis estadísticos, pero no todos
\item
  en otro momento, explicar:
\end{itemize}

\begin{enumerate}
\def\labelenumi{(\arabic{enumi})}
\setcounter{enumi}{2}
\tightlist
\item
  El análisis estadístico previsto para cada pregunta de investigación
  (esto puede requerir, por ejemplo, información sobre la unilateralidad
  de las pruebas, los criterios de inferencia, las correcciones para
  pruebas múltiples, los criterios de selección de modelos, las
  distribuciones previas, etc.).
\end{enumerate}

\textbf{\emph{Sección de método}}

El manuscrito describe completamente\ldots{}

\begin{enumerate}
\def\labelenumi{(\arabic{enumi})}
\setcounter{enumi}{3}
\item
  la justificación del tamaño de la muestra utilizado (por ejemplo, un
  análisis de potencia a priori).
\item
  el diseño, los procedimientos y los materiales del estudio para
  permitir una réplica independiente.
\item
  las medidas de interés (por ejemplo, la amabilidad) y sus
  operacionalizaciones (por ejemplo, un cuestionario que mide la
  amabilidad).
\item
  ¿algún cambio en el prerregistro (como cambios en los criterios de
  elegibilidad, en los límites de pertenencia al grupo o en los
  procedimientos experimentales)?
\end{enumerate}

\textbf{\emph{Sección de resultados y discusión}}

El manuscrito\ldots{}

\begin{enumerate}
\def\labelenumi{(\arabic{enumi})}
\setcounter{enumi}{7}
\tightlist
\item
  distingue explícitamente entre la ``confirmación'' (es decir,
  preestablecido) y ``exploratorio'' (es decir, no preestablecidos).
\end{enumerate}

\textbf{\emph{Sección de disponibilidad de datos, código y materiales}}

Se han hecho públicas las siguientes\ldots{}

\begin{enumerate}
\def\labelenumi{(\arabic{enumi})}
\setcounter{enumi}{8}
\item
  los datos (procesados) en los que se han basado los análisis del
  manuscrito.
\item
  todo el código y el software (que no esté protegido por derechos de
  autor).
\item
  todas las instrucciones, los estímulos y los materiales de las pruebas
  (que no estén protegidos por derechos de autor).
\item
  El manuscrito incluye una declaración sobre la disponibilidad y
  localización de todos los elementos de la investigación, incluidos los
  datos, materiales y códigos pertinentes para su estudio.
\end{enumerate}

Tanto la versión completa de 36 items, como la recortada de 12 están
disponibles para rellenar en línea.
\href{http://www.shinyapps.org/apps/TransparencyChecklist/}{Aquí} se
puede encontrar la lista online, es una aplicación de uso simple, además
que permite generar el reporte final de manera automática.

\section*{Anexos}\label{anexos}
\addcontentsline{toc}{section}{Anexos}

\markright{Anexos}

\begin{table}

\caption{\label{tbl-tabqrp}Algunas situaciones de QRP}

\centering{

\begingroup\fontsize{10}{12}\selectfont

\begin{tabu} to \linewidth {>{\raggedright\arraybackslash}p{5cm}>{\raggedright\arraybackslash}p{5cm}}
\toprule
Dimensión & Práctica\\
\midrule
 & Ignorar ciertos aspectos de los requerimientos de las personas participantes.\\
\cmidrule{2-2}
 & Pasar por alto el uso de datos cuestionables o de interpretaciones cuestionables que otros hacen.\\
\cmidrule{2-2}
 & Cambiar partes de un estudio como respuesta a la presión de una fuente de financiación\\
\cmidrule{2-2}
 & Eliminar observaciones de los análisis basados en la intuición de que eran inexactos.\\
\cmidrule{2-2}
 & Redondear un valor p (por ejemplo, reportar que un p-value de 0,054 es menor a 0,05)\\
\cmidrule{2-2}
 & Eliminación, adición o alteración de datos después de pruebas de hipótesis.\\
\cmidrule{2-2}
 & Supresión selectiva o adición de variables.\\
\cmidrule{2-2}
\multirow{-8}{5cm}[3.5\dimexpr\aboverulesep+\belowrulesep+\cmidrulewidth]{\raggedright\arraybackslash Diseño y procesamiento} & Invertir la dirección o reformular hipótesis para respaldar los datos\\
\cmidrule{1-2}
 & Ampliar de manera innecesaria la bibliografía de un estudio.\\
\cmidrule{2-2}
 & Tergiversar los logros de la investigación.\\
\cmidrule{2-2}
 & Exagerar la importancia y la relevancia práctica de los resultados.\\
\cmidrule{2-2}
 & Retener detalles de la metodología de investigación (e.g. no reportar todas las variables dependientes de un estudio)\\
\cmidrule{2-2}
 & Retener resultados de la investigación (e.g. no presentar datos que contradicen una propia investigación previa).\\
\cmidrule{2-2}
 & Establecer publicaciones o brindar apoyo a publicaciones que no cumplen el proceso de control de calidad de la investigación\\
\cmidrule{2-2}
 & Publicar los mismos datos o resultados en dos o más publicaciones.\\
\cmidrule{2-2}
 & Selectivamente reportar estudios que "funcionaron".\\
\cmidrule{2-2}
 & Reportar hallazgos inesperados como previstos desde el principio.\\
\cmidrule{2-2}
\multirow{-10}{5cm}[4.5\dimexpr\aboverulesep+\belowrulesep+\cmidrulewidth]{\raggedright\arraybackslash Redacción, reporte y publicación} & Afirmar que los resultados no se ven afectados por variables demográficas cuando uno no está realmente seguro (o sabe que lo hacen).\\
\cmidrule{1-2}
 & Manipular la autoría o denigrar el papel de otros investigadores en las publicaciones.\\
\cmidrule{2-2}
 & Asignar inapropiadamente los crédios de autoría.\\
\cmidrule{2-2}
 & Volver a publicar partes sustanciales de publicaciones propias anteriores, incluidas las traducciones, sin citar debidamente el original ("autoplagio").\\
\cmidrule{2-2}
 & Citar de forma selectiva para mejorar los propios resultados o para complacer a los editores, los revisores o los colegas.\\
\cmidrule{2-2}
 & Usar ideas de otros sin obtener permiso o dar crédito\\
\cmidrule{2-2}
\multirow{-6}{5cm}[2.5\dimexpr\aboverulesep+\belowrulesep+\cmidrulewidth]{\raggedright\arraybackslash Citación y autoría} & Uso no autorizado de información confidencial en relación con la propia investigación.\\
\cmidrule{1-2}
 & Permitir que los patrocinadores pongan en peligro la independencia en el proceso de investigación con el fin de introducir sesgos.\\
\cmidrule{2-2}
 & Acusar a un investigador de conducta indebida u otras infracciones de forma maliciosa.\\
\cmidrule{2-2}
 & Retrasar u obstaculizar inadecuadamente el trabajo de otros investigadores.\\
\cmidrule{2-2}
 & Emplear la experiencia profesional propia para alentar a que se incumpla la integridad de la investigación.\\
\cmidrule{2-2}
 & Ignorar supuestos incumplimientos de la integridad de la investigación cometidos por terceros o encubrir reacciones inadecuadas a conductas indebidas.\\
\cmidrule{2-2}
 & No divulgar adecuadamente la participación en empresas cuyos productos se basan en la investigación de uno.\\
\cmidrule{2-2}
\multirow{-7}{5cm}[3\dimexpr\aboverulesep+\belowrulesep+\cmidrulewidth]{\raggedright\arraybackslash Relaciones con otros} & Relaciones con estudiantes, sujetos de investigación o clientes que pueden ser interpretadas como cuestionables.\\
\cmidrule{1-2}
 & \\
\bottomrule
\end{tabu}
\endgroup{}

}

\end{table}%

\bookmarksetup{startatroot}

\chapter{Datos abiertos}\label{datos-abiertos}

Si alguna vez has intentado conseguir datos de una investigación social
producidos por otros profesionales, probablemente sabes cuán difícil es.
Si no lo has intentado, es algo como Mickey Mouse intentando pasar a la
siguiente habitación, puesto que nos enfrentaremos consecutivamente a
una barrera tras otra.

Imagine que el investigador responsable de un Fondecyt, Dr.~González,
desea estudiar la calidad de vida de los inmigrantes. En su revisión
bibliográfica encuentra una investigación con datos que podrían ayudar
al desarrollo del proyecto. Logra conseguir el contacto de la
investigadora responsable y esta pese a su voluntad de compartir los
datos le advierte que tardara unas semanas en ello, puesto que está muy
ocupada y no sabe la ubicación actual de los mismos. Pese a perder
valioso tiempo de su proyecto, el Dr.~González logra acceder a la base
de datos, no obstante, esta se encuentra dispersa en distintas hojas de
cálculo. Al abrirlos e intentar conectarlos, se da cuenta que no
comprende el significado de cada variable, por lo cual requiere un libro
de códigos de los datos. Vuelve a contactar a la investigadora, le
solicita el cuestionario y el significado de cada variable.
Lamentablemente, el equipo que creo la base de datos señala que no posee
documentación sobre el significado de cada variable. Sin más, al no
poder utilizar debidamente los datos producidos por la investigación
anterior, el Dr.~González se resigna, y decide volver a invertir
recursos en una encuesta, la que probablemente tampoco quede a
disposición de futuros investigadores.

En suma, al intentar conseguir datos de otras investigaciones nos
enfrentamos a diversas barreras, como pueden ser la posibilidad de
contactar al equipo, la voluntad del equipo de compartir sus datos, la
calidad de la documentación de los datos, los formatos que pueden estar
en versiones pagadas, entre otras.

¿Qué podemos hacer los investigadores sociales para evitar estas
situaciones? Para ello, el movimiento de la ciencia abierta incentiva a
los investigadores a publicar sus datos de investigación (Open Data).
Más aún, se señala la importancia de publicarlos cumpliendo estándares
internacionales sobre los datos y la documentación, de modo tal que
cualquier investigador sea capaz de reutilizar los datos sin la
necesidad de contactarse con el equipo que los produjo y con la
información suficiente para reconocer o citar el aporte del equipo
productor de los datos.

A nivel internacional existe una amplia preocupación por democratizar el
conocimiento. En esta línea, la Organización para la Cooperación y el
desarrollo Económicos (OCDE 2020), ha incluido dentro de las condiciones
para mantenerse en la organización la obligación de que cada país
incluya políticas de ciencia abierta, transparentando los procesos y
resultados, para todas las investigaciones financiadas públicamente. A
nivel nacional tanto el Instituto Nacional de Estadísticas (INE),
dependiente del Ministerio de Economía, como la Agencia Nacional de
Investigación y Desarrollo (ANID 2020a) dependiente del Ministerio de
Ciencia, han incorporado políticas y prácticas propias de la ciencia
abierta.

En vista de estos cambios realizados por las instituciones chilenas de
investigación, resulta evidente que aprender prácticas de ciencia
abierta sobre cómo compartir los análisis y los datos producidos se
vuelve una necesidad para los investigadores Chilenos que trabajan en
instituciones estatales o financiadas con fondos públicos. En vista de
esta necesidad el objetivo de este documento es facilitar los
conocimientos necesarios para abrir la información producida por las
investigaciones.

\section{Ventajas de la apertura de datos de investigación
social}\label{ventajas-de-la-apertura-de-datos-de-investigaciuxf3n-social}

Las ventajas de la apertura de datos son varias. A continuación nos
centramos en los tres puntos que consideramos más relevantes: las
características éticas, la contribución a la calidad y la eficiencia
científica y otros incentivos personales importantes de considerar.

\begin{itemize}
\item
  \textbf{Ética}:

  \begin{itemize}
  \item
    Es justo que el público general tenga acceso a los datos producidos,
    especialmente cuando estos son producidos con fondos públicos (Bueno
    2017). Cabe destacar que Chile gasta aproximadamente \$668.551 MM de
    pesos en investigaciones científicas, lo cual equivale a un 0,35\%
    del producto interno bruto.
  \item
    La apertura de los datos fomenta la ética investigativa y la
    confiabilidad, reduce el fraude y aumenta el valor de las ciencias
    sociales para los políticos y el público en general (Breznau 2019).
  \end{itemize}
\item
  Calidad y eficiencia científica:

  \begin{itemize}
  \item
    Dejar la base de datos a libre disposición permite hacer
    evaluaciones sobre la rigurosidad de los resultados mediante la
    reproducibilidad, mejorando la calidad y la confianza en la ciencia
    (UNESCO 2020).
  \item
    Fomenta que más investigadores utilicen los datos y produzcan
    información (Whyte y Pryor 2011), aumentando la colaboración y con
    ello la innovación científica, según señala el Foro Abierto de
    Ciencias de Latinoamérica y el Caribe (CILAC) (Ramírez y Samoilovich
    2019).
  \item
    Considerando que los recursos públicos asignados a investigación son
    escasos, la apertura de las bases de datos permite su reutilización
    y por ello ahorra recursos en la construcción de bases de datos,
    fomentando la eficiencia fiscal (Gómez, Méndez Rodríguez, y
    Hernández Pérez 2016).
  \item
    Facilita la preservación de información para estudios históricos que
    recopilen evidencia de larga data. Almacenar los datos permite
    resguardar el conocimiento producido a generaciones futuras.
  \item
    Poder acceder a los productos, procesos y discusiones propias de la
    investigación cualitativa, puede ayudar pedagógicamente a la
    formación universitaria, dando una idea más clara de lo que implica
    una investigación (L. Bishop y Kuula-Luumi 2017).
  \end{itemize}
\item
  Incentivos personales:

  \begin{itemize}
  \item
    Publicar los datos de la investigación fomenta un mayor impacto y
    visibilidad. Como señala la evidencia producida por H. A. Piwowar y
    Vision (2013) una investigación que comparte sus datos puede ser
    citada un 30\% más que una que no lo hace.
  \item
    Compartir el diseño de investigación, las hipótesis y los datos
    puede fomentar comentarios constructivos por parte de otros/as
    investigadores, respecto a la dirección de la investigación,
    ayudando a mejorar la calidad de la publicación final (Sharan 2020)
    y con ello, mejorar también la visibilidad e impacto de la
    investigación.
  \end{itemize}
\end{itemize}

Pese a las ventajas otorgadas por la apertura de la información
producida, los investigadores poseen reticencias a la hora de publicar
abiertamente ``sus'' datos (Ferguson 2014). Estas reticencias de los
investigadores, según Sharan (2020), pueden ser consideradas más bien
mitos, los cuales se sustentan en prenociones que no corresponden a la
realidad. A continuación, se destacan algunos de estos mitos y las
razones de por qué podemos considerarlos como tal.

\textbf{Mito 1: ``Preocupación a las consecuencias negativas'' (mal uso,
consecuencias legales o comerciales).}

En parte este problema se soluciona mediante el uso adecuado de
licencias, las cuales pueden restringir el uso para ciertos fines.
Además, para disminuir el posible mal uso, la preparación de datos
incluye hacer anónimos los datos. Respecto a las consecuencias legales,
hay que considerar que en general estas investigaciones son financiadas
con fondos públicos o por instituciones humanitarias, los cuales en el
contexto actual fomentan en general la apertura de los datos.

\textbf{Mito 2: Temor a la falta del reconocimiento debido de su
trabajo.}

Como señalamos anteriormente, lejos de quitarle merito a su trabajo y
disminuir su reconocimiento, el preparar y publicar los datos de modo
adecuado puede ayudar a difundir la investigación y aumentar el número
de citas. Además, usted también puede nutrirse de las investigaciones
sobre sus datos.

\textbf{Mito 3: Disgusto frente a la carga de trabajo que implica
preparar los datos para su publicación.}

Sin duda esta es una aprensión comprensible, no obstante, cada vez
existen más herramientas que facilitan la labor de la preparación de
datos cualitativos y cuantitativos. Por ejemplo, este documento le
ayudará a disminuir dicha carga. Además, hay que considerar que si bien
aumenta la carga de trabajo actual, disminuye la futura en tanto la
buena documentación de los datos también le permitirán a usted volver a
trabajar con ellos en el futuro o con datos de otra investigación.

\textbf{Mito 4: Desconocimiento de cómo y dónde compartir los datos.}

Esta razón, igual de comprensible, es parte de una cultura académica
poco acostumbrada al trabajo colaborativo. Por ello, debemos aprender
sobre plataformas que faciliten compartir nuestros datos de
investigación. Para ello, este trabajo ofrece una guía para preparar los
datos y sobre como subirlos a la plataforma de ciencia abierta Open
Science Framework.

En miras de las dudas más comprensibles de los investigadores señaladas
(Ferguson 2014; Sharan 2020) y de la necesidad tanto ética como legal de
avanzar hacia la apertura de datos en el contexto chileno, el presente
documento busca facilitar la introducción al almacenamiento y
publicación abierta de datos. Se busca entregar información para
resolver los cuatro puntos señalados, dando cuenta de lo que debe
hacerse para disminuir los riesgos, fomentar el reconocimiento del
trabajo propio, facilitar el mejoramiento de la calidad de los datos y
entregar información sobre dónde y cómo compartir los datos.

Para ello, este documento ofrece una propuesta de pasos a seguir para
mejorar la calidad de los materiales de investigación producidos antes
de publicarla en la web. Esta propuesta busca conciliar los estándares
considerados óptimos en materia de almacenamiento y preservación de
datos con la realidad de las capacidades y herramientas de los
investigadores chilenos en ciencias sociales. Dicho de otro modo,
proponemos una lista de tareas a realizar para cumplir con los
\emph{mínimos} necesarios para mejorar la calidad de los datos a
publicar.

\section{Consejos para la Apertura de Datos de Investigación
Social}\label{consejos-para-la-apertura-de-datos-de-investigaciuxf3n-social}

A continuación, se presenta una sencilla pauta sobre cómo y dónde
publicar abiertamente información cualitativa o cuantitativa producida
por investigaciones sociales, cumpliendo con estándares internacionales
de almacenamiento. Esta pauta ayuda a cumplir con objetivos señalados
por el Consejo Internacional para la Ciencia (ICSU 2014), los cuales
promueven el acceso oportuno a los registros científicos sin barreras,
mejorando su calidad y de modo que sea perdurable en el tiempo.
Igualmente, proponemos una guía de almacenamiento que permite cumplir
con los principios FAIR (Localizable, Accesible, Interoperable y
Reutilizable, en español) promovidos por organizaciones científicas del
Estado Chileno, americanas y europeas (ANID 2020a; Ramírez y Samoilovich
2019; EC 2016). Los estándares FAIR tienen el objetivo de hacer los
datos fáciles de encontrar en la web y que estén en formatos que
cualquier investigador pueda utilizar (FAIR 2020).

Esta pauta para la publicación de datos fue creada en base a otros
materiales de apoyo con objetivos similares, como la guía de preparación
de datos creada por el Consorcio interuniversitario para la
investigación política y social (ICPSR), el manual de Autoevaluación
para proveedores de repositorios abiertos (RISE) del Digital Curation
Center (DDC) (DDC 2017), los planes europeos de manejo de datos (DMP) y
los consejos de distintos investigadores tanto para datos cuantitativos
como cualitativos (CCSDS 2012;
\textbf{kapiszewski\_transparency\_2019?}).

Para publicar los datos de investigación proponemos los siguientes tres
momentos:

\subsection{Preparar}\label{preparar}

\subsubsection{Datos cualitativos:}\label{datos-cualitativos}

Considerando que al estar presente en la entrevista se cuenta con una
mayor información que solo leyendo su transcripción, cualquier anotación
que dé cuenta del ambiente anímico de la entrevista o del gesto que
acompaña alguna frase es bienvenida. También es importante subir la
pauta de la entrevista si es que existe. Asimismo, aquellos
investigadores que hacen un análisis mediante codificación tienen
registro del proceso por el cual llegaron a los códigos utilizados para
la categorización de la información. Además, esto puede ser
complementado con la descripción de las discusiones que surgieron entre
el equipo de investigación para establecer tales códigos y el esquema de
análisis (\textbf{kapiszewski\_transparency\_2019?}).

Igualmente necesario es cambiar en el texto y/o audio lo que sea
necesario para que los sujetos de investigación no puedan ser
identificados. También es conveniente compartir el cuaderno de campo de
la investigación.

Esta información bien almacenada no solo ayudará a las ciencias sociales
por su apoyo con información a otros investigadores, sino que
conjuntamente es un gran aporte a la formación de los estudiantes sobre
cómo investigar cualitativamente, pues estos materiales permiten una
aproximación más concreta al proceso de investigación cualitativo (L.
Bishop y Kuula-Luumi 2017).

\textbf{Confidencialidad en datos cualitativos}

Respecto a la confidencialidad de los datos, ICPSR recomienda que: antes
de enviar datos cualitativos a un archivo, los depositantes de datos
deben tener cuidado de eliminar la información que permita identificar a
cualquiera de sus sujetos de investigación. Este proceso se puede hacer
menos arduo creando un esquema para anonimizar antes de la recopilación
de datos y anonimizando los datos a medida que se crean los archivos
cualitativos para el análisis.

Algunas de las modificaciones que se pueden hacer a los datos
cualitativos para asegurar confidencialidad del entrevistado, según
(Marz y Dunn 2000), son: (1) Reemplazar nombres reales con texto
generalizado; (2) Reemplazar fechas; (3) Eliminar elementos únicos y/o
publicitados.

Dado que los investigadores están más familiarizados con sus datos, se
les pide que utilicen su juicio sobre si cierta información cualitativa
en combinación con el resto del texto o información cuantitativa
relacionada podría permitir la identificación de un individuo.

Los depositantes de datos deben documentar cualquier modificación para
enmascarar información confidencial en los datos cualitativos. Esto
garantizará que el personal del archivo no realice cambios innecesarios
en las modificaciones del investigador cuando realice su revisión de
confidencialidad. Por tanto, la información también se pondrá a
disposición de los usuarios secundarios de los datos para ayudarles a
utilizarlos.

\subsubsection{Datos cuantitativos:}\label{datos-cuantitativos}

\textbf{Bases de datos}

Para resguardar la calidad de nuestros datos debemos asegurar una
correcta \textbf{codificación de las respuestas y nombramiento de las
variables}, ambas bases fundamentales de los datos (filas y columnas).
Para esto, ICPSR propone, entre otros, los siguientes puntos:

\textbf{Errores de codificación}

Verifique cuidadosamente la coherencia entre las respuestas del
cuestionario y los valores en la base de datos para el primer 5 a 10 por
ciento de los registros de datos creados y luego elija registros
aleatorios para controles de calidad. Posteriormente, puede realizar
análisis descriptivos de distribución para evaluar si existen valores
atípicos atribuibles a errores de codificación (por ejemplo 66 en la
variable hijos en vez de 6). El uso de computadores y programas de
encuesta y codificación puede ayudar a disminuir estos errores. Algunos
de los servicios recomendados son Qualtrics o SurveyMonkey, entre otros.

\textbf{Recodificación automática}

Al recodificar la fuente original de datos (sin procesar), es
recomendable dejar que la computadora realice codificaciones y
rectificaciones complejas. Por ejemplo, para crear una serie de
variables que describen la estructura familiar, escriba un código de
computadora para realizar la tarea.

Los códigos de computadora no solo son precisos si las instrucciones son
precisas, sino que también se puede cambiar fácilmente para corregir un
error lógico o de programación. Incluya en la documentación los códigos
utilizados para la recodificación.

\textbf{Consistencia}

Evalúe la coherencia entre las variables, identificando a quienes poseen
combinaciones incoherentes. Por ejemplo, si alguien señala que su hijo
no asiste a la escuela y luego responde preguntas sobre la escuela. Esto
también es posible de realizar a través de filtros en un cuestionario
bien diseñado.

\textbf{Identificadores individuales y grupales}

Proporcione variables identificadoras suficientes. Es fundamental que
cada sujeto posea un identificador único. Además, si la encuesta es
longitudinal, se puede proporcionar junto al identificador de
encuestado, un identificador por cada ocasión que contesta la encuesta.
Otros identificadores dependen del tema del estudio, por ejemplo, si se
trabaja con escuelas, verifique que cada escuela tiene un identificador
id-escuela. Si trabaja con encuestados de modo tal que dos o más son de
la misma familia y cada encuestado corresponde a un núcleo familiar,
indique un único identificador para cada familia.

\textbf{Nombres de Variables}

El nombre de la variable será con lo que más se trabajará en los datos,
por ende, deben ser claros y utilizables por distintos software.

Existen distintos estándares para elegir los nombres de las variables.

\begin{enumerate}
\def\labelenumi{\arabic{enumi}.}
\item
  El primero consiste en asignar un numero único anteponiendo una V de
  modo tal que, siendo n el número de variables, las variables se
  nombran como \emph{Vn} según su posición (p ej. V0001, v0002,\ldots{}
  Vn). Se antepone la V debido a que los software en general no permiten
  nombres de variables con solo caracteres numéricos.
\item
  El segundo modo utiliza letras y números para agrupar las variables
  según escalas o temas (por ejemplo Q1, Q2a, Q2b). Si bien es un
  sistema que entrega más información, no informa sobre el contenido.
\item
  El tercero consiste en utilizar abreviaturas mnemotécnicas, es decir,
  nombres cortos de variables que representan el significado sustantivo
  de las variables facilitando su memorización y comprensión. Por
  ejemplo \emph{educpadr} como ``Educación del Padre''. Este tipo de
  nombres podrían ayudar a disminuir los errores en los análisis
  producidos por agregar una variable incorrecta en el código. El
  problema es que con la limitación de caracteres de los software es
  difícil generar abreviaturas arbitrarias que sean ampliamente
  reconocibles por un público diverso.
\item
  El cuarto consiste en abreviaciones compartidas y registradas. Un
  sistema de raíces y sufijos. Por ejemplo, todas las variables que
  tienen que ver con la educación pueden tener la raíz ED, y podría
  expresarse ``Educación del Padre'' como FAED, siendo esta nomenclatura
  previamente documentada. Esto implica una planificación previa y
  capacidad de organización para compartir las abreviaturas, así como
  herramientas para facilitar el encontrar las abreviaturas correctas en
  la biblioteca o documento de sufijos y prefijos.
\end{enumerate}

En consideración de estas opciones expuestas por ICPSR, se recomienda
utilizar la tercera, puesto que cumple con la cualidad de la primera y
la segunda de identificar las variables de modo único, a la vez que
cumple con el criterio de hacer más comprensible y fácil de recordar.

Junto a lo señalado por ICPSR, consideramos que al crear un nombre de la
variable este debe ser utilizable por los distintos software comúnmente
utilizados como SPSS, STATA y R. En vista de lo anterior sugerimos:

\begin{itemize}
\item
  Dos variables no pueden tener el mismo nombre
\item
  No utilizar más de 12 caracteres en el nombre
\item
  Empezar con una letra
\item
  Deben ser solo alfanuméricos (Números y letras, sin símbolos . ; , :
  `` \$ @)
\item
  En minúscula
\item
  No utilizar la letra ñ, remplazarlo por gn (agnos, en vez de años)
\item
  No utilizar palabras con tildes (tipo á; ü; â)
\item
  Remplazando espacios por guión bajo. (edad\_rec)
\end{itemize}

\textbf{Etiquetas de variables}

Las variables deben ser correctamente etiquetadas. Las etiquetas deben
partir con el número del ítem en el cuestionario para poder asociarlo.
Luego debe darse información sobre el contenido de la variable o
ingresar directamente la pregunta realizada al encuestado.

Considerando las limitaciones de caracteres de los softwares, en base a
manuales universitarios de SPSS y STATA, se sugiere que las etiquetas de
las variables no superen los 120 caracteres.

\textbf{Codificación}

\begin{itemize}
\item
  \textbf{Variables de identificación}. Proporcione campos al comienzo
  de cada registro para acomodar todas las variables de identificación.
  Las variables de identificación a menudo incluyen un número de estudio
  único y un número de encuestado para representar cada caso.
\item
  \textbf{Categorías de código}. Las categorías de códigos deben ser
  mutuamente excluyentes, exhaustivas y estar definidas con precisión.
  Cada respuesta de la entrevista debe encajar en una y solo una
  categoría. La ambigüedad provocará dificultades de codificación y
  problemas con la interpretación de los datos.
\item
  \textbf{Conservación de la información original}. Codifique tantos
  detalles como sea posible. Registrar datos originales, como edad e
  ingresos, es más útil que colapsar o poner entre corchetes la
  información. Con datos originales o detallados, los analistas
  secundarios pueden determinar otros paréntesis significativos por sí
  mismos en lugar de limitarse a los elegidos por otros.
\item
  \textbf{Preguntas cerradas}. Las respuestas a las preguntas de la
  encuesta que están precodificadas en el cuestionario deben conservar
  este esquema de codificación en los datos legibles por máquina para
  evitar errores y confusiones.
\item
  \textbf{Preguntas de final abierto}. Para los ítems abiertos, los
  investigadores pueden usar un esquema de codificación predeterminado o
  revisar las respuestas iniciales de la encuesta para construir un
  esquema de codificación basado en las categorías principales que
  surgen. Cualquier esquema de codificación y su derivación deben
  informarse en la documentación del estudio.
\item
  \textbf{Respuestas codificadas por el usuario}. Cada vez más, los
  investigadores envían el texto completo de las respuestas a las
  preguntas abiertas a los archivos para que los usuarios puedan
  codificar estas respuestas ellos mismos. Debido a que dichas
  respuestas pueden contener información confidencial, deben ser
  revisadas por riesgo de divulgación y, si es necesario, tratadas por
  archivos antes de su publicación.
\item
  \textbf{Comprobar codificación}. Es una buena idea verificar el código
  de algunos casos durante el proceso de codificación, es decir, repetir
  el proceso con un codificador independiente. Por ejemplo, si se asigna
  más de un código a la respuesta de una entrevista, esto resalta
  problemas o ambigüedades en el esquema de codificación. Esta
  codificación de verificación proporciona un medio importante de
  control de calidad en el proceso de codificación.
\item
  \textbf{Serie de respuestas}. Si una serie de respuestas requiere más
  de un campo, organizar las respuestas en clasificaciones importantes
  significativas es útil. Respuestas dentro de cada especialidad o
  categoría se les asigna el mismo primer dígito. Los dígitos
  secundarios pueden distinguir respuestas específicas dentro de las
  categorías principales. Tal esquema de codificación permite el
  análisis de los datos utilizando agrupaciones amplias o categorías más
  detalladas.
\end{itemize}

\textbf{Identificar Casos perdidos}

ICPSR no establece un modo determinado de identificar los casos
perdidos, aunque señala las ventajas y desventajas de distintos tipos de
codificación. Igualmente sugiere distintos tipos de casos perdidos que
deben ser identificados. Cabe destacar que, como regla general para la
preservación, los casos perdidos se deben codificar del modo más similar
a las categorías de las variables, de modo tal que una variable numérica
de un digito se indica con (8,9) y una variable categórica con
alternativas de texto con (``No sabe'', ``No responde'')

\begin{itemize}
\item
  Rechazo / Sin respuesta. El sujeto se negó explícitamente a responder
  una pregunta o no la respondió cuando debería haberlo hecho.
\item
  No lo sé. El sujeto no pudo responder una pregunta, ya sea porque no
  tenía una opinión o porque la información requerida no estaba
  disponible (por ejemplo, un encuestado no pudo proporcionar los
  ingresos familiares en dólares del año anterior).
\item
  Error de proceso. Por alguna razón, no hay respuesta a la pregunta,
  aunque el sujeto proporcionó una. Esto puede resultar de un error del
  entrevistador, codificación incorrecta, falla de la máquina u otros
  problemas.
\item
  No aplica. Al sujeto nunca se le hizo una pregunta por alguna razón. A
  veces, esto se debe a patrones de omisión después de preguntas de
  filtro, por ejemplo, a los sujetos que no están trabajando no se les
  pregunta sobre las características del trabajo. Otros ejemplos de
  inaplicabilidad son los conjuntos de elementos solicitados solo de
  submuestras aleatorias y los solicitados a un miembro de un hogar,
  pero no a otro.
\item
  Sin coincidencia. Esta situación surge cuando los datos se obtienen de
  diferentes fuentes (por ejemplo, un cuestionario de encuesta y una
  base de datos administrativa) y no se puede localizar la información
  de una fuente.
\item
  Datos no disponibles. La pregunta debería haberse formulado al
  encuestado, pero por otro motivo distinto de los enumerados
  anteriormente, no se dio ni registró ninguna respuesta.
\end{itemize}

Considerando las ventajas y desventajas de las distintas formas de
codificación se sugiere utilizar valores perdidos con valores altos en
negativo de modo tal que sean estándar para todas las variables y no
sean confundible con los valores posibles de otras variables. Se
recomienda utilizar la recodificación presente en la
Tabla~\ref{tbl-tabna}, usando valores numéricos o caracteres según
corresponda.

\begin{longtable}[]{@{}ll@{}}
\caption{Recomendación de recodificación para valores
perdidos.}\label{tbl-tabna}\tabularnewline
\toprule\noalign{}
Código de texto & Código numérico \\
\midrule\noalign{}
\endfirsthead
\toprule\noalign{}
Código de texto & Código numérico \\
\midrule\noalign{}
\endhead
\bottomrule\noalign{}
\endlastfoot
No responde & -999 \\
No sabe & -998 \\
Error de Proceso & -997 \\
No aplica & -996 \\
Sin coincidencia & -995 \\
No disponible & -994 \\
\end{longtable}

Para obtener información adicional sobre datos georreferenciados e
imputaciones revise directamente la guía ofrecida por ICPSR disponible
en este
\href{https://www.icpsr.umich.edu/files/deposit/dataprep.pdf}{enlace}

\subsection{\texorpdfstring{\textbf{Documentar}}{Documentar}}\label{documentar}

Existen distintos estándares sobre qué información incorporar junto con
los datos, destacando la importancia de que estos materiales sean
legibles por humanos y por inteligencia artificial (FAIR 2020). A
continuación, se propone un conjunto de documentos comunes que deben
estar incluidos, junto con una guía de como escribirlos en formato
abiertos para cumplir con las recomendaciones del libro \emph{Managing
and sharing research data: a guide to good practice} (Corti 2019).

\textbf{Readme}

El readme es un documento que debe responder las siguientes preguntas
sobre la producción de la información (Tierney y Ram 2020).

\begin{itemize}
\item
  Quién produjo la información
\item
  Cuál es el contenido de los datos
\item
  Cuándo fue producida la información
\item
  Dónde fue recolectada
\item
  Por qué se recopiló
\item
  Cómo se recopiló
\end{itemize}

Para crear este documento se suele utilizar el formato y lenguaje
Markdown de extensión .md. Este formato posee la ventaja de no ser un
formato propietario que permite incluir de modo sencillo distintos
elementos útiles como enlaces, imágenes y tablas. Le recomendamos
realizar la escritura en este formato. Para ello se puede escribir el
documento en RStudio señalando como extensión .md o para quienes no
manejen el software, pueden crear el documento en la página
\href{https://dillinger.io/}{Dillinger}.

\textbf{Licencia}

La licencia se puede crear fácilmente al publicar los datos en el Open
Science Framework. En la siguiente sección se señala cómo hacerlo.

\textbf{Metadatos}

La creación de los metadatos requiere previamente que se cree el
identificador y la página web en Open Science Framework que almacenara
los datos. Por ello la creación de los metadatos será posterior a la
publicación. Cabe desacatar que el modo de crear metadatos difiere según
si el tipo de material es una transcripción de entrevista en PDF o una
base de datos en formato sav, stata o rda.

\textbf{Documentos especiales según tipo de metodología}

Para publicar datos cuantitativos se recomienda recopilar y producir los
siguientes materiales con el objetivo de que los usuarios de la base de
datos cuenten con información suficiente para utilizarla correctamente.
Los siguientes documentos deben ser subidos en PDF y Block de notas.txt
con codificación utf8.

\begin{itemize}
\tightlist
\item
  Cuestionario
\item
  Consentimiento informado
\item
  Libro de códigos
\item
  Ficha técnica
\item
  Manual de usuario
\item
  Descriptivos (Optativo)
\item
  Publicaciones asociadas a los datos (Optativo)
\item
  Descripción e información detallada para metadatos.
\end{itemize}

Para publicar datos cualitativos se sugiere documentar la siguiente
información.

\begin{itemize}
\item
  Transcripciones
\item
  registros audiovisuales (de ser necesario)
\item
  Investigar métodos y prácticas que estén completamente documentados
\item
  Copia en blanco del formulario de consentimiento informado con el
  número de aprobación del IRB
\item
  Detalles sobre el escenario de las entrevistas
\item
  Detalles sobre la selección de los sujetos de la entrevista
\item
  Instrucciones dadas a los entrevistadores
\item
  Instrumentos de recopilación de datos como cuestionarios de
  entrevistas
\item
  Medidas tomadas para eliminar identificadores directos en los datos
  (por ejemplo, nombre, dirección, etc.)
\item
  Cualquier problema que surgió durante el proceso de selección y/o
  entrevista y cómo se manejaron
\item
  Lista de entrevistas
\end{itemize}

\subsection{\texorpdfstring{\textbf{Publicar}}{Publicar}}\label{publicar}

\textbf{OSF y plataformas para publicar datos}

La plataforma web Open Science Framework (OSF) ofrece gratuitamente
servicios de infraestructura digital que permite un espacio de registro
para las distintas etapas de un proyecto de investigación. Actualmente
existen otras plataformas con objetivos similares como Zenodo,
Dataverse, GitHub, Mendeley, Figshare, Dryad o ICPSR, y si bien todas
son buenas herramientas, se remienda utilizar OSF por diversos motivos
señalados por Kryvokhyzha (2019) y un documento informativo de la
Librería de Universidades de la Universidad de OKLAHOMA (Universitarias
2020). En primer lugar, a diferencia de Dryad, OSF es una plataforma
gratuita, con mejor estructura de repositorios y con posibilidades de
corregir errores. En segundo lugar, Figshare tiene la capacidad de
estructurar los repositorios en distintos componentes, pero no está
optimizado para descargar muchos archivos a la vez y, además es una
empresa con fines de lucro. Por su parte OSF, permite estructurar de
diversos modos los repositorios, esta optimizado para descargas y es una
organización sin fines de lucro que es financiada por el Centro para la
Ciencia Abierta \href{https://www.cos.io/}{(cos)} con recursos para 50
años más. En tercer lugar, a diferencia de Zenodo y GitHub, Open Science
Framework si cuenta con estadísticas de descarga que nos permiten
evaluar la visibilidad de los datos. En último lugar, GitHub si bien es
una alternativa que ha sido utilizada para este fin, en realidad más que
una plataforma para subir datos es una plataforma para crear códigos de
forma colaborativa, usualmente de aplicaciones, además no nos permite
crear de modo automático un DOI, lo cual es fundamental para cumplir con
los principios FAIR.

No obstante, si el equipo investigador lo considera conveniente podría
subir sus datos a múltiples plataformas. En esta línea, además de OSF,
recomendamos almacenar los datos paralelamente en
\href{https://www.openicpsr.org/openicpsr/workspace?path=openICPSR}{OpenICPSR}
puesto que esto permitirá conectar nuestros datos con el buscador de
ICPSR, el cual tiene amplia visibilidad dentro del campo de las ciencias
sociales a nivel internacional.

En esta página, un usuario puede crear un repositorio denominado
proyecto, que puede contener, al mismo tiempo, componentes que pueden
ser investigaciones o datos específicos dentro del proyecto. Se pueden
crear más componentes dentro de los ya creados del proyecto. Esta
estructura permite almacenar conjuntamente algunos trabajos
relacionados.

La plataforma promueve que se registren productos de la investigación de
las distintas etapas del proceso. En primer lugar, posee un espacio para
los ``preregistros'' que son un documento en el cual se expone
brevemente el diseño de la investigación, las hipótesis y la
metodología, lo cual aumenta la rigurosidad de las investigaciones.

Ambas evaluaciones de OSF señalan que un problema es el límite de
almacenamiento de solo 5gb por archivo. Este límite ha sido modificado
el 5 de noviembre del 2020, agregando un límite de 5gb a los proyectos
que estén privados y 50gb a los públicos. No obstante, el problema del
espacio se puede resolver conectando OSF con otros servicios de
almacenamiento con mayor capacidad.

\textbf{Crear Metadatos para datos cualitativos o cuantitativos.}

\textbf{Paso 1}

Para crear los metadatos de forma muy sencilla se puede utilizar un
archivo csv, que se puede editar en cualquier Hoja de cálculo. A
continuación, entregamos un enlace para descargar un csv con los campos
para rellenar los metadatos, cumpliendo con los campos utilizados por
Dataverse Harvard y los ``Social Science and Humanities Metadata'' de
ICPSR.

Para crear los metadatos usted solo debe escribir en las casillas de
abajo de las categorías, la respuesta para cada uno de los campos,
señalando el identificador, el Título de los datos, los autores, entre
otros. Es necesario que rellene los campos en inglés. Con una traducción
simple de \href{https://translate.google.com/}{Google Traductor} o
\href{https://www.deepl.com/es/translator}{DeepL} es suficiente.

Para descargar el archivo csv editable seleccione
\href{https://raw.githubusercontent.com/franciscomeneses/CADIS/master/metadata.csv}{Descargar
Metadata.cvs}.

Cabe destacar que, para el primer campo correspondiente al
Identificador, antes del código entregado por OSF debe estar escrito
https://doi.org/ como se muestra en el documento descargado. Esto es
para cumplir con los estándares FAIR.

\textbf{Paso 2}

Después de haber creado el documento csv con nuestra información solo
debemos seleccionar el siguiente enlace e ir al sitio web:
\href{https://csvjson.com/csv2json}{csvjson.com}. En este sitio tenemos
que apretar Select a file\ldots{} y buscar el documento csv creado.
Posteriormente, se debe señalar output: ``Array'' y ``Minify'', como se
señala en la Figura~\ref{fig-csvjson}. Con estas opciones apretamos el
botón morado \textgreater Convert bajo el cuadrado que posee nuestro
documento CSV. Cuando esté listo apretamos el botón Download, como se
señala en la Figura~\ref{fig-csvjson}. El documento se descargará con el
nombre csvjson, debe cambiarlo a metadata.

\begin{figure}

\centering{

\includegraphics[width=0.8\linewidth,height=\textheight,keepaspectratio]{images/csvtosjon.png}

}

\caption{\label{fig-csvjson}Creación de documento de metadatos}

\end{figure}%

\textbf{Paso 3}

Finalmente debemos agregar los archivos csv y json al repositorio,
recuerde cambiar el nombre del documento a metadata.

\bookmarksetup{startatroot}

\chapter{Análisis reproducible}\label{anuxe1lisis-reproducible}

En la discusión sobre los problemas de transparencia en torno a los
procedimientos de investigación, se vuelve necesario precisar de qué
manera es entendido el concepto de reproducibilidad en la ciencia. En
esta línea, la laxitud con que se ha empleado el término ha llevado a
definiciones poco claras, lo cual ha generado una tendencia a confundir
lo que refiere a la transparencia de un proceso único que ya ha sido
realizado, con un proceso nuevo y que puede realizarse de manera
reiterativa, obteniendo los mismos resultados. Por este motivo, esta
sección propone dar luces respecto a cómo entendemos el concepto de
reproducibilidad, en contraste con el de replicabilidad en la ciencias
sociales.

La discusión en torno a cómo se entiende la \textbf{reproducibilidad},
habitualmente lleva al contraste respecto al concepto de
\textbf{replicabilidad}. Al respecto Earth y Behavioral (2019) mencionan
que, por un lado, con el incremento de las herramientas computacionales
a principios de los años 90', el término de ``investigación
reproducible'' era concebido como las investigaciones que proveían un
compendio detallado de la documentación, código y datos que permitieran
obtener los mismos resultados publicados por los autores, enfatizando
que los análisis fueran transparentes y claros con el objetivo de ser
verificados por sus pares. Por otro lado, los autores sostienen que, en
otras disciplinas, el concepto de reproducibilidad era asociado a
investigaciones independientes entre sí en términos de los datos
empleados, los materiales, métodos e implementación de un estudio, lo
cual estaría orientado a robustecer o cuestionar la evidencia previa
(Earth y Behavioral 2019, pp 33-34). Actualmente, a esta última práctica
se la entiende como replicabilidad de una investigación y no debe ser
confundida con el concepto de reproducibilidad (Barba 2018).

Barba (2018) sugiere que la confusión entre reproducibilidad y
replicabilidad ha contribuido a obstaculizar las prácticas en ambas
dimensiones. En una revisión reciente realizada por la autora se han
identificado al menos tres escenarios o versiones de cómo se entienden
ambos conceptos en una amplia gama de disciplinas que van desde las
ciencias sociales hasta estudios clínicos en las ciencias médicas. El
primer escenario (A), y a la vez el más común, es donde el uso de ambos
conceptos es indistinto, contribuyendo a la ya mencionada confusión. El
segundo escenario (B1) es cuando la reproducibilidad es entendida como
la situación en que los datos originales y el código de análisis son
empleados para \textbf{regenerar} los resultados originales, mientras
que la replicabilidad es entendida cuando investigadores o equipos
independientes utilizan datos nuevos para obtener los mismos resultados
que la investigación previa. Finalmente, un tercer escenario (B2) es
cuando la reproducibilidad es entendida cuando investigadores o equipos
independientes obtienen los mismos resultados empleando sus propios
datos y métodos, mientras que la replicabilidad es entendida cuando
investigadores o equipos independientes llegan a los mismos resultados
empleando los artefactos digitales \footnote{Barba (2018) lo define como
  un compendio que detallar la estrategia de medición, diseño del
  estudio o código de análisis originales de un autor} originales del
autor con menores o mayores modificaciones, los cuales han sido puestos
previamente a disposición de sus pares. La Figura N° @ref(fig:scenarios)
ilustra cómo podemos entender los escenarios B1 y B2 en relación a la
distinción entre reproducibilidad y replicabilidad. El color rojo, tanto
en los datos como en los métodos, indica que los componentes empleados
son idénticos a los del estudio original. Por otro lado, el color azul
indica que tanto los datos como los métodos son distintos a los del
estudio original. Finalmente, el color morado en los métodos se entiende
como un punto intermedio y refiere a cuando se han empleado métodos que
siguen las indicaciones del estudio original, pero que han incorporado
modificaciones, nuevos métodos u otras innovaciones metodológicas
(p.~ej. métodos nuevos, pruebas de robustez u otros).

\begin{figure}

\centering{

\includegraphics[width=0.75\linewidth,height=\textheight,keepaspectratio]{images/reproducibility.png}

}

\caption{\label{fig-scenarios}Escenarios B1 y B2 en reproducibilidad y
replicabilidad.}

\end{figure}%

En las ciencias sociales, el debate en torno a la investigación
reproducible y la replicabilidad no ha estado ausente. Como fue reseñado
en el capítulo de transparencia, existen casos icónicos en torno a
prácticas cuestionables de investigación que han afectado la confianza
en la investigación científica, lo cual ha contribuido a incrementar los
esfuerzos por una ciencia social abierta y reproducible (Breznau 2021;
B. A. Nosek et~al. 2015). En los tres escenarios descritos por Barba
(2018), las ciencias sociales han experimentado de manera diversa el
ajuste hacia una cultura de mayor apertura y precisión en torno a los
problemas de la crisis de reproducibilidad, principalmente a través del
estudio sistemático de dicha problemática, dentro de lo cual la
psicología ha sido pionera en proveer evidencia para este debate (e.g.
Open Science Collaboration 2015; Gilbert et~al. 2016). Al respecto D.
Bishop (2019) sostiene que una de las principales amenazas para el
progreso de la ciencia en general ha sido la falta de reproducibilidad
de los resultados (\emph{irreproducibility}), lo cual ha afectado
principalmente la robustez y credibilidad de la evidencia reportada por
las investigaciones, problema que también ha sido identificado en las
ciencias sociales, principalmente por la falta de precisión en los
procedimientos y las barreras de acceso a materiales clave del proceso
de análisis (Freese y Peterson 2017).

Entonces, retomando la distinción clave entre lo que entendemos por
\textbf{reproducibilidad} y \textbf{replicabilidad}, en su revisión,
Barba (2018) sugiere que una manera de entender y distinguir ambos
conceptos de manera minimalista puede descansar en el carácter de los
\emph{datos} y los \emph{métodos}. Al respecto B. A. Nosek et~al. (2015)
sostiene que en lo que refiere a estas dos dimensiones, los niveles en
que una publicación los incorpora es gradual y puede entenderse como un
continuo o espectro (Peng 2011) y, por tanto, el nivel en que se cumplen
bajo determinados criterios nos permite definir el carácter de una
investigación en términos de su reproducibilidad. Por ejemplo, la Figura
N° @ref(fig:espectro) nos muestra cómo podemos caracterizar una
investigación publicada en torno al acceso y vinculación entre código y
datos. Por un lado, se observa que en el polo donde únicamente
disponemos de la publicación, se entiende como la ausencia de
reproducibilidad. Por otro lado, en la medida que incrementa el acceso a
los materiales, y se explicita el enlace entre ellos, se puede
caracterizar a una publicación como reproducible. \footnote{En la figura
  original, Peng (2011) muestra el polo derecho como el mejor escenario
  y lo clasifica como \emph{Full replication}, sugiriendo que el mejor
  estándar para poner a prueba los hallazgos de una investigación
  científica es la replicación, pero en la ausencia de dicha posibilidad
  la reproducibilidad de los resultados debiese ser un estándar mínimo}

\begin{figure}

\centering{

\includegraphics[width=0.75\linewidth,height=\textheight,keepaspectratio]{images/repro-spectrum.png}

}

\caption{\label{fig-espectro}Espectro de Reproducibilidad. Traducción
propia en base a (\textbf{peng\_Reproducible\_2011?})}

\end{figure}%

Como sugieren B. A. Nosek et~al. (2015), el problema de la ausencia o
falta de reproducibilidad debe ser abordado a través de un cambio en las
prácticas de investigación, para lo cual se requiere de una disposición
por parte de la comunidad científica, es decir, que se le atribuya un
\emph{sentido} positivo a estas prácticas. Sin embargo, Peng (2011)
sostiene que una de las principales barreras para promover estas
prácticas ha sido la falta de mecanismos que faciliten la distribución
de la investigación reproducible, como también la poca claridad respecto
de los estándares asociados a ello. Siguiendo esta autocrítica de
algunos sectores dentro de la comunidad científica, dentro de los
últimos años han surgido iniciativas como, por ejemplo, el Open Science
Framework, al alero del \href{https://www.cos.io/}{Center for Open
Science}, desde donde se busca contribuir con herramientas para el
entrenamiento y educación de la comunidad científica en general, como
también proveer de una infraestructura tecnológica que facilite la
transición cultural hacia una ciencia abierta, transparente y
reproducible (B. A. Nosek et~al. 2015). Por este motivo, proponemos
revisar tres iniciativas internacionales que han puesto sus esfuerzos en
la promoción de estos principios, con particular atención en la
reproducibilidad de la investigación científica, y en particular de las
ciencias sociales empíricas cuantitativas. Dentro de estas iniciativas
encontraremos esfuerzos orientados a la educación y entrenamiento,
herramientas tecnológicas y fortalecimiento de redes de colaboración.

\section{¿Qué se ha hecho?}\label{quuxe9-se-ha-hecho}

\subsection{Berkeley Initiative for Transparency in the Social Sciences
(BITSS)}\label{berkeley-initiative-for-transparency-in-the-social-sciences-bitss}

\begin{center}

\includegraphics[width=0.75\linewidth,height=\textheight,keepaspectratio]{images/BITSS_logo_horizontal.png}
\end{center}

\textbf{Objetivos y visión}

Esta iniciativa busca promover la credibilidad en la evidencia generada
por las ciencias sociales a través de mecanismos de avanzada para la
transparencia, reproducibilidad y prácticas éticas en la investigación
social empírica. Desde esta premisa, se ha desarrollado y puesto a
disposición de la comunidad científica una serie de herramientas en
colaboración con estudiantes, investigadores, entidades académicas y
fundaciones de la sociedad civil al alero de tres principios
orientadores:

\begin{enumerate}
\def\labelenumi{\arabic{enumi}.}
\item
  Generar evidencia en torno a problemas y soluciones a través de los
  investigadores y la comunidad de BITSS, quienes han liderado
  investigaciones meta-analíticas, con atención en las ciencias
  sociales.
\item
  Incrementar el acceso a la enseñanza de la ciencia abierta, a través
  del fortalecimiento de prácticas para reconocer y conducir
  investigación social transparente y reproducible a través del
  entrenamiento de investigadores jóvenes, acceso a materiales, apoyo
  financiero y la consolidación de una red de colaboración.
\item
  Fortalecer el ecosistema científico, estableciendo condiciones para
  investigadores e instituciones para contribuir a un cambio efectivo y
  equitativo en las normas que permitan una consolidación de una
  política interna orientada a la ciencia abierta y al desarrollo de
  protocolos en esta dirección.
\end{enumerate}

Desde sus inicios, se han desarrollado una serie de componentes que
buscan promover y dar soluciones a los problemas de transparencia y
reproducibilidad en las ciencias sociales. En particular, nos interesa
destacar algunas de las contribuciones en este ámbito que serán
presentadas a continuación, las cuales se pueden resumir en Evidencia,
Educación y Recursos.

\textbf{Contribución}

En el ámbito de Evidencia, desde BITSS se ha realizado un esfuerzo por
producir y sistematizar evidencia centralizadamente. En este contexto
existe la \href{https://www.bitss.org/research-library/}{Research
Library}, una base de datos de publicaciones científicas que engloba una
serie de investigaciones meta-analíticas en el ámbito de las ciencias
sociales, contribuyendo con un cuerpo de evidencia sistemática en torno
a los problemas y soluciones sobre transparencia y reproducibilidad en
las ciencias sociales sin precedentes. En este apartado, tanto los
colaboradores como investigadores de BITSS ponen a disposición de la
comunidad científica las investigaciones que han sido financiadas a
través de las Social Science Meta-Analysis and Research Transparency
(\href{https://www.bitss.org/ssmart-grants/}{SSMART}) grants, las cuales
constituyen fondos orientados a contribuir a la investigación empírica
en torno a la transparencia y reproducibilidad en disciplinas como la
economía, ciencia política, psicología y ciencias sociales afines.

Desde la Educación y Entrenamiento podemos identificar la articulación
de una serie de \emph{Training activities} desarrolladas por BITSS.
Dentro de los objetivos de estas actividades encontramos dos aspectos
que se buscan abordar desde esta dimensión. Por un lado, se encuentra el
promover una visión crítica de los investigadores en las ciencias
sociales, esto considera un entendimiento de los principales problemas
asociados a la investigación social de calidad al alero de los
principios de la ciencia abierta, dentro de lo cual podemos encontrar
los sesgos y prácticas referidas a las presiones por publicar, prácticas
cuestionables de investigación, reproducibilidad y privacidad de datos.
Por otro lado, se han propuesto promover el manejo de técnicas de
investigación para la transparencia y reproducibilidad, principalmente a
través de actividades de entrenamiento con un foco en el aprendizaje e
implementación de herramientas y métodos. En esta línea destacan dos
contribuciones que se fundamentan en estos principios, las cuales serán
descritas a continuación.

\textbf{Research Transparency and Reproducibility Training}

Una de las contribuciones señaladas es el Research Transparency and
Reproducibility Training (RT2), el cual constituye uno de los
principales eventos académicos realizados anualmente por BITSS, teniendo
por objetivo el poner a disposición de estudiantes e investigadores una
mirada general de las herramientas y prácticas actuales para la
transparencia y la reproducibilidad en la investigación empírica en
ciencias sociales. Los contenidos de RT2 abordan una serie de tópicos de
manera transversal que pueden ser ilustrados en seis puntos:

\begin{itemize}
\item
  \textbf{Amenazas} para la credibilidad en la ciencia y la
  reproducibilidad, junto con su relación con el \emph{ethos}
  científico: Conducta y valores en la ciencia.
\item
  \textbf{Mejoras} en las especificaciones de los diseños de
  investigación: pre-registros y plan de pre-analysis en investigación
  con datos experimentales y observacionales.
\item
  \textbf{Ética e investigación abierta}: estándares éticos para la
  ciencia abierta, manejo de datos y autoría de fuentes de información
  abiertas (citación).
\item
  \textbf{Herramientas y métodos} para la investigación reproducible y
  colaboración: control de versiones y reportes dinámicos.
\item
  \textbf{Sistematización de evidencia}, reproducibilidad e
  interpretación: métodos para investigación meta-analítica y revisiones
  sistemáticas, transparencia y reproducibilidad usando datos
  administrativos; y replicabilidad en la investigación.
\item
  \textbf{Software} para la Ciencia Abierta e innovaciones
  metodológicas.
\end{itemize}

\textbf{MOOC: Transparent and Open Social Science Research}

Otra de las contribuciones es el Transparent and Open Social Science
Research que corresponde a un curso gratuito online de cinco semanas el
cual aborda los fundamentos conceptuales y las principales herramientas
para promover una ciencia social abierta y transparente. La Tabla
@ref(tab:mooc) muestra el contenido de las sesiones, las cuales se basan
en un curso de nivel de grado dictado por el director de BITSS Ted
Miguel en la Universidad de California Berkeley.

\begin{longtable}[]{@{}
  >{\raggedright\arraybackslash}p{(\linewidth - 2\tabcolsep) * \real{0.0530}}
  >{\raggedright\arraybackslash}p{(\linewidth - 2\tabcolsep) * \real{0.9470}}@{}}
\caption{(\#tab:mooc) Cursos por semana en el MOOC de
BITSS}\tabularnewline
\toprule\noalign{}
\begin{minipage}[b]{\linewidth}\raggedright
Semana
\end{minipage} & \begin{minipage}[b]{\linewidth}\raggedright
Contenido
\end{minipage} \\
\midrule\noalign{}
\endfirsthead
\toprule\noalign{}
\begin{minipage}[b]{\linewidth}\raggedright
Semana
\end{minipage} & \begin{minipage}[b]{\linewidth}\raggedright
Contenido
\end{minipage} \\
\midrule\noalign{}
\endhead
\bottomrule\noalign{}
\endlastfoot
1 & \href{http://bitss.org/week-1/}{Introducción a la transparencia y
reproducibilidad de la investigación} \\
2 & \href{http://www.bitss.org/week-2-publication-bias/}{Sesgo de
publicación} \\
3 &
\href{http://www.bitss.org/week-3-pre-registration-pre-analysis-plans-and-meta-analysis/}{Pre-registro,
Plan de Pre-Análisis; y Meta-análisis} \\
4 &
\href{http://www.bitss.org/week-4-replication-and-open-data/}{Replicación
y Datos Abiertos} \\
5 &
\href{http://www.bitss.org/week-5-transparent-data-visualization-and-looking-forward}{Visualización
de Datos transparente y Viendo hacia adelante} \\
\end{longtable}

Una de las principales características de este curso introductorio es la
sistematización de aspectos claves para la ciencia abierta con un foco
particular en las ciencias sociales. Adicionalmente, tiene el objetivo
de introducir conceptualmente a los problemas que se han visto presentes
en las ciencias y busca dar respuestas a prácticas a través de
herramientas y métodos concretos para solucionarlo. Finalmente,
constituye un esfuerzo breve y preciso, dado que las sesiones semanales
poseen una duración promedio de unos treinta minutos y se encuentran
dosificadas en videos de corta duración subtitulados.

En el ámbito de los Recursos que provee BITTS, podemos encontrar una
librería de recursos o simplemente la
\href{https://www.bitss.org/resource-library/}{\emph{Resource Library}},
la cual incluye una multiplicidad de recursos de aprendizaje digitales
en torno a la transparencia y reproducibilidad, ordenados según (i)
Tópico, (ii) Tipo y (iii) Disciplina de las ciencias sociales. La Figura
@ref(fig:resources) muestra cómo se visualizan los tópicos disponibles
en la librería, lo cual puede ser ordenado según tipo y disciplina.

\begin{figure}

\centering{

\includegraphics[width=0.75\linewidth,height=\textheight,keepaspectratio]{images/resource-library.PNG}

}

\caption{\label{fig-resources}Librería de Recursos de BITSS}

\end{figure}%

\subsection{Proyecto TIER (Teaching Integrity in Empirical
Research)}\label{proyecto-tier-teaching-integrity-in-empirical-research}

\begin{center}

\includegraphics[width=0.75\linewidth,height=\textheight,keepaspectratio]{images/tier-logo.jpg}
\end{center}

\textbf{Objetivos y visión}

El proyecto TIER es una iniciativa respaldada por la
\href{https://sloan.org/}{Fundación Alfred Sloan} que se propone
contribuir a un cambio en las normas y conductas profesionales en torno
a la transparencia y reproducibilidad en la investigación empírica en
las ciencias sociales.

Uno de los principios orientadores de sus actividades es el proveer
formación en herramientas para la documentación oportuna de
procedimientos que involucren datos estadísticos a través de rutinas y
referencias que garanticen la \textbf{reproducibilidad} de estos. La
idea subyacente que motiva estas acciones es que los autores puedan
concebir la documentación como un componente esencial de la
\textbf{comunicación} de sus resultados con sus pares, como también al
público no especializado, de modo tal que estas medidas contribuyan a
incrementar la confianza y credibilidad en la evidencia científica. En
esta línea, su declaración de principios sostiene que su objetivo se
puede considerar como logrado cuando:

\begin{quote}
(\ldots) no proporcionar documentación de replicación para un estudio
empírico se considere tan \textbf{aberrante} como escribir un artículo
teórico que no contenga pruebas de las proposiciones, un artículo
experimental que no describa las condiciones de tratamiento o un
artículo de revisión de leyes que no cite los estatutos legales o las
decisiones judiciales. (traducción propia)
\end{quote}

\textbf{Contribución}

Es necesario tener presente que uno de los principales campos de acción
del proyecto TIER es la \textbf{Educación} y \textbf{Entrenamiento}
hacia cientistas sociales en formación, tomando en consideración que es
en el ciclo formativo inicial donde se deben impulsar la adopción de
prácticas integrales para la investigación social empírica. En esta
línea, uno de los elementos destacables es la sección de herramientas
para la enseñanza titulada ``TIER in the Classroom'', sus contenidos
referidos a temas de \textbf{reproducibilidad} se pueden resumir de la
siguiente manera:

\begin{itemize}
\item
  \href{https://www.projecttier.org/tier-classroom/soup-nuts-exercises/}{\textbf{\emph{Soup-to-Nuts
  Exercises}}}: No existe una traducción en el español, no obstante la
  expresión ``Soup-to-Nuts'' refiere a un proceso de ``inicio-a-fin''.
  Como lo dice, esta sección muestra ejercicios orientados a la
  reproducibilidad de los análisis pasando por (1) los datos, (2)
  procesamiento, (3) análisis y (4) reporte. La idea fuerza de este
  ejercicio es introducir a estudiantes a los principios y prácticas
  fundamentales de la investigación social transparente y reproducible
  para que los implementen en sus tareas o informes.
\item
  \href{https://www.projecttier.org/tier-classroom/course-materials/}{\textbf{Materiales
  para clases}}: Esta sección está fuertemente orientada al análisis
  estadístico y a los métodos cuantitativos. Se presentan una serie de
  veinticuatro cursos de pregrado y postgrado que incorporan en su
  currículum los principios de transparencia y reproducibilidad en la
  enseñanza de los métodos de manera transversal. Los materiales de cada
  curso se encuentran disponibles para libre descarga, incorporando
  ejercicios de análisis estadístico (R, Stata, SPSS), reportes
  dinámicos (R Markdown, Markstat) y sus respectivos \emph{syllabus}.
\item
  \href{https://www.projecttier.org/tier-classroom/student-work/\#student-papers}{\textbf{Trabajos
  estudiantiles}}: En esta sección se incorporan una serie de trabajos
  estudiantiles/papers, los cuales están acompañados de una completa
  documentación basada en el
  \href{https://www.projecttier.org/tier-protocol/}{Protocolo TIER (ver
  detalle abajo)}. El objetivo es presentar modelos de trabajos
  realizados con análisis reproducibles, de modo tal que quien esté
  interesado en emplear la estructura de un proyecto pueda observar un
  trabajo real e, idealmente, logre reproducir completamente sus
  resultados.
\end{itemize}

Una de las contribuciones más relevantes del proyecto TIER es la
elaboración de \textbf{estándares} para la \textbf{organización},
\textbf{publicación} y \textbf{comunicación} de proyectos de
investigación empírica cuantitativa reproducible. Al respecto, existen
dos esfuerzos orientados a este fin:

Por un lado tenemos el
\href{https://www.projecttier.org/tier-protocol/specifications-3-0/\#overview-of-the-documentation}{Protocolo
TIER}, el cual constituye una serie de especificaciones respecto a los
contenidos de la documentación para la replicación de un estudio, el
cual está orientado a ser empleado para la enseñanza de la investigación
que incorpore la reproducibilidad de los análisis. En este caso es
importante precisar, como ya hemos identificado en un principio, que el
concepto de \textbf{replicación} se emplea como sinónimo de
\textbf{reproducibilidad}, entendiendo este último como la conjunción de
datos y métodos originales que nos permitan \textbf{regenerar} los
resultados de un estudio que ha sido publicado. Por lo tanto, cuando en
TIER se habla de replicación se refiere a esta idea. La documentación
debe incluir una serie de elementos descritos a continuación.

\begin{itemize}
\item
  Datos empleados por el proyecto
\item
  Rutinas de código escrito en el software empleado para la preparación
  y análisis estadístico. Esto se incluye dado que el objetivo es
  proveer los datos brutos a procesar, junto con todas las instrucciones
  que permiten \textbf{regenerar} los resultados reportados en el
  estudio.
\item
  Fuentes de información que contribuyan a comprender detalladamente
  cada sección del estudio de inicio a fin.
\end{itemize}

Por otro lado tenemos el
\href{https://www.projecttier.org/tier-protocol/dress-protocol/}{Protocolo
DRESS} (Documenting Research in the Empirical Social Sciences). Al igual
que el Protocolo TIER, se incorporan los mismos estándares para la
documentación de una investigación transparente que incorpore la
reproducibilidad de los análisis. Sin embargo, este se encuentra
adaptado a los propósitos de los \textbf{investigadores profesionales},
más que para el uso de los estudiantes durante su formación en
investigación.

\subsection{UK Reproducibility Network
(UKRN)}\label{uk-reproducibility-network-ukrn}

\begin{center}

\includegraphics[width=0.75\linewidth,height=\textheight,keepaspectratio]{images/UKRN-Logo.png}
\end{center}

\textbf{Objetivos y visión}

La UK Reproducibility Network (UKRN) es un consorcio institucional del
Reino Unido que tiene por objetivo promover los principios y prácticas
de la ciencia abierta con una mirada local, es decir, en las
instituciones nacionales y sus investigadores. Para contribuir a este
objetivo se realizan esfuerzos en torno a la investigación de los
factores que determinan una investigación abierta y robusta, promoviendo
el entrenamiento a través de actividades abiertas y diseminando las
buenas prácticas para una ciencia abierta. En particular, se proponen
profundizar en los factores que determinan la carencia de
\textbf{reproducibilidad} y \textbf{replicabilidad}, para lo cual se
busca:

\begin{itemize}
\item
  Desarrollar aproximaciones que contrarresten esta falta de
  transparencia.
\item
  Incrementar la confianza y la calidad de la investigación científica.
\item
  Abordar de manera transversal estos problemas en las distintas
  disciplinas científicas.
\item
  Avanzar hacia un cambio cultural en la ciencia y transformar las
  prácticas de quienes la desarrollan.
\end{itemize}

La UKRN se caracteriza por un trabajo en red, es decir, por un
importante componente de vinculación entre instituciones de
investigación asociadas a universidades como también a oficinas
gubernamentales que desarrollan investigación (ver
\href{https://www.ukrn.org/stakeholders/}{External Stakeholders}). En
esta línea, existen diversas iniciativas apoyadas por la UKRN que
promueven el entrenamiento, metodologías y recursos tecnológicos para la
ciencia abierta. A continuación se presentarán algunas de las
contribuciones más relevantes realizadas por la red, como también
algunas de las iniciativas externas que han sido respaldadas por la
UKRN.

\textbf{Contribución}

En el ámbito de la \textbf{Educación} y \textbf{Entrenamiento}, es
posible identificar, por un lado, las contribuciones realizadas
directamente por la UKRN y, por otro lado, las iniciativas que son
respaldadas por la red y que promueven la formación en torno a los
principios y prácticas de la ciencia abierta, particularmente en la
etapa temprana de la carrera de investigación.

Respecto a una de las iniciativas elaboradas por los académicos e
investigadores involucrados en la UKRN, encontramos uno de los
principales recursos virtuales en un breve curso online que aborda una
serie de tópicos relevantes para la promoción de la ciencia abierta,
dentro de lo cual encontramos el uso de pre-prints, autorías, registered
reports, datos abiertos y reproducibilidad. A continuación se puede
observar la lista de sesiones que han sido desarrolladas en torno a
estos temas.

\url{https://www.youtube.com/embed/videoseries?list=PLehgGZrvxReSisRauEJ2DxPAIYVVEr1qC}

Junto con las sesiones, existe una serie de recursos compartidos a
través de un proyecto abierto en el \href{https://osf.io/x8j9q/}{Open
Science Framework}. Aquí es posible acceder a documentos breves que
abordan los tópicos de cada sesión, además de
\href{https://osf.io/qw9ck/}{recursos adicionales} sobre uso de software
de código abierto y repositorios.

Un ámbito de desarrollo ha sido la disposición de \textbf{recursos
tecnológicos} que promuevan y faciliten las prácticas en ciencia
abierta. Una de las iniciativas impulsadas es el
\href{https://openresearchcalendar.org/}{\textbf{Open Research
Calendar}}, el cual consiste en un instrumento colaborativo y abierto
que busca brindar a la comunidad de investigadores interesados en temas
relacionados a la ciencia abierta un flujo constante de actualizaciones
en torno a workshops y conferencias a nivel mundial que abordan tópicos
sobre ciencia abierta unificados en un calendario. El carácter
\textbf{colaborativo} de esta herramienta permite que usuarios
previamente registrados y validados puedan contribuir con información
que se centraliza en el calendario de eventos, precisando los contenidos
y redireccionando a la inscripción y/o enlace para las actividades que
se realizan a través de internet. Para facilitar la experiencia de
usuario, el calendario se integra con Google Calendar que puede
sincronizarse con la agenda personal y que se van actualizando
automáticamente.

Otra herramienta tecnológica patrocinada por la UKRN es la plataforma
\href{https://science-octopus.org/}{Octopus}. A la fecha, la plataforma
se presenta como una aplicación en desarrollo y abierta a comentarios de
los usuarios. En términos generales se propone ser una alternativa para
contribuir a la apertura de publicaciones. El detalle se presenta así:

\begin{quote}
(\ldots) sustituir a las revistas y los artículos como lugar para
establecer la prioridad y registrar su trabajo con todo detalle, Octopus
es de uso gratuito y publica todo tipo de trabajos científicos, ya sea
una hipótesis, un método, datos, un análisis o una revisión por pares
(traducción propia).
\end{quote}

La Figura @ref(fig:octopus) ilustra un ejemplo de cómo se ve un proyecto
en Octopus. Vemos que existen siete componentes que buscan representar
el flujo de una investigación. Entendiendo que los procesos de
investigación no son lineales y tienden a existir iteraciones entre
teoría y métodos, la virtud del registro y publicación de un proyecto
permite que otros puedan conocer y evaluar nuestras hipótesis, plan de
análisis, resultados y versiones de un artículo, así como también la
vinculación entre cada sección.

\begin{figure}

\centering{

\includegraphics[width=0.75\linewidth,height=\textheight,keepaspectratio]{images/octopus.PNG}

}

\caption{\label{fig-octopus}Ejemplo de un trabajo registrado en
desarrollo en octopus.org}

\end{figure}%

Para publicar debemos
\href{https://science-octopus.org/publish}{ingresar} con una cuenta de
ORCID. Si no tienes una cuenta puedes crear un perfil
\href{https://orcid.org/}{aquí}. Luego, se deben seguir tres pasos. El
primero es elegir qué tipo de componente se desea publicar (Problema,
Hipótesis, Métodos, etc). Segundo, dar detalles sobre el componente y
con qué otros proyectos se relaciona. Finalmente, contribuir con un
borrador de escritura que luego será publicado.

\bookmarksetup{startatroot}

\chapter{Publicaciones libres}\label{publicaciones-libres}

Un agricultor hace uso de sus \emph{herramientas, talentos y
conocimientos} para aprovechar lo mejor posible las ventajas del bien
común que es la tierra. Una vez que termina la temporada, el agricultor
entrega la cosecha a un distribuidor que se encarga de comercializar el
fruto de su trabajo y el de otros agricultores. ¿Cuánto recibe realmente
el agricultor por su trabajo? \emph{¿Cuáles son las posibilidades reales
que tiene la ciudadanía de acceder al bien común de manera autónoma?}

El problema del agricultor, el distribuidor y la ciudadanía no es una
historia del pasado. El producto de un trabajo que hace uso de un bien
común, y el cómo este se distribuye y retribuye dentro de la sociedad,
es un tema que cada vez suscita más interés no solo en los gobiernos y
la ciudadanía, sino que también en la ciencia. Esto debido a que, como
el agricultor elaboró cosechas a partir de la tierra, el científico
construirá evidencia nueva a partir de conocimientos previos. Ambos
crean productos que son bienes esenciales para la ciudadanía a partir de
bienes comunes de la humanidad. El problema para ambos está en que parte
de ese quehacer ha sido privatizado, restringiendo a la ciudadanía del
libre acceso a tales bienes.

Esta problemática también es extrapolable a la situación que ocurre con
las patentes de vacunas y, en específico, con las patentes de las
\emph{vacunas de COVID-19}, donde gran parte de
\href{https://en.unesco.org/news/unesco-calls-covid-19-vaccines-be-considered-global-public-good}{los
organismos internacionales llaman a que estas sean consideradas
\emph{bienes públicos}}. Así como la privatización al acceso de un
conocimiento médico produce desigualdad entre los países más y menos
ricos, también el acceso al conocimiento por parte de la sociedad en
general -evidencia de cómo mejorar el cumplimiento de las medidas de
cuidado, por ejemplo- implicó un daño invaluable en cómo se produjo y se
distribuyó la información para combatir la pandemia.

Así, a pesar de que con \emph{la globalización y la era digital} la
labor científica ha podido crear conocimientos con mayor facilidad y
divulgarlos de manera inmediata, el desconocimiento y los mitos sobre
las leyes que amparan la propiedad intelectual han sido el principal
obstáculo para dar el paso hacia la apertura de la creación científica
(Fernández, Graziosi, y Martínez 2018). El miedo a ser sancionado por la
editorial, el temor al plagio y la pérdida de reconocimiento autoral,
destacan entre las principales razones, sin mencionar el dominio que
poseen las revistas científicas sobre el conocimiento que se genera y
que es publicado por ellas. Dicho esto, la apuesta por destinar los
esfuerzo hacia una libre circulación del conocimiento apunta a la
necesidad de reapropiarse de los beneficios, resultados y saberes
científicos. En este sentido, Banzato (2019) hace un llamado a los
organismos de América Latina para generar espacios de evaluación y
difusión que sirvan para la democratización del conocimiento, siendo
esta una estrategia cultural y política que busca promover los procesos
de producción y reproducción social del saber.

La apertura de nuestras investigaciones traen más beneficios que
dificultades: no solo contribuimos a nutrir el conocimiento colectivo
sobre un problema, sino que incluso podemos alcanzar mejor visibilidad
de nuestro trabajo científico. Por eso, en el siguiente capítulo te
presentaremos aspectos que debes considerar para lograr una exitosa
apertura de tus publicaciones y resultados de investigación.

\section{Propiedad intelectual}\label{propiedad}

Esta sección busca comprender la importancia de la propiedad intelectual
abordando los elementos centrales de la producción científica y su
resguardo legal. Es importante conocer los elementos que nos servirán en
el momento de definir el carácter y el modo en que permitiremos que
nuestras creaciones sean utilizadas y difundidas. La propiedad
intelectual no es un elemento accesorio, pues \emph{``nos habla de
facultades exclusivas que posee el autor sobre su obra y que le permiten
gozar de su uso y explotación''} (Loredo 2012). Para nuestros efectos,
se concibe como el peldaño inicial para abordar los esfuerzos hacia la
publicación libre.

\subsection{Sobre la Propiedad Intelectual y las
Licencias}\label{sobre-la-propiedad-intelectual-y-las-licencias}

La \emph{Propiedad Intelectual (PI)} es una normativa que permite el
reconocimiento autoral de toda invención de la mente humana (World
Intelectual Propiety Organization 2020), su objetivo es fomentar la
creación e innovación al mismo tiempo que regula el uso y difusión de
las obras. La PI es un tipo de ley cuya rigurosidad varía según el país
y funciona en base a dos grandes acepciones:

\begin{enumerate}
\def\labelenumi{\arabic{enumi}.}
\item
  \textbf{Common Law}: rige en \emph{Estados Unidos e Inglaterra} y
  regula la reproducción (copia) de la obra, su distribución, su
  exhibición pública y su transformación ya sea traducida o adaptada
  (Fernández 2009). A partir de lo que señala Omatos (2013) se comprende
  que es una facultad que poseen los autores y su particularidad es que
  puede ser transferida ya sea por voluntad propia o por prescripción,
  siendo así ampliados los derechos del uso de la obra bajo distintos
  tipos de licencias de \emph{Copyrigth, Creative Commons, Copyleft y
  Dominio Público}.
\item
  \textbf{Derechos de autor}: rige en \emph{Europa continental y
  Latinoamérica} y regula las facultades patrimoniales y el derecho
  moral que posee el autor sobre su obra, esto quiere decir que se
  reconoce la paternidad sobre la creación y se protege la integridad
  original del trabajo (Loredo 2012). El derecho moral, a diferencia del
  Common Law es irrenunciable e inalienable, por lo que no se puede
  traspasar.
\end{enumerate}

A modo de contexto, en el caso de Chile se han suscrito tratados
internacionales que permiten el resguardo internacional de cualquier
obra intelectual que se haya elaborado en territorio nacional, los más
importantes son dos:

\begin{enumerate}
\def\labelenumi{\arabic{enumi}.}
\item
  \textbf{Tratado de Berna} (rectificado en 1973): establece
  indicaciones a los países asociados para resguardar las obras
  internacionales en el propio territorio.
\item
  \textbf{Convenio de OMPI} (rectificado en 1975): establece una suerte
  de guía a los países adherentes que sirve para la elaboración de leyes
  que atiendan el resguardo de la propiedad intelectual dentro del mundo
  digital (Fernández 2009).
\end{enumerate}

\subsection{Tipos de Licencias}\label{tipos-de-licencias}

Las licencias sobre la propiedad intelectual regulan los derechos
patrimoniales de la obra y establecen las \emph{reglas} para su uso. Las
licencias más conocidas son:

\begin{enumerate}
\def\labelenumi{\arabic{enumi}.}
\item
  \textbf{Copyright:} Esta licencia permite que el autor se reserve
  todos los derechos sobre la obra y solo se puede hacer uso de ella
  bajo permiso del autor.
\item
  \textbf{Creative Commons (CC):} Esta licencia, inspirada en la General
  Public License - GPL (Licencia Pública General en español),\footnote{La
    \emph{General Public License} es una licencia elaborada por el
    sistema operativo \href{https://www.gnu.org/licenses/}{GNU} y su
    objetivo es permitir el uso de software y códigos de libre acceso.
    GNU es un sistema libre que busca ser compatible con Unix, otro
    sistema operativo que se caracteriza por ser portable, multitarea y
    multiusuario (Blanca 2019).} tiene el propósito de desarrollar
  herramientas digitales estandarizadas que faciliten la distribución de
  la obra, pues entrega al autor, empresa o institución la
  responsabilidad de autorizar reglamentariamente el modo de uso,
  difusión y citación de su obra. Existen seis tipos de licencias CC
  (véase Figura N° @ref(fig:licencias)) y cada una de ellas tiene como
  base la \emph{CC-BY} que permite hacer uso de la obra con la
  correspondiente atribución, el resto de licencias funcionan a modo de
  capas que se superponen a la principal. De este modo, cada capa
  entrega una especificidad sobre cómo utilizar la obra Creative Commons
  (2019).
\end{enumerate}

\begin{figure}[H]

{\centering \includegraphics[width=0.75\linewidth,height=\textheight,keepaspectratio]{images/licencias.png}

}

\caption{Capas de las licencias Creative Commons. Elaboración propia.}

\end{figure}%

\begin{enumerate}
\def\labelenumi{\arabic{enumi}.}
\setcounter{enumi}{2}
\item
  \textbf{Copyleft:} Este tipo de licencia proviene del movimiento
  \href{https://open-access.net/en/information-on-open-access/history-of-the-open-access-movement}{\emph{Open
  Access}} y se orienta a abrir el uso, aplicación, distribución y
  creación de obras. Además de permitir el uso libre, indica la
  obligación de que todo proyecto que nazca a partir del original
  contenga los principios del acceso abierto.
\item
  \textbf{Dominio Público:} Si bien no corresponde a una licencia como
  tal, es un estado en el que la obra no posee protección patrimonial
  pues ha prescrito el plazo de su protección.
\end{enumerate}

\subsection{El conocimiento es poder pero ¿de
quién?}\label{el-conocimiento-es-poder-pero-de-quiuxe9n}

Tanto organismos de investigación pública como universidades crean
conocimiento científico por medio de la investigación y la docencia con
el fin de aportar al bien público. Si bien la apertura de las
publicaciones puede ser lo ideal para tales objetivos, en ocasiones la
confidencialidad de los resultados científicos permite a sus autores
obtener beneficios económicos de su trabajo. Sin importar cuál sea el
camino, la propiedad intelectual juega un papel importante al orientar
la toma de decisiones en torno al desarrollo, difusión y utilización del
conocimiento intelectual (WIPO 2020). Por ello contar con una política
intelectual de calidad es el primer paso para gestionar estratégicamente
la difusión y transferencia de los resultados científicos.

A continuación, se presentan dos experiencias chilenas que se consideran
como buenas prácticas en términos de políticas institucionales, pues
promueven la apertura de publicaciones.

\subsubsection*{\texorpdfstring{\textbf{Política de Acceso Abierto de
ANID}}{Política de Acceso Abierto de ANID}}\label{poluxedtica-de-acceso-abierto-de-anid}
\addcontentsline{toc}{subsubsection}{\textbf{Política de Acceso Abierto
de ANID}}

La \href{https://www.anid.cl/}{\emph{Agencia Nacional de Investigación y
Desarrollo (ANID)}} que nace en 2020 como una estructura que reemplaza a
la \emph{Comisión Nacional de Investigación Científica y Tecnológica
(CONICYT)}, es hoy la institución que encabeza el \emph{Ministerio de
Ciencia, Tecnología, Conocimiento e Innovación}. Siguiendo el legado de
su antecesora, su objetivo es apoyar con recursos al desarrollo de la
ciencia y la innovación en Chile. Desde el 2021, bajo el principio de
que todo conocimiento generado con fondos públicos debe resultar en
beneficios para la sociedad, ANID ha implementado una \emph{Política de
Acceso Abierto} con el objetivo de asegurar la disponibilidad del
conocimiento científico que resulte de investigaciones financiadas por
la institución (Agencia Nacional de Investigación y Desarrollo 2020).
Esta política busca ser un curso de acción \emph{progresivo}
implementado en dos fases:

\begin{itemize}
\item
  \textbf{Fase I}: En el plazo inicial de dos años se pretende
  incentivar la apertura de las publicaciones y sus datos. Esta primera
  fase servirá para la recopilación de antecedentes sobre el uso y
  gastos asociados a la investigación abierta. En esta primera etapa de
  la política toman gran relevancia los principios \emph{FAIR}. FAIR es
  la abreviatura en lengua inglesa para referirse a Findability,
  Accessibility, Interoperability y Reusability (Encontrable, Accesible,
  Interoperable y Reutilizable). El diseño de estos principios tienen el
  objetivo de ser una guía medible para que los investigadores realicen
  una eficaz gestión de los datos y para que posteriormente puedan ser
  replicados (Wilkinson y Dumontier 2016).
\item
  \textbf{Fase II}: Los resultados de la primera fase servirán para que
  en la segunda se implemente -de manera más rigurosa- la publicación
  libre mediante la \emph{``Ruta Dorada''}, la cual corresponde a un
  formato de publicación que permite eliminar los periodos de embargo,
  dejando las publicaciones disponibles en acceso abierto de manera
  inmediata tras el pago del \emph{Article Processing Charges - APC} a
  las editoriales. El APC es una tarifa que costea el procesamiento del
  artículo final para que se adapte al diseño de la revista y, en
  ocasiones, son las propias instituciones públicas las que costean el
  APC, devolviendo así un monto no menor a las editoriales. En la
  sección de
  \href{https://lisa-coes.github.io/open-access/herramientas-para-publicar.html}{Herramientas
  Para Públicar} profundizaremos en ello.
\end{itemize}

ANID ha detectado de manera temprana el conflicto que acarrea la
industria editorial, que lidera un mercado donde los conocimientos se
tranzan como un bien. Esto beneficia principalmente a las editoriales
con un alto margen de ganancias y, por ello, en la propuesta de ANID se
especifica que:

\begin{quote}
\emph{``Esta política busca ser un curso de acción que esté en constante
revisión y que permita, de manera progresiva, avanzar hacia un sistema
transparente y abierto, donde el acceso, la re-utilización y la
constante oferta de nueva información y datos científicos contribuyan de
manera real y concreta al desarrollo social, cultural, científico y
económico de Chile''} (ANID 2020b).
\end{quote}

\subsubsection*{\texorpdfstring{\textbf{Repositorio Institucional,
Universidad de
Chile}}{Repositorio Institucional, Universidad de Chile}}\label{repositorio-institucional-universidad-de-chile}
\addcontentsline{toc}{subsubsection}{\textbf{Repositorio Institucional,
Universidad de Chile}}

La Universidad de Chile es a la fecha la institución de educación
superior pionera en desarrollar un repositorio institucional abierto.
Este recoge documentos digitales e impresos con previa autorización,
tales como tesis de pregrado y postgrado, pre y post-print, libros y
capítulos, material didáctico y presentaciones, informes técnicos,
recursos audiovisuales e imágenes. El
\href{http://repositorio.uchile.cl/}{Repositorio Académico de la
Universidad de Chile} conserva, difunde y proporciona acceso a material
científico generado por docentes e investigadores de la institución y
cuenta actualmente con más de 68.000 publicaciones.

Este repositorio académico hace uso de un protocolo de interoperabilidad
que permite que se conecte con otros repositorios, con el propósito de
incrementar la visibilidad de los documentos. Bajo este objetivo, los
autores deben proteger sus obras con licencias \emph{Creative-Commons},
de este modo aseguran el reconocimiento e identificación de la propiedad
intelectual y favorece la visibilidad del trabajo.

Ambos ejemplos de políticas institucionales proveen una exitosa
colaboración entre el mundo científico y el público general, ya que
orientan la toma de decisiones al finalizar el ejercicio investigativo y
permiten la apertura del conocimiento. Según Alperin, Babini, y Fischman
(2014), en América Latina se ha desarrollado el ejercicio del Open
Access mediante el financiamiento -casi exclusivo- de agencias estatales
y agencias de cooperación internacional. Sus resultados se publican
principalmente en revistas locales o repositorios regionales. Un ejemplo
de ello es Argentina, país donde el 80\% de los artículos científicos se
encuentran indexados en revistas locales (Alperin, Babini, y Fischman
2014), ya que la nación se ha inclinado en promover políticas de
autoarchivo como la \emph{Ley Nacional de Repositorios} promulgada en
2013 y la creación del \emph{Sistema Nacional de Repositorios Digitales}
creado en 2011 (Monti y Unzurrunzaga 2020).

La evidencia da cuenta que para el caso Latinoamericano son los
organismos universitarios y de investigación pública los responsables de
desarrollar eficaces políticas de ciencia abierta con el objetivo de
aportar a la libre circulación de los resultados científicos, pues como
se ilustra en la Figura N° @ref(fig:beneficios), dentro de sus
beneficios no solo se evidencia la mayor exposición de los trabajos
científicos, sino que también existe un aporte en términos de desarrollo
a nivel país, influencia de políticas, entre otros aspectos positivos.

\begin{figure}[H]

{\centering \includegraphics[width=1\linewidth,height=\textheight,keepaspectratio]{images/OpenAccess.png}

}

\caption{Beneficios del Open Access. Traducción propia en base a
Kingsley, D. \& Brown, S. (2021).}

\end{figure}%

\section{Formas open access}\label{formas-open-access}

El quehacer científico de las ciencias sociales es una arena de lucha
continua, donde investigadores deben mantenerse en constante
entrenamiento para no perder el reconocimiento académico que brinda
oportunidades laborales y estatus social. \emph{Publica o perece} es más
una realidad que una frase retórica, pues los científicos sociales están
en la obligación de investigar y publicar para mantener su carrera
laboral a flote. Este agotador panorama da pie a lo que muchos
investigadores de las ciencias sociales han denominado como la
\emph{crisis de la ciencia} (Breznau 2021), un nuevo panorama donde el
conocimiento científico es menos confiable producto de prácticas
antiéticas en el proceso de investigación. Falta de consistencia, sesgos
no controlados, presión editorial, entre otros factores, son los que
influyen en un tipo de ciencia con problemas de credibilidad, como es el
emblemático caso de Diederik Stapel que se mencionó en el capítulo I.

Quizás muchos confiarán la solución de la crisis a las instituciones
revisoras de las revistas, quienes aplican rigurosos métodos de control
de calidad. Sin embargo, del papel a la práctica hay un gran trecho. La
falta de transparencia en la manipulación de información, los métodos de
análisis y la recolección de datos hacen imposible corroborar la
veracidad total de un trabajo. Actualmente, muchas editoriales han
abordado el problema mediante la adopción del Open Access para promover
la transparencia de los datos y la apertura de las publicaciones, pues
con esto se permite usar la reproducción de los procesos como un
mecanismo de control. Por ello, en esta sección conoceremos las
particularidades del Open Access, las rutas por las cuales compartir
publicaciones, los repositorios y las revistas que recomendamos.

\subsection{¿Qué es el Open Access?}\label{quuxe9-es-el-open-access}

\emph{Open Access - OA} (Acceso Abierto en español) es un movimiento
cuyo objetivo es promover el acceso libre a la producción de
conocimiento científico digital y hacer uso de él sin restricciones de
copyright. Esto en base a estándares internacionales que permiten
instalar la posibilidad de que cualquier persona pueda reutilizar la
información de la investigación en cuanto datos, procesos y resultados.

Hace algún tiempo, científicos advirtieron del problema asociado al
oligopolio del conocimiento científico del que disfrutan revistas y
editoriales de gran renombre, sin mencionar conflictos entre
investigadores en torno a la competencia de estatus académico, la
manipulación de resultados, los problemas de intersubjetividad y la
falta de transparencia en el proceso de investigación (Breznau et~al.
2021). De este modo, los primeros vestigios del Open Access aparecieron
a la par con la creación de Internet, pero no fue hasta el 2002 con la
declaración de
\href{https://www.budapestopenaccessinitiative.org/translations/spanish-translation/}{Budapest}
que se definió el concepto de Open Access y que posteriormente sería
reforzado con las declaraciones que dieron lugar en
\href{https://ictlogy.net/articles/bethesda_es.html}{Bethesda 2003} y
\href{https://openaccess.mpg.de/67627/Berlin_sp.pdf}{Berlín 2003}. Uno
de los aspectos destacables de la primera declaración indica que:

\begin{quote}
\emph{Retirar las barreras de acceso a esta literatura acelerará la
investigación, enriquecerá la educación, compartirá el aprendizaje de
los ricos con los pobres y el de los pobres con el de los ricos, hará
esta literatura tan útil como sea posible y sentará los cimientos para
unir a la humanidad en una conversación intelectual común y búsqueda del
conocimiento} (Initiative 12 de Septiembre, 2012).
\end{quote}

Según Melero y Abad (2008), las causas que gatillan el desarrollo del
Open Access son principalmente dos:

\begin{enumerate}
\def\labelenumi{\arabic{enumi}.}
\item
  La constitución de Internet y las nuevas tecnologías como el medio
  para la divulgación científica permitió abaratar costos de impresión y
  acelerar el transporte de los conocimientos.
\item
  Los elevados precios de suscripción a revistas científicas digitales
  se presentaron como barreras económicas y de copyright para acceder al
  conocimiento y resultados de investigaciones.
\end{enumerate}

\subsubsection*{¿Por qué hacer uso del Open Access en mis
publicaciones?}\label{por-quuxe9-hacer-uso-del-open-access-en-mis-publicaciones}
\addcontentsline{toc}{subsubsection}{¿Por qué hacer uso del Open Access
en mis publicaciones?}

El estudio de H. Piwowar, Priem, y Orr (2019) demuestra que desde la
década de los noventa ha habido un aumento constante de publicaciones
que se adscriben al Open Access y, si bien el fenómeno se desarrolla en
mayor medida dentro de las ciencias exactas, las ciencias sociales no se
quedan atrás (Swan 2013). Según Hernández (2016), los beneficios
atribuídos a las prácticas científicas de carácter abierto no difieren
sobre la disciplina y traen consigo efectos positivos ligados a:

\begin{itemize}
\tightlist
\item
  \emph{Mayor accesibilidad y conservación de los productos
  científicos.}
\item
  \emph{Difusión rápida e inmediata de las publicaciones.}
\item
  \emph{Facilita la comunicación directa de los conocimientos
  científicos ayudando a avanzar en el mejoramiento de la calidad en la
  investigación.}
\item
  \emph{Abre la posibilidad de reutilizar información, datos y procesos
  para nuevos proyectos.}
\end{itemize}

\subsection{Las distintas rutas para publicar de manera
abierta}\label{las-distintas-rutas-para-publicar-de-manera-abierta}

Swan (2013) -representante de UNESCO- define dos rutas principales
(tipos de repositorios) donde depositar publicaciones con el rótulo del
Open Access (Dorado y Verde), pero en la práctica se pueden diferenciar
en cinco. Según Abierto (2019) las principales características de cada
una son:

\begin{enumerate}
\def\labelenumi{\arabic{enumi}.}
\item
  \textbf{Repositorios de Ruta Dorada}: Corresponden a repositorios que
  publican trabajos de investigación en forma gratuita para el lector,
  por lo que se puede leer cualquier artículo inmediatamente después de
  ser publicado. Cabe destacar que, para su publicación, detrás hay un
  cobro \emph{Article Processing Charges - APC (Cargo por Procesamiento
  de Artículo en español)} cuyo costo es pagado por los autores o las
  instituciones que les financian.
\item
  \textbf{Repositorios de Ruta Híbrida}: Gran parte de las editoriales
  académicas utilizan este modelo con el objetivo de ofrecer Acceso
  Abierto al mismo tiempo en el que mantienen su modelo de negocio
  habitual basado en suscripciones. Esto permite a los autores optar por
  pagar una cuota de publicación y tener sus artículos con Acceso
  Abierto en revistas por suscripción.
\item
  \textbf{Repositorios de Ruta Verde}: Consiste en un proceso de
  ``auto-archivo o auto-depósito'' donde los autores comparten sus
  post-print (artículo final) en sus páginas personales o en revistas
  que no son gratuitas, pero que poseen repositorios de acceso libre.
  Estos repositorios antes mencionados se adhieren a un conjunto de
  reglas y técnicas de inter-operación, formando una red interconectada
  de bases de datos.
\item
  \textbf{Repositorios de Ruta Diamante}: Es un modelo que intenta dar
  solución a los inconvenientes que surgen en las rutas anteriores. Este
  posibilita la revisión por pares de los trabajos enviados por medio de
  investigadores que trabajan como voluntarios, obteniendo de esta forma
  reconocimiento académico y social. Esta ruta es considerada como el
  único modelo que garantiza la sostenibilidad de la publicación de
  acceso abierto.
\item
  \textbf{Repositorios de Ruta Bronce}: Según H. Piwowar, Priem, y Orr
  (2019) esta es una nueva categoría donde los editores alojan artículos
  embargados en sitios web de su preferencia, de modo que los trabajos
  carecen de licencias que permitan su uso autorizado.
\end{enumerate}

\begin{figure}[H]

{\centering \includegraphics[width=1\linewidth,height=\textheight,keepaspectratio]{images/future_oa.png}

}

\caption{Proyección de la cantidad de visitas y artículos publicados por
cada ruta del Open Access. Traducción propia en base al estudio de H.
Piwowar, Priem, y Orr (2019).}

\end{figure}%

La Figura N° @ref(fig:FutureOA) ilustra el modelo elaborado por H.
Piwowar, Priem, y Orr (2019). Este estudio se basa en el análisis de 70
millones de artículos publicados entre 1950 y 2019 que comparten el
rasgo de haber sido publicados en acceso abierto y estar alojados en la
base de datos de \href{https://unpaywall.org/}{Unpaywall}. La técnica
utilizada permitió estimar la trayectoria de los artículos hasta el
2025, poniendo énfasis en el porcentaje de visitas y el porcentaje de
artículos publicados según cada ruta.

De esta manera, en la Figura N° @ref(fig:FutureOA) se proyecta que para
el 2025, los artículos publicados por la vía dorada alcanzarán un
aumento de sus visitas en un 33\%, la ruta verde un 19\% y la ruta
híbrida un 10\%; mientras que la ruta bronce seguirá la misma
trayectoria pero en menor medida. En cuanto al porcentaje de artículos
publicados en acceso abierto, esta Figura indica que tanto la ruta
verde, dorada e híbrida tendrán un aumento desde el 2019 al 2025,
mientras que la ruta bronce en sus categorías inmediata y retrasada se
mantendrán uniforme sin aumento ni descenso. Lo importante aquí es un
notable aumento de la publicación en abierto en 13 puntos porcentuales,
mientras que se estima el descenso de la publicación cerrada en 13
puntos porcentuales.

\subsection{ResearchGate, Academia.edu y Sci-Hub: ¿Plataformas de aceso
abierto?}\label{researchgate-academia.edu-y-sci-hub-plataformas-de-aceso-abierto}

De seguro conoces las plataformas de \emph{ResearchGate y Academia.edu}
o quizás en alguna ocasión has recurrido a estas páginas para compartir
tus publicaciones o descargar archivos científicos para uso personal.
Sobre estas plataformas hay que mencionar que \emph{no} son un ejemplo
de Open Access, debido a que funcionan como redes sociales similares a
\emph{Facebook o LinkedIn}, pues su objetivo es crear redes de contactos
entre colegas e investigadores con intereses afines. Según Fortney y
Gonder (2015), ResearchGate y Academia.edu, al ser empresas con fines de
lucro, no cumplen con la característica principal de los repositorios de
Open Access, los cuales son administrados por universidades, agencias
gubernamentales, asociaciones internacionales o editoriales que persigan
el objetivo de poner los conocimientos al servicio y disponibilidad
pública.

Estas páginas al ser de tipo auto-archivo presentan similitudes con los
repositorios de ruta verde, ya que los documentos se guardan en sitios
personales creados por el propio autor, sin embargo, dadas las
características de ResearchGate y Academia.edu, estas no permiten
corroborar los estándares de calidad y seguridad de las publicaciones,
ni tampoco se asegura su permanencia perpetua. Dicho lo anterior,
publicar la producción científica en cualquier plataforma digital no
implica estar desarrollando prácticas acordes al Open Access.

Las restricciones de pago empujaron al desarrollo de sitios piratas que
funcionan como bases de datos integrando diversas publicaciones, estos
son denominados como \emph{la vía negra del Acceso Abierto} que buscan
hacer frente a las barreras de pago (Björk 2017). Un claro ejemplo es el
sitio web \emph{Sci-Hub} que desde el 2011 posibilita el libre acceso a
literatura científica (tanto pagada como libre) y para el 2021 alberga
más de 85 millones de documentos. Su eslogan \emph{Sci-Hub para remover
todas las barreras en el camino de la ciencia} se instaura bajo tres
principios: (1) El conocimiento es para todos, (2) No al Copyright y (3)
Apoyo al Open Access. Sin embargo, Monti y Unzurrunzaga (2020) indica
que el carácter \emph{``ilegal''} del sitio es la razón de que
exponentes de la ciencia lo restrinjan, dado que no tiene la injerencia
para modificar los aspectos legales de las publicaciones, transgrede los
derechos de explotación indicados por los autores y puede conducir a un
desprestigio del Open Access.

Es evidente que lo atractivo de estos sitios es la posibilidad de
acceder de manera libre a las publicaciones. Un hecho corroborable es la
investigación de Bohannon (2016), en ella se ha concluido que el 25\% de
los países que registran la mayor cantidad de descargas desde Sci-hub,
corresponden a las naciones más ricas del mundo en Europa y Estados
Unidos. Björk (2017) es el primer autor en considerar la plataforma como
una vía dentro del Acceso Abierto, y es quizás una reflexión de orden
ético-moral la decisión de integrar (o no) los sitios con estas
características dentro del Open Access.

\subsection{Revistas de acceso
abierto:}\label{revistas-de-acceso-abierto}

Con la invención de Internet también aumentaron las oportunidades de
divulgación científica y, si bien el concepto de Open Access no fue
precisado hasta los 2000, en la década anterior ya venían gestándose
proyectos como \emph{The public-access computer systems review (1990)},
una revista electrónica tipo boletín que era distribuida por correo
electrónico. Posteriormente, en 1991 revistas como \emph{Surfaces y
Psycoloquy} fueron pioneras en instalar la metodología de acceso
gratuito a sus publicaciones. Hoy en día son muchos los sitios de
repositorios que cumplen con políticas de Acceso Abierto, pero con
licencias particulares que es necesario conocer para adentrarnos en su
uso. Uno de los sitios que se recomienda usar es el de
\href{https://socopen.org/}{SocArxiv}, pues cuenta con un servicio
gratuito de acceso, una base de datos completa y que permite a los
autores obtener un DOI para su publicación. Sin mencionar la posibilidad
de un contador de descargas y otras particularidades que no poseen las
plataformas de ResearchGate y Academia.edu (véase Tabla N°
@ref(tab:coom)). En la siguiente sección abordaremos en mayor
profundidad su uso.

\begin{longtable}[]{@{}
  >{\raggedright\arraybackslash}p{(\linewidth - 8\tabcolsep) * \real{0.2649}}
  >{\raggedright\arraybackslash}p{(\linewidth - 8\tabcolsep) * \real{0.1854}}
  >{\raggedright\arraybackslash}p{(\linewidth - 8\tabcolsep) * \real{0.1788}}
  >{\raggedright\arraybackslash}p{(\linewidth - 8\tabcolsep) * \real{0.1854}}
  >{\raggedright\arraybackslash}p{(\linewidth - 8\tabcolsep) * \real{0.1854}}@{}}
\caption{(\#tab:coom) Tabla comparativa del uso de SocArXiv versus otros
sitios}\tabularnewline
\toprule\noalign{}
\begin{minipage}[b]{\linewidth}\raggedright
Características
\end{minipage} & \begin{minipage}[b]{\linewidth}\raggedright
ResearchGate, Academia.edu
\end{minipage} & \begin{minipage}[b]{\linewidth}\raggedright
Sitio web personal
\end{minipage} & \begin{minipage}[b]{\linewidth}\raggedright
Repositorio Institucional
\end{minipage} & \begin{minipage}[b]{\linewidth}\raggedright
SocArXiv
\end{minipage} \\
\midrule\noalign{}
\endfirsthead
\toprule\noalign{}
\begin{minipage}[b]{\linewidth}\raggedright
Características
\end{minipage} & \begin{minipage}[b]{\linewidth}\raggedright
ResearchGate, Academia.edu
\end{minipage} & \begin{minipage}[b]{\linewidth}\raggedright
Sitio web personal
\end{minipage} & \begin{minipage}[b]{\linewidth}\raggedright
Repositorio Institucional
\end{minipage} & \begin{minipage}[b]{\linewidth}\raggedright
SocArXiv
\end{minipage} \\
\midrule\noalign{}
\endhead
\bottomrule\noalign{}
\endlastfoot
Acceso libre para leer & Requiere Registro & Si & Si & Si \\
Servicio de acceso publico - sin lucro & No & Solo si es alojado por tu
universidad & Solo si es alojado por tu universidad & Si \\
Metadatos completos, incluidos coautores, DOI, ORCID, etc. & Quizás & No
& Quizás & Si \\
Enlace para el repositorio para base de datos, codigo, etc. & No & Solo
si tu lo construyes & Quizás & Si \\
Url persistente en todas las versiones & - & No & No & Si \\
Entrega DOI para el paper & - & No & Solo Algunos & Si \\
Contador de descargas & Solo Algunos & Quizás & Quizás & Si \\
Contribuye al futuro de la comunicación académica abierta & No &
Débilmente & Quizás & Si \\
\end{longtable}

Fuente: Traducción propia a partir de tabla comparativa de
\href{https://socopen.org/frequently-asked-questions/}{SocArXiv}.

\section{Herramientas para publicar}\label{herramientas-para-publicar}

Imagine que usted es parte de un equipo de investigación que finalizó un
estudio sobre las repercusiones socio-económicas que tuvo la pandemia de
Covid-19 en el país, desarrollando un aporte significativo para la toma
de decisiones en materias de políticas económicas y levantando las bases
para nuevas líneas de investigación. Dado el aporte que significan para
el bien público los resultados de esta investigación, es necesaria su
divulgación inmediata, pero es una realidad que la publicación regular
en cualquier editorial toma un periodo largo de revisión y aceptación,
es por ello que una alternativa es la pre-publicación del documento en
un repositorio de Acceso Abierto. En la siguiente guía abordaremos un
marco de referencia para empezar a adoptar prácticas ligadas al Open
Access.

\subsection{El camino hacia la
publicación}\label{el-camino-hacia-la-publicaciuxf3n}

Todo proyecto de investigación tiene la potencialidad de ser publicado
pero, como ya se ha ido ilustrando en las secciones anteriores, publicar
de manera abierta trae consigo beneficios de visibilidad e impacto
social. Para desarrollar una estrategia eficaz de publicación abierta
hay que seguir un flujo que se divide tres pasos, como se ilustra en la
Figura N° @ref(fig:flujo). Este flujo inicia con la elección de la
revista de Acceso Abierto, sigue con la identificación de la etapa del
artículo y finaliza con la selección del servicio de publicación.

\begin{figure}[H]

{\centering \includegraphics[width=1\linewidth,height=\textheight,keepaspectratio]{images/flujo_lisa.png}

}

\caption{Guía paso a paso para publicar libre (Elaboración propia).}

\end{figure}%

\subsubsection*{Paso 1: Pre-rint
(Pre-publicación)}\label{paso-1-pre-rint-pre-publicaciuxf3n}
\addcontentsline{toc}{subsubsection}{Paso 1: Pre-rint (Pre-publicación)}

Es lo que conocemos comúnmente como el borrador del artículo, el mismo
que enviamos al proceso de revisión por pares. Este documento tiene la
posibilidad de ser publicado en cualquier repositorio de ruta verde
(\emph{pre-print server}) con el objetivo de ser difundido abiertamente
para permitir su disponibilidad inmediata.

La gran particularidad de este método de pre-publicación es que se
obtiene un código alfanumérico que identifica y ubica el documento
dentro de Internet, esto se conoce como un \emph{Digital Object
Identifier - DOI} (Identificador de Objeto Digital en español), por lo
que se convierte de manera inmediata en un documento referenciable y sin
posibilidad de plagio. Algunas personas tendrán el temor de que al
pre-publicar un informe de investigación, las posibilidades de que el
documento sea aceptado y posteriormente publicado en una revista es
menor. Sin embargo, los datos empíricos demuestran que en realidad esto
no es un impedimento para su posterior publicación. En la investigación
de Abdill y Blekhman (2019) que analizó cerca de 37.648 pre-print
alojados en bioRxiv -una extensión centrada en biología del sitio
ArXiv-, una de las grandes conclusiones tiene que ver con que la tasa de
preprints publicados en 2018 alcanzó un máximo de 2,100 por mes y, de
manera más específica, en la Figura N° @ref(fig:preprint) se puede
identificar a las diez principales revistas que han publicado la mayor
cantidad de pre-print. En esta figura, cada barra indica la cantidad de
documentos publicados por cada revista.

\begin{figure}[H]

{\centering \includegraphics[width=0.75\linewidth,height=\textheight,keepaspectratio]{images/nature_preprints.png}

}

\caption{Destino final de los Preprint. Traducción propia en base al
estudio de Abdill y Blekhman (2019)}

\end{figure}%

Para el caso de las ciencias sociales, el servicio de pre-print que
recomendamos utilizar es \href{https://socopen.org/}{\textbf{SocArxiv}}
y particularmente para la disciplina de la psicología
\href{https://psyarxiv.com/}{PsyArxiv}. Estos servicios permiten subir
archivos PDF desde la plataforma
\href{https://help.osf.io/hc/en-us/categories/360001530554-Preprints}{Open
Science Framework - OSF} (Marco de Ciencia Abierta en español) -un
software de código abierto que permite la reproducibilidad de los
procesos- y de este modo se obtiene el beneficio de acceder a un DOI que
asegure su correcta citación. Ambos repositorios son extensiones de
\href{https://arxiv.org/}{ArXiv}, la plataforma pionera de acceso
abierto que en la actualidad alberga alrededor de dos millones de
artículos categorizados en ocho áreas temáticas. ArXiv no posee tarifas
asociadas a la publicación de los artículos, puesto que los documentos
se someten a un proceso de clasificación y no a una revisión por pares.

\subsubsection*{Paso 2: Post-print}\label{paso-2-post-print}
\addcontentsline{toc}{subsubsection}{Paso 2: Post-print}

La segunda etapa es clave para el camino hacia la publicación, pues el
post-print corresponde al artículo aceptado tras la revisión por pares,
pero cuyo formato de presentación no ha sido adaptado al requerido. Por
lo tanto no ha sido publicado de manera oficial por la revista y para
lograr ello, interviene un equipo especializado de la editorial que se
encarga de tales aspectos, ya sean márgenes, tipos de citación,
estructura, entre otros.

Hoy en día si bien son varias las editoriales que entregan la
posibilidad de publicar el post-print en cualquier repositorio abierto,
esto es solo tras el periodo de embargo, el cual consiste en un tiempo
determinado donde la editorial se reserva los derechos patrimoniales del
artículo para su distribución. Dicho esto, recomendamos tener
conocimiento de las posibilidades que tiene el autor de publicar un
documento previo a la publicación oficial.

\subsubsection*{Paso 3: Print}\label{paso-3-print}
\addcontentsline{toc}{subsubsection}{Paso 3: Print}

En la última etapa del flujo, recomendamos optar por abrir la
publicación del \emph{Print} (artículo final). El artículo final que es
publicado oficialmente por una revista permite que la editorial conserve
para sí los beneficios de los derechos patrimoniales de la obra,
mientras que los equipos de investigación solo conservan el derecho al
reconocimiento. Publicar en una revista de reconocimiento que integre
políticas de acceso abierto brinda la posibilidad de que una vez
finalizado el periodo de embargo, los autores puedan abrir sus
artículos, pero no todas las revistas tienen las mismas políticas y
directrices, por ello plantean formas y periodos distintos para el
depósito del artículo en un repositorio abierto.

Al momento de realizar el envío de un artículo a una revista cualquiera,
puede ocurrir que el autor tenga que firmar un \emph{Acuerdo de
Transferencia de Derechos (CTA por sus siglas en inglés)},
transfiriéndole al editor todo derecho sobre la obra y, por lo tanto,
imposibilitando toda acción posterior del propio creador sobre la
investigación. Para que la publicación dentro de una revista no afecte a
la posterior decisión de abrir el acceso a una investigación, en
ocasiones las editoriales plantean periodos de embargo en cuyo tiempo
gozan de los beneficios económicos sobre la obra, pero al finalizar, el
autor es libre de difundir en abierto su publicación. En estos casos los
editores tienen una licencia que sirve únicamente para publicar,
mientras que los autores deben retener para sí los derechos sobre la
obra. En síntesis, para que cualquier recurso científico sea abierto,
este debe contener una licencia que explicite a sus usuarios las
acciones que pueden realizar sobre la obra e indicar la correcta
acreditación de la fuente (Swan 2013).

\subsection{Revistas para publicar con Open
Access}\label{revistas-para-publicar-con-open-access}

Por lo general, cuando llevamos tiempo investigando y publicando tenemos
ciertas nociones de las revistas a las que les puede interesar los temas
que estamos trabajando o incluso tenemos ciertas certezas de las
posibles revistas en las cuales nuestro proyecto tiene una mayor
probabilidad de ser publicado. El paso lógico para la publicación es
elegir aquella editorial donde queremos que aparezca nuestro trabajo, ya
sea por reconocimiento, por recomendaciones, por el tema que trabajamos
o por cualquier otro motivo que tengamos. Una vez elegida la revista, se
recomienda revisar las políticas de autoarchivo y para ello recomendamos
acceder a \href{https://v2.sherpa.ac.uk/romeo/}{Sherpa Romeo} y buscar
la revista escogida. Sherpa Romeo es un sitio web que funciona como una
base de datos que recopila la información básica sobre las políticas de
autoría y acceso abierto de las principales revistas científicas de todo
el mundo que utiliza un código de cuatro colores para diferenciar los
tipos de políticas de cada revista, los que se definen en la Figura N°
@ref(fig:colores-sherpa):

\begin{figure}[H]

{\centering 
\includegraphics[width=0.75\linewidth,height=\textheight,keepaspectratio]{images/colores_sherpa.png}

}

\caption{Etiquetas de colores de Sherpa Romeo según tipo de acceso
abierto (Elaboración propia).}

\end{figure}%

El mundo del Open Access es bastante grande y en ocasiones costoso, toda
revista de ruta dorada posee costos de APC que deben ser financiados por
los investigadores o sus patrocinantes, lo que se vuelve una barrera
económica para los equipos de investigación que deseen abrir sus
publicaciones, por ello el proceso de autoarchivo surge como una
alternativa al alcance de cualquiera y, en su mayoría, cientistas
sociales utilizan la vía verde y la diamante para depositar sus
manuscritos o tesis en repositorios como
\href{https://soporte.dialnet.unirioja.es/portal/es/kb/articles/el-acceso-abierto-a-la-literatura-cient\%C3\%ADfica}{Dialnet},
\href{https://www.latindex.org/latindex/descripcion}{Latindex},
\href{https://scielo.org/es/sobre-el-scielo/declaracion-de-accesso-abierto/}{Scielo},
\href{https://www.redalyc.org/}{Redalyc} y
\href{http://biblioteca.clacso.edu.ar/accesoabierto/}{CLACSO}.
Adicionalmente, existen revistas de ciencias sociales cuyo foco se
encuentra en el diseño de metodologías de investigación cualitativa y
que, además, poseen un alto factor de impacto, como las resumidas en la
Tabla N° @ref(tab:revc), pues son indexadas a Scopus (Ulloa y Mardones
2017).

\begin{longtable}[]{@{}
  >{\raggedright\arraybackslash}p{(\linewidth - 4\tabcolsep) * \real{0.6199}}
  >{\raggedright\arraybackslash}p{(\linewidth - 4\tabcolsep) * \real{0.1462}}
  >{\raggedright\arraybackslash}p{(\linewidth - 4\tabcolsep) * \real{0.2339}}@{}}
\caption{(\#tab:revc) Descripción de Revistas de Open Access de orden
Cualitativo}\tabularnewline
\toprule\noalign{}
\begin{minipage}[b]{\linewidth}\raggedright
Revista
\end{minipage} & \begin{minipage}[b]{\linewidth}\raggedright
País
\end{minipage} & \begin{minipage}[b]{\linewidth}\raggedright
Asunto o Categoría
\end{minipage} \\
\midrule\noalign{}
\endfirsthead
\toprule\noalign{}
\begin{minipage}[b]{\linewidth}\raggedright
Revista
\end{minipage} & \begin{minipage}[b]{\linewidth}\raggedright
País
\end{minipage} & \begin{minipage}[b]{\linewidth}\raggedright
Asunto o Categoría
\end{minipage} \\
\midrule\noalign{}
\endhead
\bottomrule\noalign{}
\endlastfoot
\href{http://www.qualitativesociologyreview.org}{Qualitative Sociology
Review} & Polonia & Ciencias Sociales diversas \\
\href{https://journals.sagepub.com/home/ijq}{The International Journal
Qualitative Methods} & Canadá & Educación \\
\href{https://www.qualitative-research.net/index.php/fqs}{Forum
Qualitative Sozialforscung} & Alemania & Ciencias Sociales diversas \\
\end{longtable}

Fuente: Adaptación propia a partir de tabla extraída del estudio de
Ulloa y Mardones (2017).

\subsection{Pagos asociados a la publicación
abierta}\label{pagos-asociados-a-la-publicaciuxf3n-abierta}

El Open Access está redefiniendo la forma en que el capital científico
se define, se moviliza y se trata. Este es un nuevo paradigma que choca
con las reglas de un mercado académico que por décadas ha mantenido el
monopolio económico del conocimiento científico por medio de barreras de
pago, un modelo que no siempre es receptivo a alteraciones y novedades
(Sadaba 2014), pero que ha tenido la obligación de adaptarse al nuevo
panorama cibernético digital.

El modelo de suscripción se presenta como un muro de pago que al ser
costeado permite acceder al material científico. Detrás de cada ejemplar
físico existe un costo de adaptación, impresión y envío, lo que obligaba
a las editoriales a publicar una cantidad determinada. Sin embargo, con
el advenimiento de Internet, la digitalización y la inmediatez de la
información, las razones y argumentos que daban sustento al modelo se
volvieron obsoletas. La era digital permitió avanzar a pasos agigantados
en el desarrollo de nuevos conocimientos, en la difusión masiva e
instantánea de los trabajos académicos y, por consecuencia, abrió paso
al cuestionamiento de las suscripciones pagadas como método de acceso al
conocimiento. En este sentido, la declaración que tuvo lugar en Bethesda
(2003) fue precisa en indicar que \emph{``Internet ha cambiado
fundamentalmente las realidades prácticas y económicas relacionadas con
la distribución del conocimiento científico y el patrimonio cultural''}.

En la nueva era digital, el modelo de suscripción sigue estando
presente, pero convive con el paradigma del open access, el que intenta
superar las barreras de pago poniendo la responsabilidad de la apertura
del conocimiento en el propio autor. Algunos repositorios de ruta dorada
e híbrida solicitan un pago de APC para publicar en acceso abierto, esto
corresponde a un modelo de negocios cuyo propósito es financiar los
gastos asociados a la gestión de la publicación. En estricto rigor, no
existe un monto estandarizado del APC, por lo que su costo dependerá
únicamente de la revista. Su valor fluctúa según el pago de impuestos
adicionales, lo que se puede evidenciar en la Tabla N° @ref(tab:precio).
Según Socha (2018), de las revistas de mayor renombre en la industria
académica como \href{https://www.elsevier.com/open-science}{Elsevier},
\href{https://www.springer.com/la/open-access}{Springer},
\href{https://authorservices.wiley.com/author-resources/Journal-Authors/open-access/index.html}{Wiley}
y
\href{https://authorservices.taylorandfrancis.com/publishing-open-access/}{Taylor
\& Francis} se han adscrito al acceso abierto cobrando a los autores un
costo de APC pero manteniendo su modelo de suscripción.

\begin{longtable}[]{@{}
  >{\raggedright\arraybackslash}p{(\linewidth - 6\tabcolsep) * \real{0.2647}}
  >{\raggedright\arraybackslash}p{(\linewidth - 6\tabcolsep) * \real{0.1912}}
  >{\raggedright\arraybackslash}p{(\linewidth - 6\tabcolsep) * \real{0.1912}}
  >{\raggedright\arraybackslash}p{(\linewidth - 6\tabcolsep) * \real{0.3529}}@{}}
\caption{(\#tab:precio) Fluctuación del precio en dolares del APC en las
principales revistas científicas}\tabularnewline
\toprule\noalign{}
\begin{minipage}[b]{\linewidth}\raggedright
Repositorio
\end{minipage} & \begin{minipage}[b]{\linewidth}\raggedright
APC mínimo
\end{minipage} & \begin{minipage}[b]{\linewidth}\raggedright
APC máximo
\end{minipage} & \begin{minipage}[b]{\linewidth}\raggedright
¿Open Acces o Híbrido?
\end{minipage} \\
\midrule\noalign{}
\endfirsthead
\toprule\noalign{}
\begin{minipage}[b]{\linewidth}\raggedright
Repositorio
\end{minipage} & \begin{minipage}[b]{\linewidth}\raggedright
APC mínimo
\end{minipage} & \begin{minipage}[b]{\linewidth}\raggedright
APC máximo
\end{minipage} & \begin{minipage}[b]{\linewidth}\raggedright
¿Open Acces o Híbrido?
\end{minipage} \\
\midrule\noalign{}
\endhead
\bottomrule\noalign{}
\endlastfoot
ELSEVIER & \$100 & \$5.000 & Híbrido \\
SPRINGER & \$3.000 & \$3.000 & Híbrido \\
WILEY & \$1.300 & \$5.200 & Híbrido \\
Taylor y Francis & \$500 & \$2.950 & Open Access \\
\end{longtable}

Fuente: Adaptación propia a partir de información recopilada en
\href{https://www.openaccess.cam.ac.uk/publishing-open-access/how-much-do-publishers-charge-open-access}{University
of Cambridge (2018)}.

Spinak (2019) indica que tanto en Estados Unidos como en países Europeos
ha habido un aumento exponencial del costo del APC a través de los años.
A modo de ejemplo, el \emph{APC promedio de 319 revistas} asociadas a
las editoriales \emph{BMC, Frontiers, MDPI e Hindawi} aumentó entre 2,5
y 6 veces la inflación entre el 2012 y el 2018, al alero de un
considerable crecimiento en la cantidad de volúmenes que paso de 58.007
hasta 127.528 en aquellos años. Los editores han admitido que el costo
del APC lo fijan según el valor económico de la revista que es atribuido
según el valor de impacto.

La introducción de un sistema de pagos por procesamiento de artículos
hace frente al modelo de suscripción, pero no logra reemplazarlo. En
cambio, funciona de manera paralela dentro del mercado científico,
presentándose como una nueva barrera que obstaculiza la decisión de
abrir las publicaciones de aquellas investigaciones que son autónomas o
que dispongan de un presupuesto acotado. Los APC tienen sus propios
cuestionamientos y uno de ellos refiere a que el pago no tiene mucho que
ver con el procesamiento y posterior publicación del artículo, sino que
más bien los equipos de investigación pagan por el reconocimiento de ser
publicados en una revista de alto factor de impacto (Velterop 2018).
Pareciera ser que el nuevo ciclo de la ciencia que apuntaba a la
democratización del conocimiento se vio inmerso en las lógicas
económicas en donde las editoriales obtienen los principales beneficios
monetarios.

\emph{¿Pagar por hacer ciencia o pagar por reconocimiento?} Claro esta
que existen otras vías gratuitas para publicar de manera abierta, sin
embargo la ciencia -y principalmente las ligadas a la medicina, la
física y las matemáticas- ha tenido la necesidad imperiosa de ser
creíble, reconocida y citada, y por ello existe una mayor competencia en
torno al reconocimiento que otorga la publicación en revistas de alto
impacto. Sin embargo, el quehacer en pos del conocimiento y su
disposición como bien público para la sociedad y su desarrollo poco
tiene que ver con una distinción simbólica que sirve dentro del mundo
científico. La ciencia social, a diferencia de las ciencias exactas, se
ha podido desarrollar con mayor libertad por medio de las vías gratuitas
en importantes repositorios de ruta verde, como es el ejemplo del caso
latinoamericano que profundizaremos en la siguiente sección.

\section{¿Cuál es el destino de las ciencias sociales en el mundo del
Open
Access?}\label{cuuxe1l-es-el-destino-de-las-ciencias-sociales-en-el-mundo-del-open-access}

La dicotomía entre ciencias naturales y sociales es una discusión
epistemológica de años que no solo remite al objeto de estudio, sino que
también a las metodologías de investigación, a los criterios de validez
y quizás lo más importante, a la reflexión en torno al carácter objetivo
de la propia ciencia. Es claro que el movimiento del Open Access no
pretende dar respuestas ni reflexiones profundas sobre este tipo de
cuestionamientos de orden epistemológico, y aunque se abre en extenso
como una herramienta que sirve para la libre difusión de los
conocimientos en el plano de Internet, una cosa es clara y es que las
ciencias sociales han llegado algo tarde en comparación a las ciencias
exactas (Sadaba 2014). La comunidad científica ligada al campo de lo
social ha hecho los esfuerzos por adquirir prácticas abiertas y, de este
modo, han comenzado a trazar una próspera trayectoria en el campo del
Open Access con un fuerte componente político.

Las barreras de pago no son el único obstáculo que enfrenta la libre
circulación del conocimiento científico, Banzato (2019) indica que la
industria editorial ha instalado la idea de una ciencia del centro o
bien denominada \emph{mainstream}, un concepto consolidado globalmente
gracias al trabajo hecho por los servicios de
\href{https://clarivate.com/webofsciencegroup/release-notes/wos/new-wos-admin-faq/}{Web
Of Science WoS} y
\href{https://www.elsevier.com/solutions/scopus?dgcid=RN_AGCM_Sourced_300005030}{Scopus}
que han creado indicadores bibliométricos que segmentan a las
investigaciones científicas, fomentan el uso del inglés como la lengua
ilustre de la ciencia, entregan prestigio a aquellos trabajos que siguen
los estándares de publicación y son indexadas en las revistas de gran
renombre. En suma, generan un sistema restrictivo donde existe poca
representatividad de las revistas de ciencias sociales. Según Beigel
(2018), las bases de datos comerciales no solo contribuyeron en la
ilusión de una ciencia del \emph{centro}, sino que también desarrollaron
la idea de un tipo de ciencia \emph{``Periférica''}, que refiere al
trabajo que se reclusa en el terreno de la academia regional-local (como
es el caso Latinoamericano). Este tipo de ciencia se indexa en revistas
subvaloradas y comprendidas como endogámicas y de baja calidad, pues en
su gran mayoría producen y re-producen el conocimiento dentro de la
propia región, teniendo así un rol subsidiario en el desarrollo de la
ciencia y no protagónico como es el caso anterior.

La apuesta del Open Access de hacer frente al oligopolio de la industria
editorial y lograr un cambio sociocultural en la ciencia a menudo tiene
que luchar con condiciones sociales, culturales y económicas que
terminan modelando la forma en que los cambios tecnológicos se
convierten en cambios sociales. Por ello el Open Access es un proceso
diferencial no solo entre las disciplinas, sino que también entre las
zonas geográficas y las distintas regiones del mundo (Sadaba 2014). La
globalización y el auge de Internet por sí mismos no democratizan ni
permiten la distribución equitativa del conocimiento, por ello es
importante adscribir a las palabras de Babini (2014), coordinadora del
programa de acceso abierto de CLACSO, quien hace un llamado a la
comunidad científica de tomar responsabilidad por gestionar
comunicaciones académicas no comerciales, cuyo objetivo sea permitir que
el conocimiento se convierta en un \textbf{bien público} y, por lo
tanto, se gestione como tal.

La experiencia latinoamericana ha logrado arribar en el puerto de la
ruta verde de una forma similar al contexto anglosajón, en el estudio de
H. Piwowar, Priem, y Orr (2019) se evidencia que tanto en ciencias
sociales como en psicología se utiliza mayoritariamente la ruta verde
para la publicación abierta (véase Figura N° @ref(fig:piwowar2019)).
Para el caso Latinoamericano, se ha creado una estructura no comercial
donde la publicación libre es manejada por institutos de investigación
sin fines de lucro o por universidades estatales, dando forma a un
ecosistema sustentado por \href{https://openjournalsystems.com/}{Open
Journals Systems} (Sistema de Revistas Abiertas en español), un software
libre que entrega herramientas para la administración de plataformas
digitales como
\href{https://www.latindex.org/latindex/inicio}{Latindex},
\href{https://www.redalyc.org/}{Redalyc} y
\href{https://scielo.org/es/}{Scielo}, los principales exponentes del
Open Access en la región según Becerril García y Aguado López (2019).
Estas plataformas sirvieron de trampolín para el desarrollo de la
ciencia abierta de tipo \emph{no comercial}, una distinción propia de
Latinoamérica (Beigel 2021).

\begin{figure}[H]

{\centering \includegraphics[width=1\linewidth,height=\textheight,keepaspectratio]{images/Piwowar_future.png}

}

\caption{Piwowar, H. et al.~(2019). Porcentaje de diferentes tipos de
acceso de una muestra aleatoria de artículos y reseñas de WoS con un DOI
publicado entre 2009 y 2015 por disciplina (Traducción propia).}

\end{figure}%

Algunas personas ligaran de manera exclusiva el Open Acces a
metodologías de carácter cuantitativo, desconfiando del uso de estas
técnicas para la investigación cualitativa con el temor de la pérdida de
confidencialidad o el reproche sobre la toma de decisiones
metodológicas. La evidencia levantada por Mardones, Ulloa, y Salas
(2018) demuestra que existen experiencias cualitativas de trabajos
publicados en abierto que no tienen repercusiones negativas, aún cuando
estos declaran y describen sus diseños. Ello demuestra que la
publicación abierta no afecta la particularidad y flexibilidad del
diseño, siempre y cuando este se dirija con el fin de develar el
fenómeno de estudio.

Taylor (2020) realizó un análisis sistemático de la publicación en
abierto de libros y capítulos en el área de las ciencias sociales, pues
ha habido una falta de atención en la publicación de libros en abierto,
muy por el contrario al análisis que se realiza casi por completo a
artículos y revistas. Este estudio comprende la existencia de un gran
número de científicos sociales que se abocan a publicar en libros y,
dentro de los resultados de la investigación, se evidencia que el número
de libros y capítulos publicados en acceso abierto es mucho más bajo que
el de los artículos. Sin embargo, a través de los años ha ido en un
ligero incremento, siendo estos en su gran mayoría de las disciplinas de
la psicología y la economía, y en último lugar, textos del área de la
filosofía. Lo anterior no solo evidencia un auge constante de la
publicación de libros y capítulos en abierto, también da cuenta de que
al igual que en el panorama general, dentro de las humanidades y las
ciencias sociales ocurre una dicotomía entre las disciplinas.

Las revistas académicas y los repositorios antes mencionados son
administrados, dirigidos y financiados por instituciones académicas o
por organismos no gubernamentales sin fines de lucro -como es el caso
del Centro Latinoamericano de Ciencias Sociales (CLACSO)-, por lo que no
se incluyen tarifas ni para autores ni para lectores. América Latina ha
desarrollado una suerte de \emph{ecosistema en red}, que se configura
cómo un espacio virtual donde los miembros de las universidades ponen a
disposición el contenido intelectual en repositorios de sus casas de
estudio, los cuales a su vez se encuentran interconectados a otros
repositorios académicos. Un ejemplo claro de ello es
\href{https://www.lareferencia.info/es/}{LaReferencia} que integra más
de 1.3 millones de archivos según Becerril García y Aguado López (2019).
En términos de políticas públicas, los organismos estatales han
desarrollado leyes que promueven el uso de prácticas abiertas en el
desarrollo de investigaciones con fondos estatales, sin embargo, por lo
que señala Becerril García y Aguado López (2019), los mandatos tienden a
ser aún muy débiles y operan como sugerencias más que como obligaciones
metodológicas. Otra debilidad indicada por el mismo estudio da cuenta de
que la mayoría de los sistemas de validación social de las
investigaciones descansan en el renombre de la revista, aún cuando estas
se encuentren lejos de las discusiones locales o bien ajenas a los
contextos latinoamericanos.

El destino de las ciencias sociales en el mundo del Open Access es
bastante amplio. Los beneficios que entrega este paradigma no solo
contribuyen a la difusión y amplio reconocimiento de la obra y de su
autoría, también aportan a la apertura del conocimiento científico que
sirve a nivel local en la toma de decisiones, en la consecución de
estrategias políticas y a nivel regional en tanto aporta al desarrollo
cultural y económico de la sociedad en su conjunto. Sin duda, el camino
que han trazado otras personas sigue su curso e invita no solo a
investigadores a incurrir en este tipo de prácticas, sino que también
pone la responsabilidad en las instituciones universitarias y
gubernamentales en proporcionar las herramientas, guías y conocimientos
necesarios para que el Open Access sea una alternativa viable por la
cual la ciencia y el conocimiento se transformen.

\bookmarksetup{startatroot}

\chapter{Conclusión}\label{conclusiuxf3n}

Este libro se plantea como objetivo resaltar los principales elementos
que revisten de importancia a la Ciencia Abierta y, a partir de ello, se
pretende constituir como una guía para los científicos sociales que
quieran suscribirse bajo los principios de la Ciencia Social Abierta.
Los cuatro principios que guían este libro se posicionan como las
prácticas y experiencias de mayor relevancia para superar la crisis de
la ciencia actual (Breznau 2021), donde el conocimiento científico es
menos confiable producto de prácticas antiéticas en el proceso de
investigación.

A lo largo de este libro se han señalado las principales herramientas
que permitirán posicionar su investigación dentro de los márgenes de la
Ciencia Abierta, como los \textbf{preregistros}, \textbf{preprints}, el
uso de software libres, estándares de calidad para publicación de los
datos de la investigación, entre otros. En suma, estas herramientas
permiten que su investigación adquiera un estándar mínimo de calidad,
robustez y confianza, siempre en miras de lograr una mayor transparencia
del diseño de investigación, una correcta apertura de los datos
recopilados, un proceso de análisis que sea reproducible y una
publicación de hallazgos y resultados que esté abierta a la comunidad
científica y al público en general.

Se hace necesario avanzar en una mayor divulgación de estas prácticas y
herramientas para así lograr hacer de ellas una experiencia cotidiana
dentro de la investigación en las Ciencias Sociales empíricas. De esta
forma, se podrán elaborar investigaciones más confiables que erradiquen
las prácticas antiéticas a lo largo de todo el proceso de investigación.

Finalmente, nuestra aspiración es que este trabajo quede abierto a
comentarios, propuestas, sugerencias y revisiones. Por esta razón es
que, si bien no se ha seguido un proceso tradicional de investigación
con un diseño de investigación, recopilación de datos y análisis de
resultados, el texto que constituye a este libro se encuentra en un
repositorio público (https://github.com/lisa-coes/lisa-book). De este
modo, quien se interese por revisar, replicar o corregir lo que se
expone en este libro podrá hacerlo con la mayor facilidad y
transparencia posible.

W

\bookmarksetup{startatroot}

\chapter*{References}\label{references}
\addcontentsline{toc}{chapter}{References}

\markboth{References}{References}

\phantomsection\label{refs}
\begin{CSLReferences}{1}{0}
\bibitem[\citeproctext]{ref-abdill_tracking_2019}
Abdill, Richard J., y Ran Blekhman. 2019. {«Tracking the Popularity and
Outcomes of All {bioRxiv} Preprints»}. \emph{bioRxiv}, 515643.
\url{https://doi.org/10.1101/515643}.

\bibitem[\citeproctext]{ref-universoabierto_5_2019}
Abierto, Universo. 2019. {«Las 5 Rutas Para Llegar al Acceso Abierto:
Verde, Dorada, Bronce, H{í}brida y Diamante»}. \emph{Blog de la
biblioteca de Traducci{ó}n y Documentaci{ó}n de la Universidad de
Salamanca}.

\bibitem[\citeproctext]{ref-abrilruiz_manzanas_2019}
Abril Ruiz, Angel. 2019. \emph{{Manzanas podridas: Malas pr{á}cticas de
investigaci{ó}n y ciencia descuidada}}.

\bibitem[\citeproctext]{ref-aczel_consensusbased_2020}
Aczel, Balazs, Barnabas Szaszi, Alexandra Sarafoglou, Zoltan Kekecs,
Šimon Kucharský, Daniel Benjamin, Christopher D. Chambers, et~al. 2020.
{«A Consensus-Based Transparency Checklist»}. \emph{Nature Human
Behaviour} 4 (1): 4-6. \url{https://doi.org/10.1038/s41562-019-0772-6}.

\bibitem[\citeproctext]{ref-agencianacionaldeinvestigacionydesarrollo_consulta_2020}
Agencia Nacional de Investigación y Desarrollo, (ANID). 2020. {«Consulta
{P{ú}blica}: {Pol{í}tica Acceso Abierto} a {Informaci{ó}n
Cient{í}fica}»}.

\bibitem[\citeproctext]{ref-agnoli_questionable_2017}
Agnoli, Franca, Jelte M. Wicherts, Coosje L. S. Veldkamp, Paolo Albiero,
y Roberto Cubelli. 2017. {«Questionable Research Practices Among Italian
Research Psychologists»}. \emph{PLOS ONE} 12 (3): e0172792.
\url{https://doi.org/10.1371/journal.pone.0172792}.

\bibitem[\citeproctext]{ref-alperin_indicadores_2014}
Alperin, Juan Pablo, Dominique Babini, y Gustavo Fischman. 2014.
\emph{{Indicadores de acceso abierto y comunicaciones acad{é}micas en
Am{é}rica Latina}}. Juan Pablo Alperin. Vol. 1. Buenos Aires, Argentina:
Consejo Latinoamericano de Ciencias Sociales, CLACSO.

\bibitem[\citeproctext]{ref-anid_con_2020}
ANID. 2020a. {«Con {Amplio Consenso ANID Inicia Hoja} de {Ruta} de
{Pol{í}tica} de {Acceso Abierto}»}.

\bibitem[\citeproctext]{ref-anid_propuesta_2020}
ANID. 2020b. {«Propuesta de {Pol{í}tica} de Acceso Abierto a La
Informaci{ó}n Cient{í}fica y Adatos de Investigaci{ó}nfinanciados Con
Fondos P{ú}blicos de La {ANID}»}. ANID.

\bibitem[\citeproctext]{ref-babini_universidades_2014}
Babini, Dominique. 2014. {«{Universidades y acceso abierto: hora de
tomar protagonismo}»}. En \emph{{Foro Revista Iberoamericana de Ciencia,
Tecnolog{í}a y Sociedad}}, 1-3. 2015.

\bibitem[\citeproctext]{ref-banzato_soberania_2019}
Banzato, Guillermo. 2019. {«{Soberan{í}a del conocimiento para superar
inequidades: pol{í}ticas de Acceso Abierto para revistas cient{í}ficas
en Am{é}rica Latina}»}. \emph{Mecila Working Paper Series} 18: 1-18.

\bibitem[\citeproctext]{ref-barba_terminologies_2018}
Barba, Lorena A. 2018. {«Terminologies for {Reproducible Research}»}.
\emph{arXiv:1802.03311 {[}cs{]}}, febrero.
\url{https://arxiv.org/abs/1802.03311}.

\bibitem[\citeproctext]{ref-becerrilgarcia_end_2019}
Becerril García, Arianna, y Eduardo Aguado López. 2019. {«The {End} of a
{Centralized Open Access Project} and the {Beginning} of a
{Community-Based Sustainable Infrastructure} for {Latin America}»}.
\emph{OpenEdition Press}, 41-55.
\href{https://doi.org/10.4000/books.oep.\%209003.}{https://doi.org/10.4000/books.oep.
9003.}

\bibitem[\citeproctext]{ref-beigel_relaciones_2018}
Beigel, Fernanda. 2018. {«{Las relaciones de poder en la ciencia
mundial}»}. \emph{Nueva Sociedad} 274: 13-28.

\bibitem[\citeproctext]{ref-beigel_america_2021}
---------. 2021. {«{Am{é}rica Latina podr{í}a convertirse en l{í}der
mundial de la ciencia abierta no comercial}»}. \{Scientific Blog\}.
\emph{The Conversation}.

\bibitem[\citeproctext]{ref-bethesda_declaracion_2003}
Bethesda. 2003. {«Declaraci{ó}n de {Bethesda} Sobre Publicaci{ó}n de
Acceso Abierto»}.

\bibitem[\citeproctext]{ref-bishop_rein_2019}
Bishop, Dorothy. 2019. {«Rein in the Four Horsemen of
Irreproducibility»}. \emph{Nature} 568 (7753): 435-35.
\url{https://doi.org/10.1038/d41586-019-01307-2}.

\bibitem[\citeproctext]{ref-bishop_revisiting_2017}
Bishop, Libby, y Arja Kuula-Luumi. 2017. {«Revisiting {Qualitative Data
Reuse}: {A Decade On}»}. \emph{SAGE Open} 7 (1): 215824401668513.
\url{https://doi.org/10.1177/2158244016685136}.

\bibitem[\citeproctext]{ref-bjork_gold_2017}
Björk, Bo-Christer. 2017. {«Gold, {Green} and {Black Open Access}»}.
\emph{Learned Publishing} 30 (2): 173-75.
\url{https://doi.org/10.1002/leap.1096}.

\bibitem[\citeproctext]{ref-blanca_informatica_2019}
Blanca, José. 2019. {«{Inform{á}tica para las ciencias de la vida: Unix
y Python}»}. \{Course\}. \emph{Bioinformatics at COMAV}.

\bibitem[\citeproctext]{ref-bohannon_whos_2016}
Bohannon, John. 2016. {«Who's Downloading Pirated Papers? {Everyone}»}.
\emph{American Association for the Advancement of Science} 352 (6285):
508-12. \url{https://doi.org/10.1126/science.352.6285.508}.

\bibitem[\citeproctext]{ref-breznau_future_2019}
Breznau, Nate. 2019. {«The {Future} of {Sociology Depends} on {Open
Science}»}. Preprint. SocArXiv.
\url{https://doi.org/10.31235/osf.io/d37be}.

\bibitem[\citeproctext]{ref-breznau_does_2021}
---------. 2021. {«Does {Sociology Need Open Science}?»}
\emph{Societies} 11 (1): 9. \url{https://doi.org/10.3390/soc11010009}.

\bibitem[\citeproctext]{ref-breznau_observing_2021}
Breznau, Nate, Eike Mark Rinke, Alexander Wuttke, Muna Adem, Jule
Adriaans, Amalia Alvarez-Benjumea, Henrik Kenneth Andersen, et~al. 2021.
{«Observing {Many Researchers Using} the {Same Data} and {Hypothesis
Reveals} a {Hidden Universe} of {Uncertainty}»}. MetaArXiv.
\url{https://doi.org/10.31222/osf.io/cd5j9}.

\bibitem[\citeproctext]{ref-brodeur_star_2016}
Brodeur, Abel, Mathias Lé, Marc Sangnier, y Yanos Zylberberg. 2016.
{«Star {Wars}: {The Empirics Strike Back}»}. \emph{American Economic
Journal: Applied Economics} 8 (1): 1-32.
\url{https://doi.org/10.1257/app.20150044}.

\bibitem[\citeproctext]{ref-bueno_what_2017}
Bueno, Gema. 2017. {«What {Is Open Science}? {Introduction}.»}

\bibitem[\citeproctext]{ref-burlig_improving_2018}
Burlig, Fiona. 2018. {«Improving Transparency in Observational Social
Science Research: {A} Pre-Analysis Plan Approach»}. \emph{Economics
Letters} 168 (julio): 56-60.
\url{https://doi.org/10.1016/j.econlet.2018.03.036}.

\bibitem[\citeproctext]{ref-byington_solutions_2017}
Byington, Eliza, y Will Felps. 2017. {«Solutions to the {Credibility
Crisis} in {Management Science}»}. \emph{Academy of Management Learning
and Education, The} 16 (marzo): 142-62.
\url{https://doi.org/10.5465/amle.2015.0035}.

\bibitem[\citeproctext]{ref-camerer_evaluating_2018}
Camerer, Colin F., Anna Dreber, Felix Holzmeister, Teck-Hua Ho, Jürgen
Huber, Magnus Johannesson, Michael Kirchler, et~al. 2018. {«Evaluating
the Replicability of Social Science Experiments in {Nature} and
{Science} Between 2010 and 2015»}. \emph{Nature Human Behaviour} 2 (9):
637-44. \url{https://doi.org/10.1038/s41562-018-0399-z}.

\bibitem[\citeproctext]{ref-card_role_2011}
Card, David, Stefano DellaVigna, y Ulrike Malmendier. 2011. {«The {Role}
of {Theory} in {Field Experiments}»}. \emph{Journal of Economic
Perspectives} 25 (3): 39-62. \url{https://doi.org/10.1257/jep.25.3.39}.

\bibitem[\citeproctext]{ref-carey_fraud_2011}
Carey, Benedict. 2011. {«Fraud {Case Seen} as a {Red Flag} for
{Psychology Research}»}. \emph{The New York Times}, noviembre.

\bibitem[\citeproctext]{ref-ccsds_recommendation_2012}
CCSDS. 2012. {«Recommendation for {Space Data System Practices MAGENTA
BOOK REFERENCE MODEL FOR AN OPEN ARCHIVAL INFORMATION SYSTEM}
({OAIS})»}.

\bibitem[\citeproctext]{ref-chambers_registered_2013}
Chambers, Christopher D. 2013. {«Registered {Reports}: {A} New
Publishing Initiative at~{Cortex}»}. \emph{Cortex} 49 (3): 609-10.
\url{https://doi.org/10.1016/j.cortex.2012.12.016}.

\bibitem[\citeproctext]{ref-chambers_registered_2015}
Chambers, Christopher D., Zoltan Dienes, Robert D. McIntosh, Pia
Rotshtein, y Klaus Willmes. 2015. {«Registered Reports: Realigning
Incentives in Scientific Publishing»}. \emph{Cortex; a Journal Devoted
to the Study of the Nervous System and Behavior} 66: A1-2.
\url{https://doi.org/10.1016/j.cortex.2015.03.022}.

\bibitem[\citeproctext]{ref-chin_questionable_2021}
Chin, Jason M., Justin T. Pickett, Simine Vazire, y Alex O. Holcombe.
2021. {«Questionable {Research Practices} and {Open Science} in
{Quantitative Criminology}»}. \emph{Journal of Quantitative
Criminology}, agosto. \url{https://doi.org/10.1007/s10940-021-09525-6}.

\bibitem[\citeproctext]{ref-christensen_transparent_2019}
Christensen, Garret, Jeremy Freese, y Edward Miguel. 2019.
\emph{Transparent and Reproducible Social Science Research: How to Do
Open Science}. Oakland, California: University of California Press.

\bibitem[\citeproctext]{ref-christensen_transparency_2018}
Christensen, Garret, y Edward Miguel. 2018. {«Transparency,
{Reproducibility}, and the {Credibility} of {Economics Research}»}.
\emph{Journal of Economic Literature} 56 (3): 920-80.
\url{https://doi.org/10.1257/jel.20171350}.

\bibitem[\citeproctext]{ref-corti_managing_2019}
Corti, Louise. 2019. \emph{Managing and {Sharing Research Data}: {A
Guide} to {Good Practice}.} 2nd edition. Thousand Oaks, CA: SAGE
Publications.

\bibitem[\citeproctext]{ref-creativecommons_cc_2019}
Creative Commons, CC. 2019. {«About {CC Licenses}»}. \{\{FAQ\}\}.
\emph{What We Do}.

\bibitem[\citeproctext]{ref-ddc_using_2017}
DDC. 2017. {«Using {RISE}, the {Research Infrastructure Self Evaluation
Framework} {\textbar} {DCC}.»}

\bibitem[\citeproctext]{ref-doucouliagos_are_2013}
Doucouliagos, Chris, y T. D. Stanley. 2013. {«Are {All Economic Facts
Greatly Exaggerated}? {Theory Competition} and {Selectivity}»}.
\emph{Journal of Economic Surveys} 27 (2): 316-39.
\url{https://doi.org/10.1111/j.1467-6419.2011.00706.x}.

\bibitem[\citeproctext]{ref-earth_reproducibility_2019}
Earth, Division on, y Cognitive Board on Behavioral. 2019.
{«Reproducibility and {Replicability} in {Science}»}. \emph{Undefined}.

\bibitem[\citeproctext]{ref-ec_fair_2016}
EC. 2016. {«{FAIR Data Management} in {Horizon} 2020.»}

\bibitem[\citeproctext]{ref-editors_observational_2014}
Editors, The PLOS Medicine. 2014. {«Observational {Studies}: {Getting
Clear} about {Transparency}»}. \emph{PLOS Medicine} 11 (8): e1001711.
\url{https://doi.org/10.1371/journal.pmed.1001711}.

\bibitem[\citeproctext]{ref-elliott_taxonomy_2020}
Elliott, Kevin C. 2020. {«A {Taxonomy} of {Transparency} in {Science}»}.
\emph{Canadian Journal of Philosophy}, 1-14.
\url{https://doi.org/10.1017/can.2020.21}.

\bibitem[\citeproctext]{ref-fair_fair_2020}
FAIR, GO. 2020. {«{FAIR Principles}»}.

\bibitem[\citeproctext]{ref-fanelli_how_2009}
Fanelli, Daniele. 2009. {«How {Many Scientists Fabricate} and {Falsify
Research}? {A Systematic Review} and {Meta-Analysis} of {Survey Data}»}.
\emph{PLOS ONE} 4 (5): e5738.
\url{https://doi.org/10.1371/journal.pone.0005738}.

\bibitem[\citeproctext]{ref-ferguson_how_2014}
Ferguson, Liz. 2014. {«How and {Why Researchers Share Data} (and {Why
They Don}'t)»}.

\bibitem[\citeproctext]{ref-fernandez_derechos_2009}
Fernández, Juan Carlos. 2009. {«Derechos de {Autor}»}. \emph{Derechos de
Autor en plataformas e-learning}.

\bibitem[\citeproctext]{ref-fernandez_derechos_2018}
Fernández, Juan Carlos, Eduardo Graziosi, y Daniel Martínez. 2018.
{«Derechos de Autor y Ciencia Abierta: El Papel de La Biblioteca
Universitaria»}. En \emph{{VIII Conferencia Biredial-ISTEC}}. Lima -
Per{ú}.

\bibitem[\citeproctext]{ref-fiedler_questionable_2016}
Fiedler, Klaus, y Norbert Schwarz. 2016. {«Questionable {Research
Practices Revisited}»}. \emph{Social Psychological and Personality
Science} 7 (1): 45-52. \url{https://doi.org/10.1177/1948550615612150}.

\bibitem[\citeproctext]{ref-fortney_social_2015}
Fortney, Katie, y Justin Gonder. 2015. {«A Social Networking Site Is Not
an Open Access Repository»}. \emph{Office of Scholary Communication
University of California}.

\bibitem[\citeproctext]{ref-franco_publication_2014}
Franco, A., N. Malhotra, y G. Simonovits. 2014. {«Publication Bias in
the Social Sciences: {Unlocking} the File Drawer»}. \emph{Science} 345
(6203): 1502-5. \url{https://doi.org/10.1126/science.1255484}.

\bibitem[\citeproctext]{ref-fraser_questionable_2018}
Fraser, Hannah, Tim Parker, Shinichi Nakagawa, Ashley Barnett, y Fiona
Fidler. 2018. {«Questionable Research Practices in Ecology and
Evolution»}. \emph{PLOS ONE} 13 (7): e0200303.
\url{https://doi.org/10.1371/journal.pone.0200303}.

\bibitem[\citeproctext]{ref-freese_replication_2017}
Freese, Jeremy, y David Peterson. 2017. {«Replication in {Social
Science}»}. \emph{Annual Review of Sociology} 43 (1): 147-65.
\url{https://doi.org/10.1146/annurev-soc-060116-053450}.

\bibitem[\citeproctext]{ref-gall_credibility_2017}
Gall, Thomas, John P. A. Ioannidis, y Zacharias Maniadis. 2017. {«The
Credibility Crisis in Research: {Can} Economics Tools Help?»} \emph{PLOS
Biology} 15 (4): e2001846.
\url{https://doi.org/10.1371/journal.pbio.2001846}.

\bibitem[\citeproctext]{ref-gerber_publication_2008}
Gerber, Alan, y Neil Malhotra. 2008a. {«Publication {Bias} in {Empirical
Sociological Research}: {Do Arbitrary Significance Levels Distort
Published Results}?»} \emph{Sociological Methods \& Research} 37 (1):
3-30. \url{https://doi.org/10.1177/0049124108318973}.

\bibitem[\citeproctext]{ref-gerber_statistical_2008}
---------. 2008b. {«Do {Statistical Reporting Standards Affect What Is
Published}? {Publication Bias} in {Two Leading Political Science
Journals}»}. \emph{Quarterly Journal of Political Science} 3 (3):
313-26. \url{https://doi.org/10.1561/100.00008024}.

\bibitem[\citeproctext]{ref-gilbert_comment_2016}
Gilbert, D. T., G. King, S. Pettigrew, y T. D. Wilson. 2016. {«Comment
on "{Estimating} the Reproducibility of Psychological Science"»}.
\emph{Science} 351 (6277): 1037-37.
\url{https://doi.org/10.1126/science.aad7243}.

\bibitem[\citeproctext]{ref-gomez_datos_2016}
Gómez, Nancy Diana, Eva María Méndez Rodríguez, y Antonio Hernández
Pérez. 2016. {«Datos y Metadatos de Investigaci{ó}n En Ciencias Sociales
y Humanidades: Una Aproximaci{ó}n Desde Los Repositorios Tem{á}ticos de
Datos»}.

\bibitem[\citeproctext]{ref-hernandez_por_2016}
Hernández, Enrique. 2016. {«{\textquestiondown}{Por} Qu{é} Open
Access?»} 28 (1).

\bibitem[\citeproctext]{ref-icsu_open_2014}
ICSU. 2014. {«Open {Access} to {Scientific Data} and {Literature} and
the {Assessment} of {Research} by {Metrics}.»}

\bibitem[\citeproctext]{ref-budapestopenaccessinitiative_diez_2012}
Initiative, Budapest Open Access. 12 de Septiembre, 2012. {«{Diez a{ñ}os
desde la Budapest Open Access Initiative: hacia lo abierto por
defecto}»}. \emph{BOAI}.

\bibitem[\citeproctext]{ref-john_measuring_2012}
John, Leslie K., George Loewenstein, y Drazen Prelec. 2012. {«Measuring
the {Prevalence} of {Questionable Research Practices With Incentives}
for {Truth Telling}»}. \emph{Psychological Science} 23 (5): 524-32.
\url{https://doi.org/10.1177/0956797611430953}.

\bibitem[\citeproctext]{ref-kerr_harking_1998}
Kerr, Norbert L. 1998. {«{HARKing}: {Hypothesizing After} the {Results}
Are {Known}»}. \emph{Personality and Social Psychology Review} 2 (3):
196-217. \url{https://doi.org/10.1207/s15327957pspr0203_4}.

\bibitem[\citeproctext]{ref-kryvokhyzha_best_2019}
Kryvokhyzha, Dmytro. 2019. {«The {Best Free Research Data Repository}.»}
Bioinformatics \& Genomics Scientist.

\bibitem[\citeproctext]{ref-lindsay_seven_2020}
Lindsay, D. Stephen. 2020. {«Seven Steps Toward Transparency and
Replicability in Psychological Science.»} \emph{Canadian
Psychology/Psychologie Canadienne} 61 (4): 310-17.
\url{https://doi.org/10.1037/cap0000222}.

\bibitem[\citeproctext]{ref-loredo_derecho_2012}
Loredo, Alejandro. 2012. {«Derecho {Comparado}: {Derecho} de {Autor} y
{Copyright}. {Dos} Caminos Que Se Encuentran»}. \emph{Portal de gobierno
electr{ó}nico, inclusi{ó}n digital y sociedad del conocimiento}.

\bibitem[\citeproctext]{ref-makel_both_2021}
Makel, Matthew C., Jaret Hodges, Bryan G. Cook, y Jonathan A. Plucker.
2021. {«Both {Questionable} and {Open Research Practices Are Prevalent}
in {Education Research}»}. \emph{Educational Researcher}, marzo,
0013189X211001356. \url{https://doi.org/10.3102/0013189X211001356}.

\bibitem[\citeproctext]{ref-mardones_usos_2018}
Mardones, Rodolfo, Jorge Ulloa, y Gonzálo Salas. 2018. {«{Usos del
dise{ñ}o metodol{ó}gico cualitativo en art{í}culos de acceso abierto de
alto impacto en ciencias sociales}»}. \emph{Forum: Qualitative Social
Research} 19, n{\(^\circ\)}1: 1-18.

\bibitem[\citeproctext]{ref-marsden_introducing_2018}
Marsden, Emma, Kara Morgan-Short, Pavel Trofimovich, y Nick C. Ellis.
2018. {«Introducing {Registered Reports} at {Language Learning}:
{Promoting Transparency}, {Replication}, and a {Synthetic Ethic} in the
{Language Sciences}»}. \emph{Language Learning} 68 (2): 309-20.
\url{https://doi.org/10.1111/lang.12284}.

\bibitem[\citeproctext]{ref-marz_depositing_2000}
Marz, K., y C. S. Dunn. 2000. \emph{Depositing Data with the {Data
Resources Program} of the {National Institute} of {Justice}: {A}
Handbook}. Ann Arbor, MI: {Inter-university Consortium for Political and
Social Research}.

\bibitem[\citeproctext]{ref-melero_revistas_2008}
Melero, Remedios, y María Francisca Abad. 2008. {«Revistas {Open
Access}: Caracter{í}sticas, Modelos Econ{ó}micos y Tendencias»} 20.

\bibitem[\citeproctext]{ref-mertens_preregistration_2019}
Mertens, Gaëtan, y Angelos-Miltiadis Krypotos. 2019. {«Preregistration
of {Analyses} of {Preexisting Data}»}. \emph{Psychologica Belgica} 59
(1): 338-52. \url{https://doi.org/10.5334/pb.493}.

\bibitem[\citeproctext]{ref-merton_sociology_1973}
Merton, Robert K. 1973. \emph{The {Sociology} of {Science}:
{Theorerical} and {Empirical Investigations}}. Chicago; London: The
University of Chicago Press.

\bibitem[\citeproctext]{ref-miguel_promoting_2014}
Miguel, E., C. Camerer, K. Casey, J. Cohen, K. M. Esterling, A. Gerber,
R. Glennerster, et~al. 2014. {«Promoting {Transparency} in {Social
Science Research}»}. \emph{Science} 343 (6166): 30-31.
\url{https://doi.org/10.1126/science.1245317}.

\bibitem[\citeproctext]{ref-monti_acceso_2020}
Monti, Carolina, y Carolina Unzurrunzaga. 2020. {«{Acceso a la
literatura cient{í}fica desde Sci-Hub: An{á}lisis y reflexi{ó}n de las
descargas en Argentina}»}. \emph{Revista Hipertextos} 8 (14): 111-16.
\url{https://doi.org/10.24215/23143924e022}.

\bibitem[\citeproctext]{ref-moore_preregister_2016}
Moore, Don A. 2016. {«Preregister If You Want To»}. \emph{The American
Psychologist} 71 (3): 238-39. \url{https://doi.org/10.1037/a0040195}.

\bibitem[\citeproctext]{ref-motyl_state_2017}
Motyl, Matt, Alexander P. Demos, Timothy S. Carsel, Brittany E. Hanson,
Zachary J. Melton, Allison B. Mueller, J. P. Prims, et~al. 2017. {«The
State of Social and Personality Science: {Rotten} to the Core, Not so
Bad, Getting Better, or Getting Worse?»} \emph{Journal of Personality
and Social Psychology} 113 (1): 34-58.
\url{https://doi.org/10.1037/pspa0000084}.

\bibitem[\citeproctext]{ref-nosek_promoting_2015}
Nosek, B. A., G. Alter, G. C. Banks, D. Borsboom, S. D. Bowman, S. J.
Breckler, S. Buck, et~al. 2015. {«Promoting an Open Research Culture»}.
\emph{Science} 348 (6242): 1422-25.
\url{https://doi.org/10.1126/science.aab2374}.

\bibitem[\citeproctext]{ref-nosek_preregistration_2018}
Nosek, Brian A., Charles R. Ebersole, Alexander C. DeHaven, y David T.
Mellor. 2018. {«The Preregistration Revolution»}. \emph{Proceedings of
the National Academy of Sciences} 115 (11): 2600-2606.
\url{https://doi.org/10.1073/pnas.1708274114}.

\bibitem[\citeproctext]{ref-nosek_registered_2014}
Nosek, Brian A., y Daniël Lakens. 2014. {«Registered {Reports}: {A
Method} to {Increase} the {Credibility} of {Published Results}»}.
\emph{Social Psychology} 45 (3): 137-41.
\url{https://doi.org/10.1027/1864-9335/a000192}.

\bibitem[\citeproctext]{ref-oboyle_chrysalis_2017}
O'Boyle, Ernest Hugh, George Christopher Banks, y Erik Gonzalez-Mulé.
2017. {«The {Chrysalis Effect}: {How Ugly Initial Results Metamorphosize
Into Beautiful Articles}»}. \emph{Journal of Management} 43 (2): 376-99.
\url{https://doi.org/10.1177/0149206314527133}.

\bibitem[\citeproctext]{ref-ocde_open_2020}
OCDE. 2020. {«Open {Science} - {OECD}»}.

\bibitem[\citeproctext]{ref-omatos_aspectos_2013}
Omatos, Antonio. 2013. {«Aspectos {Legales} En La {Educaci{ó}n}»}.

\bibitem[\citeproctext]{ref-opensciencecollaboration_estimating_2015}
Open Science Collaboration. 2015. {«Estimating the Reproducibility of
Psychological Science»}. \emph{Science} 349 (6251): aac4716-16.
\url{https://doi.org/10.1126/science.aac4716}.

\bibitem[\citeproctext]{ref-peng_reproducible_2011}
Peng, R. D. 2011. {«Reproducible {Research} in {Computational
Science}»}. \emph{Science} 334 (6060): 1226-27.
\url{https://doi.org/10.1126/science.1213847}.

\bibitem[\citeproctext]{ref-piwowar_data_2013}
Piwowar, Heather A., y Todd J. Vision. 2013. {«Data Reuse and the Open
Data Citation Advantage»}. \emph{PeerJ} 1 (octubre): e175.
\url{https://doi.org/10.7717/peerj.175}.

\bibitem[\citeproctext]{ref-piwowar_future_2019}
Piwowar, Heather, Jason Priem, y Richard Orr. 2019. {«The {Future} of
{OA}: {A} Large-Scale Analysis Projecting {Open Access} Publication and
Readership»}. \emph{BioRxiv}, 795310.

\bibitem[\citeproctext]{ref-rabelo_questionable_2020}
Rabelo, André L. A., Jéssica E. M. Farias, Maurício M. Sarmet, Teresa C.
R. Joaquim, Raquel C. Hoersting, Luiz Victorino, João G. N. Modesto, y
Ronaldo Pilati. 2020. {«Questionable Research Practices Among
{Brazilian} Psychological Researchers: {Results} from a Replication
Study and an International Comparison»}. \emph{International Journal of
Psychology} 55 (4): 674-83. \url{https://doi.org/10.1002/ijop.12632}.

\bibitem[\citeproctext]{ref-ramirez_ciencia_2019}
Ramírez, Paola, y Daniel Samoilovich. 2019. {«Ciencia {Abierta}.
{Reporte Para Tomadores} de {Decisiones}»}.

\bibitem[\citeproctext]{ref-rosenthal_file_1979}
Rosenthal, Robert. 1979. {«The File Drawer Problem and Tolerance for
Null Results.»} \emph{Psychological Bulletin} 86 (3): 638-41.
\url{https://doi.org/10.1037/0033-2909.86.3.638}.

\bibitem[\citeproctext]{ref-sadaba_acceso_2014}
Sadaba, Igor. 2014. {«{EL ACCESO ABIERTO EN CIENCIAS SOCIALES}: {NOTAS
SOCIOL{Ó}GICAS SOBRE PUBLICACIONES}, {COMUNIDADES Y CAMPOS}»} 17:
93-113.

\bibitem[\citeproctext]{ref-sharan_ten_2020}
Sharan, Malvika. 2020. {«Ten {Arguments Against Open Science That You
Can Win} {\textbar} {Software Sustainability Institute}»}.

\bibitem[\citeproctext]{ref-simmons_falsepositive_2011}
Simmons, Joseph P., Leif D. Nelson, y Uri Simonsohn. 2011.
{«False-{Positive Psychology}: {Undisclosed Flexibility} in {Data
Collection} and {Analysis Allows Presenting Anything} as
{Significant}»}. \emph{Psychological Science} 22 (11): 1359-66.
\url{https://doi.org/10.1177/0956797611417632}.

\bibitem[\citeproctext]{ref-simonsohn_pcurve_2014}
Simonsohn, Uri, Leif D. Nelson, y Joseph P. Simmons. 2014. {«P-Curve:
{A} Key to the File-Drawer»}. \emph{Journal of Experimental Psychology:
General} 143 (2): 534-47. \url{https://doi.org/10.1037/a0033242}.

\bibitem[\citeproctext]{ref-socha_cuanto_2018}
Socha, Beata. 2018. {«{{\textquestiondown}Cu{á}nto cobran los
principales editores comerciales por tener un art{í}culo en acceso
abierto?}»} \emph{Universo Abierto}.

\bibitem[\citeproctext]{ref-spinak_revistas_2019}
Spinak, Ernesto. 2019. {«{Revistas que han aumentado el valor del APC
han recibido m{á}s art{í}culos}»}. \emph{Scielo en Perspectiva}.

\bibitem[\citeproctext]{ref-steneck_fostering_2006}
Steneck, Nicholas H. 2006. {«Fostering Integrity in Research:
{Definitions}, Current Knowledge, and Future Directions»}. \emph{Science
and Engineering Ethics} 12 (1): 53-74.
\url{https://doi.org/10.1007/PL00022268}.

\bibitem[\citeproctext]{ref-stewart_preregistration_2020}
Stewart, Suzanne, Eike Mark Rinke, Ronan McGarrigle, Dermot Lynott,
Carole Lunny, Alexandra Lautarescu, Matteo M. Galizzi, Emily K. Farran,
y Zander Crook. 2020. {«Pre-Registration and {Registered Reports}: A
{Primer} from {UKRN}»}. OSF Preprints.

\bibitem[\citeproctext]{ref-swan_directrices_2013}
Swan, Alma. 2013. \emph{Directrices Para Pol{í}ticasde Desarrollo y
Promoci{ó}n Del Acceso Abierto}. UNESCO.

\bibitem[\citeproctext]{ref-taylor_altmetric_2020}
Taylor, Michael. 2020. {«An Altmetric Attention Advantage for Open
Access Books in the Humanities and Social Sciences»}.
\emph{Scientometrics} 125: 2523-43.
\url{https://doi.org/10.1007/s11192-020-03735-8}.

\bibitem[\citeproctext]{ref-tierney_realistic_2020}
Tierney, Nicholas J., y Karthik Ram. 2020. {«A {Realistic Guide} to
{Making Data Available Alongside Code} to {Improve Reproducibility}»}.
\emph{arXiv:2002.11626 {[}cs{]}}, febrero.
\url{https://arxiv.org/abs/2002.11626}.

\bibitem[\citeproctext]{ref-ulloa_tendencias_2017}
Ulloa, Jorge, y Rodolfo Mardones. 2017. {«Tendencias Paradigm{á}ticas y
T{é}cnicas Conversacionales En Investigaci{ó}n {Cualitativa} En
{Ciencias Sociales}»}. \emph{Perspectivas de la Comunicaci{ó}n},
Universidad de La {Frontera}, 10, n{\(^\circ\)}1: 213-35.

\bibitem[\citeproctext]{ref-unesco_que_2020}
UNESCO. 2020. {«Qu{é} {Es Ciencia Abierta}? {UNESCO Lanza Consulta
Global}»}.

\bibitem[\citeproctext]{ref-universitarias_make_2020}
Universitarias, Bibliotecas. 2020. {«Make {Your Data Discoverable}
{\textbar} {University} of {Oklahoma Libraries}.»}

\bibitem[\citeproctext]{ref-velterop_suscripciones_2018}
Velterop, Jan. 2018. {«{De suscripciones y Tasas de Procesamiento de
Art{í}culos}»}. \{Scientific Blog\}. \emph{Scielo en Perspectiva}.

\bibitem[\citeproctext]{ref-viscusi_role_2014}
Viscusi, W. Kip. 2014. {«The {Role} of {Publication Selection Bias} in
{Estimates} of the {Value} of a {Statistical Life}»}. Working
\{\{Paper\}\} 20116. National Bureau of Economic Research.

\bibitem[\citeproctext]{ref-vivalt_heterogeneous_2015}
Vivalt, Eva. 2015. {«Heterogeneous {Treatment Effects} in {Impact
Evaluation}»}. \emph{American Economic Review} 105 (5): 467-70.
\url{https://doi.org/10.1257/aer.p20151015}.

\bibitem[\citeproctext]{ref-whyte_open_2011}
Whyte, Angus, y Graham Pryor. 2011. {«Open {Science} in {Practice}:
{Researcher Perspectives} and {Participation}»}. \emph{International
Journal of Digital Curation} 6 (1): 199-213.
\url{https://doi.org/10.2218/ijdc.v6i1.182}.

\bibitem[\citeproctext]{ref-wilkinson_fair_2016}
Wilkinson, Mark, y Michel Dumontier. 2016. {«The {FAIR Guiding
Principles} for Scientific Data Management and Stewardship»} 3 (160018).
\url{https://doi.org/10.1038/sdata.2016.18}.

\bibitem[\citeproctext]{ref-wingen_no_2020}
Wingen, Tobias, Jana B. Berkessel, y Birte Englich. 2020. {«No
{Replication}, {No Trust}? {How Low Replicability Influences Trust} in
{Psychology}»}. \emph{Social Psychological and Personality Science} 11
(4): 454-63. \url{https://doi.org/10.1177/1948550619877412}.

\bibitem[\citeproctext]{ref-wipo_frequently_2020}
WIPO. 2020. {«Frequently {Asked Questions}: {IP Policies} for
{Universities} and {Research Institutions}»}. Science
\{\{Organization\}\}. \emph{World Intelectual Property Organization}.

\bibitem[\citeproctext]{ref-worldintelectualpropietyorganization_what_2020}
World Intelectual Propiety Organization, WIPO. 2020. \emph{What Is
Intellectual Property?} WIPO. Switzerland.

\end{CSLReferences}




\end{document}
